\documentclass[11pt]{amsart}

% Encoding and fonts
\usepackage[T1]{fontenc}
\usepackage[utf8]{inputenc}
\usepackage{lmodern}

% Math packages
\usepackage{amsmath}
\usepackage{amssymb}
\usepackage{amsfonts}
\usepackage{amsthm}
\usepackage{mathtools}

% Layout and utilities
\usepackage{geometry}
\usepackage{enumitem}
\usepackage{mathrsfs}

\geometry{margin=1in}

\numberwithin{equation}{section}
\setcounter{tocdepth}{3}
\setcounter{secnumdepth}{3}

% ---------------------------------------------------------
% Theorem-like environments
% ---------------------------------------------------------

\theoremstyle{plain}
\newtheorem{theorem}{Theorem}[section]
\newtheorem{proposition}[theorem]{Proposition}
\newtheorem{lemma}[theorem]{Lemma}
\newtheorem{corollary}[theorem]{Corollary}
\newtheorem{conjecture}[theorem]{Conjecture}

\theoremstyle{definition}
\newtheorem{definition}[theorem]{Definition}
\newtheorem{experiment}[theorem]{Experiment}
\newtheorem{heuristic}[theorem]{Heuristic}
\newtheorem{question}[theorem]{Question}

\theoremstyle{remark}
\newtheorem{remark}[theorem]{Remark}
\newtheorem{warning}[theorem]{Warning}

\newcommand{\R}{\mathbb{R}}
\newcommand{\C}{\mathbb{C}}
\newcommand{\N}{\mathbb{N}}
\newcommand{\Qp}{\mathbb{Q}_{>0}}
\newcommand{\Hil}{\mathcal{H}}
\newcommand{\Sch}{\mathcal{S}}
\newcommand{\Four}{\mathcal{F}}
\newcommand{\Dom}{\mathcal{D}}
\newcommand{\Bop}{\mathcal{B}}
\newcommand{\id}{\mathrm{Id}}
\newcommand{\ad}{\operatorname{ad}}
\newcommand{\Reop}{\operatorname{Re}}
\newcommand{\Om}{\Omega}
\newcommand{\twconv}{\star_{\vartheta}}
\newcommand{\deltaone}{\delta_{1}}
\newcommand{\ellft}{\lambda}
\newcommand{\Log}{\operatorname{Log}}

% hyperref must be loaded last
\usepackage[hidelinks]{hyperref}

\title[Heisenberg Uncertainty and M\"obius Dynamics XII]{
Heisenberg Uncertainty and Twisted M\"obius Dynamics XII:\
Twisted Curvature, Replica Filtering, Mertens Record Laws,\\
Critical-Scale Limits, Perturbative Blindness, and Global Heisenberg Geometry
}

\author{}
\date{}

\begin{document}

\begin{abstract}
We study an arithmetic specialization of the Weyl canonical commutation
relations obtained from the map
\(n\mapsto (\ell\log n,\kappa\Omega(n))\).  The resulting Weyl cocycle
defines a twisted Dirichlet convolution and a deformed M\"obius inverse
\(z_\vartheta\).  We prove an exact closed formula for its second deformation
curvature, derive its Dirichlet generating function as a prime-log covariance
determinant, and isolate the logarithmic branching inherited from the prime
zeta function.  We then compare the twisted inverse with a norm-convergent
ordered Euler object.  The odd infinitesimal Euler-locality defect has a
Dirichlet series \(\mathcal O(s)\) which, for \(\Re s>1\), admits both an
absolutely convergent prime-pair expansion and a Fourier-distance integral.

The principal new analytic result carried into Version 9 is a domain-complete treatment
of that integral.  Absolute convergence, local uniform convergence,
differentiation, and all sum--integral interchanges are justified in their
stated domains.  After continuation from \(\Re s>1\), every contribution
arriving at a scaled zero \(s_0=\rho/r\) is included, including rational
collisions.  For the normalized energy
\(\mathcal I(s)=2\pi\mathcal O(s)/\zeta(s)\), the singular--singular block attached to \((\rho,r)\) has
logarithmic-over-pole coefficient
\(-\pi m_\rho^2\varphi(r)r^{-2}\sum_{d\mid r}\mu(d)\).
Consequently this coefficient vanishes identically for every \(r>1\) and
survives for \(r=1\).  We prove that no other contour singularity can restore
a term of the same order.  At a pure replica the next coefficient is a
strictly negative divisor energy \(\beta_r\), for which we give both an
integral and a finite divisor-sum formula, together with the exact radical law
\(\beta_r=(\operatorname{rad}r/r)\beta_{\operatorname{rad}r}\).  We also
record the exact distinction between the normalized energy \(\mathcal I\) and
\(\mathcal O\) itself: multiplication by \(\zeta\) lowers the singular order
at an actual zeta zero.  Version 9 further polarizes the arithmetic curvature
into a positive Gram kernel and combines it, after the exact transverse
primitivization of the prime-zeta replicas, with Weil's quadratic functional.
We prove that the resulting curvature--zero kernel is positive semidefinite
for every compactly supported smooth Weil test function if and only if RH
holds.  This is an exact reformulation, not a proof of RH: the converse already
contains Weil positivity on a single diagonal entry.  The genuinely new target
is therefore an unconditional Gram realization of the combined kernel from
the Weyl/cocycle side before the zero set is used.  No such realization is
claimed here.

Version 10 adds a separate Mertens-side development.  Let
\(A^*(x)=\max_{n\le x}|M(n)|\) and
\(\Theta=\sup_{\zeta(\rho)=0}\Re\rho\).  Combining the classical
oscillation theorem of Kaczorowski--Pintz with the standard upper bound for
Mertens sums to the right of the zero supremum gives
\(A^*(x)=x^{\Theta+o(1)}\).  From this we prove uniform separation of
scales, a logarithmic historical \(L^p\)-collapse, the record-time law
\(r_k=k^{1/\Theta+o(1)}\), and
\(\int_1^{r_k}M(t)dt/t=o(M(r_k)\log r_k)\).  We then lift the weighted
Mertens path to an auxiliary Heisenberg area \(\mathcal A_t\).  Its
order-zero term satisfies the exact identity
\(2\mathcal A_0=MI-\int M^2dt/t\), while its first transverse derivative
is the original arithmetic cocycle.  A standard Selberg--Delange argument
gives \(\sum_{n\le x}\mu(n)\Omega(n)\sim x/(\log x)^2\), and hence
the local first-jet law
\(\Delta\mathcal A'_0(N)=\mu(N)N/(2\log N)+o(N/\log N)\).  These
record laws use classical analytic inputs; their Heisenberg reformulation is
recorded as structural mathematics, not as a new restriction on zeta zeros.

Version 11 closed the fixed-jet gap by passing to the complete squarefree
parameter family
\(P_x(z)=\sum_{n\le x}\mu(n)^2z^{\Omega(n)}\), for which
\(F_t(x)=P_x(-e^t)\).  Uniform Selberg--Delange theory near \(z=-1\)
proves a unique simple real critical root \(\xi_*(x)\to-1\) and the sharp
localization
\(1+\xi_*(x)=M(x)(\log x)^2/x\,(1+o(1))\), without an uncontrolled Taylor
tail.  On the critical parameter scale \(t=u/\log\log x\), the complete
renormalized family converges uniformly on compact subsets of the complex
\(u\)-plane to \(u e^{-u}\).  Cauchy's formula then controls the whole jet
tower and its Taylor tails uniformly.  We also prove a deterministic
factorization law for the noncritical roots, uniform dominance of the
critical root in all logarithmic-derivative traces, an inner Mertens scale
\(t\asymp M/M_\Omega\), and the corresponding critical Heisenberg-area law.
Finally, a Perron--Mellin representation is derived for the global quadratic
first-jet increments.  This last transform is recorded as the principal open
route for testing whether the accumulated Heisenberg area sees zeta
singularities directly rather than only through Mertens record times.

Version 12 closes a further false lead.  Differentiating the exact Heisenberg
increment to arbitrary fixed order and combining it with the compact-uniform
critical-scale theorem proves that every fixed transverse jet has a universal
leading term generated by the local Selberg--Delange singularity at \(s=1\).
In particular the second jet is no escape: its leading scale is
\(-\mu(N)N\log\log N/\log N\).  More generally the whole perturbative window
\(t=O(1/\log\log N)\) is governed, at leading order, by the same universal
profile \(ue^{-u}\).  Thus the dependence on the zero supremum \(\Theta\) in
record evaluations enters through the record locations, not through any fixed
local jet.  This isolates the genuinely nonlocal candidates: the accumulated
area, the Perron--Mellin transform, nonperturbative deformations, and remainders
beyond the complete local Selberg--Delange expansion.
\end{abstract}

\maketitle

\tableofcontents

\bigskip

\begin{remark}[Version 12: mathematical status]
\label{rem:status}
The canonical Weyl calculus, Robertson uncertainty inequality, general
2-cocycle machinery, prime-zeta continuation, and distance-energy identities
are standard and are cited to primary or standard references.  The operator
series based only on translations have direct precedents in the spectral
operator literature; in particular, Herichi--Lapidus rigorously study
\(\zeta(\partial)\) and its M\"obius inverse
\cite{HerichiLapidus2014}.  The notebook-specific content is the arithmetic
specialization of the cocycle, the exact twisted M\"obius curvature, the
ordered Euler-locality defect, and the replica-filtering identities proved
below.

Version 10 retains the complete replica analysis introduced in Version 8, the positivity layer of Version 9, and adds a Mertens-record/Heisenberg-jet layer.  Version 11 adds the squarefree deformation and removes the
previous fixed-jet/Taylor-tail gap by proving compact-uniform critical-scale
limits and Cauchy bounds for the full jet tower.  Version 12 proves that the
second jet and, in fact, every fixed transverse jet are still governed at
leading order by the same local Selberg--Delange mechanism at \(s=1\).  It
therefore closes the proposed ``higher fixed jet'' escape route and sharpens
the boundary: only nonlocal accumulation, nonperturbative deformation, or the
global Perron--Mellin transform can plausibly carry new spectral information.
Version 9 added a positive-kernel formulation tied to Weil's criterion.  Every theorem-like statement used in the main chain is
accompanied by a proof.  Domains of unbounded operators, domains of absolute
convergence, differentiability, normal convergence, exchanges of sums and
integrals, boundary terms, branch choices, contour singularities, and possible
rational collisions are stated explicitly.  The phrase ``principal
singularity'' always refers to the displayed asymptotic scale and never to an
unstated global ordering of multivalued germs.
\end{remark}

\section{Introduction: the Version 12 structure}
\label{sec:introduction}

Version 1 placed a standard Heisenberg pair next to a zeta multiplier but did
not couple the two structures.  Version 2 corrected that defect: the exact
canonical commutation relation deformed the arithmetic product itself through
a nontrivial cocycle.  The Version 2 audit nevertheless identified two
remaining weaknesses.  First, the actual uncertainty principle was not stated
prominently.  Second, several domain, inverse, and convergence claims required
repair.

Version 12 starts from the self-adjoint position and momentum operators
\begin{equation}
(Q\psi)(q)=q\psi(q),
\qquad
(P\psi)(q)=-i\hbar\psi'(q),
\label{eq:intro_QP}
\end{equation}
with the canonical commutation relation
\begin{equation}
[Q,P]=i\hbar I
\label{eq:intro_CCR}
\end{equation}
on the Schwartz core.  For every normalized Schwartz state, the Heisenberg--Robertson uncertainty inequality is \cite{Robertson1929}
\begin{equation}
\boxed{\Delta_\psi Q\,\Delta_\psi P\ge\frac{\hbar}{2}.}
\label{eq:intro_heisenberg_uncertainty}
\end{equation}
Here the uncertainty is a mathematical spread intrinsic to the state, not an
experimental error.

The arithmetic specialization assigns to an integer \(n\) the phase-space
increments
\begin{equation}
x_n=\ell\log n,
\qquad
y_n=\kappa\Omega(n),
\label{eq:intro_phase_space_coordinates}
\end{equation}
and the self-adjoint generator
\begin{equation}
G_n=y_nQ-x_nP.
\label{eq:intro_Gn}
\end{equation}
The arithmetic specialization of the standard Robertson inequality gives, for two integers,
\begin{equation}
\boxed{
\Delta_\psi G_m\,\Delta_\psi G_n
\ge
\frac{\hbar\ell\kappa}{2}
\left|\Omega(m)\log n-\log m\,\Omega(n)\right|.}
\label{eq:intro_arithmetic_uncertainty}
\end{equation}
For distinct primes \(p,q\), this becomes
\begin{equation}
\boxed{
\Delta_\psi G_p\,\Delta_\psi G_q
\ge
\frac{\hbar\ell\kappa}{2}
\left|\log\frac{q}{p}\right|.}
\label{eq:intro_prime_uncertainty}
\end{equation}
Thus the same arithmetic symplectic area that later appears in the twisted
M\"obius product already appears as a lower bound forced directly by the
Heisenberg uncertainty principle.

The associated Weyl operators are
\begin{equation}
W_n=\exp\left(\frac{i}{\hbar}G_n\right),
\label{eq:intro_Wn}
\end{equation}
and satisfy
\begin{equation}
W_mW_n
=
\chi_\vartheta(m,n)W_{mn},
\label{eq:intro_twisted_product}
\end{equation}
where
\begin{equation}
\chi_\vartheta(m,n)
=
\exp\left[
\frac{i\vartheta}{2}
\bigl(\Omega(m)\log n-\log m\,\Omega(n)\bigr)
\right],
\qquad
\vartheta=\frac{\ell\kappa}{\hbar}.
\label{eq:intro_cocycle}
\end{equation}
The cocycle is proved below to be cohomologically nontrivial for every
\(\vartheta\ne0\).  It induces an associative noncommutative Dirichlet
product.  M\"obius inversion in that product defines coefficients
\(z_\vartheta(n)\).  For prime powers and distinct prime pairs one obtains
\begin{align}
z_\vartheta(p^r)&=1,
\label{eq:intro_prime_power_formula}\\
z_\vartheta(pq)
&=2\cos\left(\frac{\vartheta}{2}\log\frac{q}{p}\right)-1.
\label{eq:intro_prime_pair_formula}
\end{align}
The same determinant-like arithmetic quantity therefore controls both the
uncertainty lower bound \eqref{eq:intro_arithmetic_uncertainty} and the phase
responsible for the deformed inverse.


The deformation also differentiates the twisted inverse itself at the commuting
point. For
\begin{equation}
n=\prod_{j=1}^{r}p_j^{a_j},
\qquad
T(a)=\frac{a(a+1)}2,
\label{eq:intro_factorization_T}
\end{equation}
one obtains
\begin{equation}
\boxed{
-z_0''(n)=\frac12\sum_{1\le i<j\le r}
T(a_i)T(a_j)\log^2\!\frac{p_j}{p_i}.}
\label{eq:intro_curvature_formula}
\end{equation}
For square-free \(n=p_1\cdots p_r\),
\begin{equation}
\boxed{
-z_0''(n)=\frac12\left[r\sum_{p\mid n}(\log p)^2-(\log n)^2\right]
=\frac{r^2}{2}\operatorname{Var}_{p\mid n}(\log p).}
\label{eq:intro_variance_formula}
\end{equation}
Summing this curvature does not return a two-variable M\"obius correlation.
For \(\Re s>1\),
\begin{equation}
\mathcal K(s)=-\sum_{n\ge1}\frac{z_0''(n)}{n^s}
=\frac{\zeta(s)}2\bigl(S_0(s)S_2(s)-S_1(s)^2\bigr),
\label{eq:intro_curvature_determinant}
\end{equation}
where
\begin{equation}
S_j(s)=\sum_p\frac{p^{-s}}{(1-p^{-s})^2}(\log p)^j.
\label{eq:intro_Sj}
\end{equation}
Section~\ref{sec:curvature_dirichlet_series} rewrites these functions through
\(\log\zeta(rs)\) and \(\zeta'(rs)/\zeta(rs)\), producing, after a branch is fixed by continuation from \(\Re s>1\), a local
continued germ in \(\Re s>1/2\) and an exact singularity criterion for
nontrivial zeta zeros to the right of the critical line.

Finally, a bounded self-adjoint M\"obius Hamiltonian is defined for
\(\sigma>1\).  Its Schr\"odinger evolution and the exact Heisenberg equations
for \(Q(t)\) and \(P(t)\) are derived with domains stated explicitly.  The
paper does not claim an advance toward RH.  It records exact reformulations and
local singularity statements involving the classical zero set, but no new
restriction on the location of those zeros is proved.

\begin{remark}[Scope of Version 8]
\label{rem:v3_scope}
The operator deformation is developed in the half-planes stated for each
operator. Version 8 additionally develops a direct local relation between the
second deformation curvature and the classical zero set in
Section~\ref{sec:curvature_zero_detection}.
\end{remark}



The new positivity layer in Version 9 has two branches.  First, the exact
cocyclic curvature is polarized into a Gram kernel
\(K_{\mathcal C}\succeq0\).  Second, the transverse operator
\(\Pi=i\partial_t|_{t=0}\) sends \(B(s,t)\) exactly to
\(\zeta'(s)/\zeta(s)\), eliminating every grouped dilation sector except the
primitive one.  The corresponding zero kernel is then coupled to
\(K_{\mathcal C}\) entrywise.  Positivity of the resulting kernel for the
full Weil test class is proved equivalent to RH.  The equivalence is not
advertised as progress on RH, because a single positive diagonal curvature
entry already recovers Weil's classical criterion; the open problem is to
obtain that positivity directly from the CCR/cocycle construction without
using the zero representation.



The Version 10 Mertens layer is independent of the zero-kernel reformulation.
Write
\begin{equation}
M(x)=\sum_{n\le x}\mu(n),
\qquad
A^*(x)=\max_{n\le x}|M(n)|,
\qquad
\Theta=\sup_{\zeta(\rho)=0}\Re\rho.
\label{eq:intro_record_definitions}
\end{equation}
Using the classical oscillation theorem of Kaczorowski--Pintz together with
the standard upper estimate associated with the zero supremum, Section
\ref{sec:mertens_records_heisenberg} proves
\begin{equation}
\boxed{A^*(x)=x^{\Theta+o(1)}}
\label{eq:intro_Astar_theta}
\end{equation}
and therefore, for every fixed \(0<\alpha<1\),
\begin{equation}
\boxed{A^*(x^\alpha)=o(A^*(x)).}
\label{eq:intro_scale_separation}
\end{equation}
If \(r_k\) is the first point at which \(|M|\) attains the integer height
\(k\), then
\begin{equation}
\boxed{r_k=k^{1/\Theta+o(1)}}
\label{eq:intro_record_time_law}
\end{equation}
and the full logarithmic past is negligible relative to the record height:
for every \(p>0\),
\begin{equation}
\boxed{
\frac1{\log r_k}\int_1^{r_k}
\left|\frac{M(t)}{M(r_k)}\right|^p\frac{dt}{t}\longrightarrow0.}
\label{eq:intro_record_Lp}
\end{equation}
In particular, with \(I(x)=\int_1^xM(t)dt/t\),
\begin{equation}
\boxed{I(r_k)=o(M(r_k)\log r_k).}
\label{eq:intro_record_integral}
\end{equation}

The same section constructs an auxiliary Weyl path whose ordered symplectic
area is \(\mathcal A_t(N)\).  At order zero and first transverse order one
obtains, respectively,
\begin{align}
2\mathcal A_0(N)
&=M(N)I(N)-\int_1^N\frac{M(t)^2}{t}\,dt,
\label{eq:intro_area_energy_identity}\\
\mathcal A'_0(N)
&=\frac12\sum_{m<n\le N}\mu(m)\mu(n)
\bigl(\Omega(m)\log n-\log m\,\Omega(n)\bigr).
\label{eq:intro_area_first_jet}
\end{align}
Thus the original arithmetic cocycle is exactly the first transverse jet of
an area whose order-zero term contains the logarithmic history of Mertens.
A standard Selberg--Delange input then yields the local reconstruction law
\begin{equation}
\boxed{
\Delta\mathcal A'_0(N)
=
\frac{\mu(N)N}{2\log N}+o\!\left(\frac{N}{\log N}\right).}
\label{eq:intro_first_jet_reconstruction}
\end{equation}

\section{Notation, domains, and symbol index}
\label{sec:notation}

This section fixes the notation used throughout the notebook.  Reusing a
symbol with two unrelated meanings is avoided in Version 10.

\subsection{Basic arithmetic and analytic notation}
\label{subsec:notation_arithmetic}

\begin{description}[leftmargin=3.2cm,style=nextline]
\item[\(\mu(n)\)] the M\"obius function.
\item[\(\Omega(n)\)] the completely additive number of prime factors of
\(n\), counted with multiplicity.
\item[\(\omega(n)\)] the number of distinct prime divisors of \(n\).
\item[\(v_p(n)\)] the exponent of \(p\) in \(n\).
\item[\(\operatorname{rad}(n)\)] the square-free kernel
\(\prod_{p\mid n}p\).
\item[\(\varphi(n)\)] Euler's totient and \(J_2(n)\) the second Jordan
totient.
\item[\(T(a)\)] the triangular weight \(a(a+1)/2\).
\item[\(\zeta(s)\)] the Riemann zeta function and
\(L(s)=\zeta'(s)/\zeta(s)\) its logarithmic derivative.
\item[\(P(s)\)] the prime-zeta function \(\sum_p p^{-s}\) in its initial
half-plane \(\Re s>1\), followed by analytic continuation when explicitly
stated.
\item[\(\rho\)] a nontrivial zero of \(\zeta\); \(m_\rho\) denotes its
multiplicity.
\item[\(s,t\)] Write \(s=\sigma+i\tau\).  The real variable \(t\) is reserved for the
Fourier/vertical shift in \(B(s,t)\).
\end{description}

\subsection{Operator and deformation notation}
\label{subsec:notation_operator}

\begin{description}[leftmargin=3.2cm,style=nextline]
\item[\(\Hil\)] \(L^2(\mathbb R,dq)\).
\item[\(Q,P\)] position and momentum.  Their commutator identity is asserted
on the Schwartz core \(\Sch(\mathbb R)\), not on the whole intersection of
maximal domains.
\item[\(W(x,y)\)] the Weyl displacement and \(W_n\) its arithmetic
specialization.
\item[\(\vartheta\)] the dimensionless deformation parameter
\(\ell\kappa/\hbar\).
\item[\(\chi_\vartheta(m,n)\)] the arithmetic Weyl 2-cocycle.
\item[\(\star_\vartheta\)] twisted Dirichlet convolution.
\item[\(z_\vartheta(n)\)] the coefficientwise twisted inverse of \(\mu\).
\item[\(\mathcal C(n)\)] the deformation curvature \(-z_0''(n)\).
\item[\(\mathcal K(s)\)] the Dirichlet series of \(\mathcal C(n)\).
\end{description}

\subsection{Euler-locality and replica notation}
\label{subsec:notation_replica}

\begin{description}[leftmargin=3.2cm,style=nextline]
\item[\(E_\vartheta\)] the increasing-prime ordered Euler object.
\item[\(R_\vartheta\)] the Euler-locality defect \(z_\vartheta-E_\vartheta\).
\item[\(\Phi_{\rm ord}(n)\)] the oriented prime-pair phase of the ordered
Euler object.
\item[\(\mathcal O(s)\)] the odd Euler-locality Dirichlet defect, initially
holomorphic for \(\Re s>1\).
\item[\(u_p(s)\)] \(p^{-s}/(1-p^{-s})=1/(p^s-1)\).
\item[\(A(s)\)] \(\sum_pu_p(s)=B(s,0)\), initially for \(\Re s>1\).
\item[\(B(s,t)\)] \(\sum_pu_p(s)p^{it}\), initially for \(\Re s>1\).
\item[\(\mathcal I(s)\)] the normalized Fourier energy
\(2\pi\mathcal O(s)/\zeta(s)\) in \(\Re s>1\).  Its continued germ is the
object whose leading local term is naturally of
\(z^{-1}\log z\)-type.
\item[\(s_0\)] a fixed point in the strip, and \(z=s-s_0\) the local
coordinate.
\item[\(\mathcal R(s_0)\)] the finite set of dilation indices
\(N\ge1\) for which \(\zeta(Ns_0)=0\).
\item[\(q_N\)] the multiplicity of the zero \(Ns_0\), for
\(N\in\mathcal R(s_0)\).
\item[\(\beta_r\)] the strictly negative divisor energy attached to a
single pure replica of index \(r>1\).
\item[\(\Phi_{pq}(n)\)] the prime-pair curvature feature used to polarize \(\mathcal C(n)\).
\item[\(K_{\mathcal C}(m,n)\)] the positive-semidefinite curvature Gram kernel.
\item[\(\Pi\)] the transverse primitive-replica filter \(i\,\partial_t|_{t=0}\).
\item[\(G^\#(s)\)] the Weil involution \(\overline{G(1-\overline{s})}\).
\item[\(W_g(s)\)] the Weil weight \(G(s)G^\#(s)\) attached to a test function \(g\).
\item[\(K_g^Z(m,n)\)] the zero kernel built from primitive zeta zeros and the Weil weight.
\item[\(\mathscr K_g(m,n)\)] the entrywise product \(K_{\mathcal C}(m,n)K_g^Z(m,n)\).
\end{description}



\subsection{Version 12 squarefree-deformation notation}
\label{subsec:notation_v11}

\begin{description}[leftmargin=3.2cm,style=nextline]
\item[\(P_x(z)\)] the finite squarefree generating polynomial
\(\sum_{n\le x}\mu(n)^2z^{\Omega(n)}\).
\item[\(F_t(x)\)] the signed deformation
\(P_x(-e^t)=\sum_{n\le x}\mu(n)e^{t\Omega(n)}\).
\item[\(\xi_*(x)\)] the unique zero of \(P_x\) converging to \(-1\).
\item[\(h_x\)] the critical displacement \(1+\xi_*(x)\).
\item[\(L_x\)] \(\log\log x\) in the critical-scale subsections.
\item[\(M_j(x)\)] \(\sum_{n\le x}\mu(n)\Omega(n)^j\), the \(j\)-th
transverse moment.
\item[\(\Phi_x(u)\)] the critically renormalized deformation
\((\log x)^2L_xF_{u/L_x}(x)/x\).
\item[\(T_m(x)\)] the root trace
\(\sum_j(1+\xi_{x,j})^{-m}\), defined when \(M(x)\ne0\).
\end{description}

\begin{warning}[Normalized energy versus \(\mathcal O\)]
\label{warn:I_vs_O_notation}
The logarithmic-over-pole cancellation is a statement about
\(\mathcal I=2\pi\mathcal O/\zeta\).  At an actual zero \(s=\rho\), the
factor \(\zeta(s)\) lowers the singular order of \(\mathcal O\).  Thus a
formula containing \(\log(s-\rho)/(s-\rho)\) must not be attributed to
\(\mathcal O\) itself at \(r=1\).
\end{warning}

\section{Exact canonical kinematics}
\label{sec:kinematics}

\subsection{Hilbert space, position, and momentum}
\label{subsec:QP_definitions}

\begin{definition}[Hilbert space]
\label{def:hilbert_space}
Fix $\hbar>0$ and set
\begin{equation}
\Hil=L^2(\R,dq)
\label{eq:hilbert_space}
\end{equation}
with its standard complex Hilbert-space inner product.
\end{definition}

\begin{definition}[Position]
\label{def:position}
The position operator $Q$ is
\begin{equation}
(Q\psi)(q)=q\psi(q)
\label{eq:Q_action}
\end{equation}
with maximal domain
\begin{equation}
\Dom(Q)=\{\psi\in L^2(\R):q\psi(q)\in L^2(\R)\}.
\label{eq:Q_domain}
\end{equation}
\end{definition}

\begin{definition}[Momentum]
\label{def:momentum}
The momentum operator $P$ is
\begin{equation}
(P\psi)(q)=-i\hbar\,\psi'(q)
\label{eq:P_action}
\end{equation}
with domain
\begin{equation}
\Dom(P)=H^1(\R).
\label{eq:P_domain}
\end{equation}
\end{definition}

\begin{proposition}[Self-adjointness and common core]
\label{prop:self_adjoint_QP}
The operators $Q$ and $P$ are self-adjoint.  The Schwartz space
$\Sch(\R)$ is a common invariant core for both.
\end{proposition}

\begin{proof}
$Q$ is the maximal multiplication operator by the real function $q$.  Under
the unitary Fourier transform with convention
\begin{equation}
(\Four\psi)(p)
=
\frac{1}{\sqrt{2\pi\hbar}}
\int_{\R}e^{-ipq/\hbar}\psi(q)\,dq,
\label{eq:fourier_transform}
\end{equation}
$P$ is unitarily equivalent to multiplication by the real variable $p$.
The core statement follows by standard cutoff and mollification arguments.
See, for example, \cite{ReedSimon1980}.
\end{proof}

\begin{proposition}[Exact canonical commutation relation]
\label{prop:CCR}
For every $\psi\in\Sch(\R)$,
\begin{equation}
[Q,P]\psi=i\hbar\psi.
\label{eq:CCR}
\end{equation}
\end{proposition}

\begin{proof}
Direct differentiation gives
\begin{align}
(QP\psi)(q)
&=-i\hbar q\psi'(q),
\label{eq:QP}\\
(PQ\psi)(q)
&=-i\hbar\psi(q)-i\hbar q\psi'(q).
\label{eq:PQ}
\end{align}
Subtracting \eqref{eq:PQ} from \eqref{eq:QP} proves \eqref{eq:CCR}.
\end{proof}

\subsection{The Heisenberg uncertainty principle and exact incompatibility}
\label{subsec:heisenberg_uncertainty}

\begin{definition}[Expectation and uncertainty]
\label{def:uncertainty}
Let \(A\) be self-adjoint and let \(\psi\in\Dom(A)\) be normalized.  Define
\begin{align}
\langle A\rangle_\psi
&=\langle\psi,A\psi\rangle,
\label{eq:expectation_definition}\\
(\Delta_\psi A)^2
&=\|(A-\langle A\rangle_\psi)\psi\|^2.
\label{eq:variance_definition}
\end{align}
\end{definition}

\begin{theorem}[Heisenberg uncertainty principle]
\label{thm:heisenberg_uncertainty}
For every normalized \(\psi\in\Sch(\R)\),
\begin{equation}
\boxed{
\Delta_\psi Q\,\Delta_\psi P
\ge
\frac{\hbar}{2}.}
\label{eq:heisenberg_uncertainty}
\end{equation}
\end{theorem}

\begin{proof}
Set
\begin{equation}
\widetilde Q=Q-\langle Q\rangle_\psi I,
\qquad
\widetilde P=P-\langle P\rangle_\psi I.
\label{eq:centered_QP}
\end{equation}
By Cauchy--Schwarz,
\begin{equation}
\Delta_\psi Q\,\Delta_\psi P
\ge
\left|\langle\widetilde Q\psi,\widetilde P\psi\rangle\right|.
\label{eq:CS_uncertainty}
\end{equation}
The modulus dominates the absolute value of the imaginary part, and on
\(\Sch(\R)\)
\begin{align}
2i\,\operatorname{Im}
\langle\widetilde Q\psi,\widetilde P\psi\rangle
&=
\langle\psi,[Q,P]\psi\rangle
\label{eq:imaginary_commutator}\\
&=i\hbar.
\label{eq:imaginary_CCR}
\end{align}
Combining \eqref{eq:CS_uncertainty}--\eqref{eq:imaginary_CCR} proves
\eqref{eq:heisenberg_uncertainty}.
\end{proof}

\begin{remark}[Meaning of the inequality]
\label{rem:uncertainty_meaning}
Equation~\eqref{eq:heisenberg_uncertainty} is the Heisenberg uncertainty
principle used in this paper.  The quantities \(\Delta_\psi Q\) and
\(\Delta_\psi P\) are intrinsic standard deviations of a state.  They are
not measurement errors.  Exact simultaneous localization would require both
to vanish, which is prohibited by \eqref{eq:heisenberg_uncertainty}.
\end{remark}

\begin{theorem}[No simultaneous exact position and momentum]
\label{thm:no_joint_eigenstate}
There is no nonzero state with simultaneous exact position and exact momentum.
More precisely, both of the following are impossible for real \(q_0,p_0\):
\begin{enumerate}[leftmargin=*]
\item a nonzero vector \(\psi\) in a common invariant domain satisfying
\begin{equation}
Q\psi=q_0\psi,
\qquad
P\psi=p_0\psi;
\label{eq:joint_eigenstate_assumption}
\end{equation}
\item a nonzero tempered distribution \(T\in\Sch'(\R)\) satisfying
\begin{equation}
qT=q_0T,
\qquad
-i\hbar T'=p_0T.
\label{eq:joint_distribution_assumption}
\end{equation}
\end{enumerate}
\end{theorem}

\begin{proof}
For vectors, \eqref{eq:joint_eigenstate_assumption} gives
\([Q,P]\psi=0\), whereas the CCR gives \([Q,P]\psi=i\hbar\psi\); hence
\(\psi=0\).  For distributions define directly
\begin{equation}
QT=qT,
\qquad
PT=-i\hbar T'.
\label{eq:distribution_QP_definition}
\end{equation}
The distributional Leibniz rule gives
\begin{equation}
[Q,P]T=i\hbar T.
\label{eq:distribution_CCR}
\end{equation}
The simultaneous eigenvalue equations would make the left-hand side zero,
so \(T=0\).
\end{proof}

\section{Weyl displacements and the arithmetic phase space}
\label{sec:weyl_arithmetic}

\subsection{Weyl operators}
\label{subsec:weyl_operators}

\begin{definition}[Weyl displacement]
\label{def:weyl_operator}
For $(x,y)\in\R^2$, define
\begin{equation}
W(x,y)
=
\exp\!\left(\frac{i}{\hbar}(yQ-xP)\right).
\label{eq:weyl_definition}
\end{equation}
This is a unitary operator on all of $\Hil$.
\end{definition}

\begin{proposition}[Coordinate action]
\label{prop:weyl_coordinate_action}
For $\psi\in L^2(\R)$,
\begin{equation}
(W(x,y)\psi)(q)
=
\exp\!\left(\frac{i}{\hbar}y\left(q-\frac{x}{2}\right)\right)
\psi(q-x)
\label{eq:weyl_coordinate_action}
\end{equation}
almost everywhere.
\end{proposition}

\begin{proof}
Since $[iyQ/\hbar,-ixP/\hbar]$ is scalar, the Baker--Campbell--Hausdorff
formula gives
\begin{equation}
W(x,y)
=
e^{-ixy/(2\hbar)}e^{iyQ/\hbar}e^{-ixP/\hbar}.
\label{eq:weyl_BCH_factorization}
\end{equation}
The last factor translates by $x$ and the middle factor multiplies by
$e^{iyq/\hbar}$, yielding \eqref{eq:weyl_coordinate_action}.
\end{proof}

\begin{proposition}[Weyl form of the CCR]
\label{prop:weyl_relation}
Let $z=(x,y)$ and $z'=(x',y')$.  Define
\begin{equation}
\omega(z,z')=yx'-xy'.
\label{eq:symplectic_form}
\end{equation}
Then
\begin{equation}
W(z)W(z')
=
\exp\!\left(\frac{i}{2\hbar}\omega(z,z')\right)
W(z+z').
\label{eq:weyl_product}
\end{equation}
Consequently,
\begin{equation}
[W(z),W(z')]
=
2i\sin\!\left(\frac{\omega(z,z')}{2\hbar}\right)W(z+z').
\label{eq:weyl_commutator}
\end{equation}
\end{proposition}

\begin{proof}
For
\begin{equation}
A_z=\frac{i}{\hbar}(yQ-xP),
\label{eq:Az_definition}
\end{equation}
the CCR gives
\begin{equation}
[A_z,A_{z'}]
=
\frac{i}{\hbar}\omega(z,z')I.
\label{eq:Az_commutator}
\end{equation}
This commutator is central, so the Baker--Campbell--Hausdorff series stops
after the commutator term and gives \eqref{eq:weyl_product}.  Reversing $z$
and $z'$ changes the sign of $\omega$, giving \eqref{eq:weyl_commutator}.
The Weyl relations are standard; see \cite{Folland1989}.
\end{proof}

\begin{proposition}[Position and momentum displacement equations]
\label{prop:weyl_QP_commutators}
On $\Sch(\R)$,
\begin{align}
[Q,W(x,y)]
&=xW(x,y),
\label{eq:Q_W_commutator}\\
[P,W(x,y)]
&=yW(x,y).
\label{eq:P_W_commutator}
\end{align}
\end{proposition}

\begin{proof}
Both identities follow by differentiating \eqref{eq:weyl_coordinate_action},
or directly from \eqref{eq:CCR} because the relevant commutators with the
exponent are scalar.
\end{proof}

\subsection{Arithmetic embedding}
\label{subsec:arithmetic_embedding}

\begin{definition}[The additive prime-factor counter]
\label{def:Omega}
For
\begin{equation}
n=\prod_p p^{v_p(n)},
\label{eq:prime_factorization}
\end{equation}
define
\begin{equation}
\Om(n)=\sum_p v_p(n).
\label{eq:Omega_definition}
\end{equation}
Then
\begin{equation}
\Om(mn)=\Om(m)+\Om(n).
\label{eq:Omega_additivity}
\end{equation}
\end{definition}

\begin{definition}[Arithmetic phase-space map]
\label{def:arithmetic_phase_space_map}
Fix two positive scales $\ell>0$ and $\kappa>0$.  For $n\in\N$, put
\begin{align}
x_n
&=\ell\log n,
\label{eq:xn_definition}\\
y_n
&=\kappa\Om(n),
\label{eq:yn_definition}\\
z_n
&=(x_n,y_n).
\label{eq:zn_definition}
\end{align}
Define the dimensionless Heisenberg area parameter
\begin{equation}
\vartheta=\frac{\ell\kappa}{\hbar}.
\label{eq:theta_definition}
\end{equation}
\end{definition}

\begin{remark}[Why these two coordinates]
\label{rem:why_coordinates}
The first coordinate is forced by multiplicative size because
$\log(mn)=\log m+\log n$.  The second coordinate is chosen as the simplest
integer-valued completely additive function that directly controls the
M\"obius sign on square-free integers:
\begin{equation}
\mu(n)\neq0
\quad\Longrightarrow\quad
\mu(n)=(-1)^{\Om(n)}.
\label{eq:mobius_Omega_relation}
\end{equation}
The construction can be generalized by replacing $\Om$ with another
completely additive arithmetic function, but the present notebook fixes $\Om$ throughout.
\end{remark}

\begin{definition}[Arithmetic Weyl operators]
\label{def:arithmetic_weyl_operator}
For $n\in\N$, define
\begin{equation}
W_n=W(x_n,y_n).
\label{eq:Wn_definition}
\end{equation}
\end{definition}

\begin{theorem}[Exact Heisenberg cocycle on the integers]
\label{thm:integer_weyl_cocycle}
For $m,n\in\N$,
\begin{equation}
W_mW_n
=
\chi_{\vartheta}(m,n)W_{mn},
\label{eq:integer_weyl_product}
\end{equation}
where
\begin{equation}
\chi_{\vartheta}(m,n)
=
\exp\!\left[
\frac{i\vartheta}{2}
\bigl(\Om(m)\log n-\log m\,\Om(n)\bigr)
\right].
\label{eq:chi_definition}
\end{equation}
Moreover,
\begin{equation}
[W_m,W_n]
=
2i\sin\!\left[
\frac{\vartheta}{2}
\bigl(\Om(m)\log n-\log m\,\Om(n)\bigr)
\right]W_{mn}.
\label{eq:integer_weyl_commutator}
\end{equation}
\end{theorem}

\begin{proof}
Complete additivity gives
\begin{equation}
z_{mn}=z_m+z_n.
\label{eq:phase_space_additivity}
\end{equation}
Substitution of \eqref{eq:xn_definition}--\eqref{eq:theta_definition} into
\eqref{eq:weyl_product} gives \eqref{eq:integer_weyl_product};
\eqref{eq:integer_weyl_commutator} follows from
\eqref{eq:weyl_commutator}.
\end{proof}

\begin{corollary}[Prime commutator]
\label{cor:prime_commutator}
For distinct primes $p,q$,
\begin{equation}
[W_p,W_q]
=
2i\sin\!\left(\frac{\vartheta}{2}\log\frac{q}{p}\right)W_{pq}.
\label{eq:prime_commutator}
\end{equation}
Thus the prime operators commute only at the resonant values
\begin{equation}
\vartheta\log\frac{q}{p}\in2\pi\mathbb Z.
\label{eq:prime_commutation_resonance}
\end{equation}
\end{corollary}

\begin{proof}
For primes \(p,q\), one has \(\Omega(p)=\Omega(q)=1\).  Substituting these
values into the integer Weyl commutator formula of
Theorem~\ref{thm:integer_weyl_cocycle} gives
\eqref{eq:prime_commutator}.  The commutator vanishes exactly when its sine
factor vanishes, which is equivalent to
\eqref{eq:prime_commutation_resonance}.
\end{proof}

\begin{remark}[Where Heisenberg enters nontrivially]
\label{rem:where_heisenberg_enters}
Equation~\eqref{eq:integer_weyl_product} is the first point where Version 2
ceases to be ordinary Dirichlet algebra in operator notation.  If the
canonical pair were commuting, the Weyl cocycle would be $1$ and the
arithmetic product below would reduce to ordinary Dirichlet convolution.
For nonzero $\vartheta$, the CCR creates order-dependent phases between
different arithmetic factors.
\end{remark}

\begin{theorem}[The arithmetic cocycle is not gauge-trivial]
\label{thm:cocycle_nontrivial}
Extend \(\log\) and \(\Omega\) additively to
\(\mathbb Q_{>0}^{\times}\).  For every real \(\vartheta\ne0\), the
cocycle \(\chi_\vartheta\) is not a scalar coboundary on
\(\mathbb N^{\times}\), and its extension is not a scalar coboundary on
\(\mathbb Q_{>0}^{\times}\).
\end{theorem}

\begin{proof}
On a commutative semigroup, every scalar coboundary
\begin{equation}
(\delta a)(m,n)=a(m)a(n)a(mn)^{-1}
\label{eq:scalar_coboundary}
\end{equation}
is symmetric in \(m,n\).  Hence a necessary condition for
\(\chi_\vartheta\) to be a coboundary is
\begin{equation}
\frac{\chi_\vartheta(m,n)}{\chi_\vartheta(n,m)}=1
\label{eq:cocycle_commutator_trivial}
\end{equation}
for all \(m,n\).  For distinct primes,
\begin{equation}
\frac{\chi_\vartheta(p,q)}{\chi_\vartheta(q,p)}
=
\exp\left(i\vartheta\log\frac{q}{p}\right).
\label{eq:prime_cocycle_commutator}
\end{equation}
If this were one simultaneously for \((p,q)=(2,3)\) and \((2,5)\), then
there would be nonzero integers \(k,l\) such that
\begin{equation}
\vartheta\log(3/2)=2\pi k,
\qquad
\vartheta\log(5/2)=2\pi l.
\label{eq:two_prime_phase_integrality}
\end{equation}
Thus \(\log(3/2)/\log(5/2)=k/l\in\mathbb Q\).  Exponentiating a
reduced rational relation would produce an equality between incompatible
prime factorizations, contradicting unique factorization.  Therefore
\eqref{eq:cocycle_commutator_trivial} cannot hold for all primes.
\end{proof}


\section{Structural filter: linear M\"obius states cannot escape second order}
\label{sec:linear_state_filter}

Let \(v_n\in\Hil\) be arbitrary vectors and define
\begin{equation}
\Psi_x=\sum_{n\le x}\mu(n)v_n.
\label{eq:linear_mobius_state}
\end{equation}

\begin{theorem}[Quadraticity of unnormalized sesquilinear moments]
\label{thm:linear_state_quadratic_obstruction}
For every linear operator \(A\) for which the expression is defined,
\begin{equation}
\langle\Psi_x,A\Psi_x\rangle
=\sum_{m,n\le x}\mu(m)\mu(n)\langle v_m,Av_n\rangle.
\label{eq:expectation_bilinear_mobius}
\end{equation}
Thus every unnormalized sesquilinear moment, including
\(\|\Psi_x\|^2\), is quadratic in the M\"obius coefficients.
Normalized expectations are quotients of such quadratic forms, and variances
or Robertson expressions are rational combinations of them; they are not,
in general, themselves quadratic forms.
\end{theorem}

\begin{proof}
Expand the inner product sesquilinearly and use that \(\mu(n)\) is real.
The assertions about normalized quantities follow directly by dividing by
\(\langle\Psi_x,\Psi_x\rangle\).
\end{proof}

\begin{remark}[Scope of the linear-state filter]
\label{rem:linear_state_filter_scope}
Equation~\eqref{eq:expectation_bilinear_mobius} is a structural warning, not a
no-go theorem for every nonlinear observable constructed from a linear
M\"obius state.  What is ruled out is only the claim that an unnormalized
sesquilinear moment escapes second order in the coefficients.
\end{remark}

\section{Arithmetic uncertainty generated by the CCR}
\label{sec:arithmetic_uncertainty}

\begin{definition}[Integer generator]
\label{def:integer_generator}
For \(n\ge1\), define on \(\Sch(\R)\) the symmetric operator
\begin{equation}
G_n^{(0)}=y_nQ-x_nP
=\kappa\Omega(n)Q-\ell\log n\,P.
\label{eq:Gn_definition}
\end{equation}
It is essentially self-adjoint on \(\Sch(\R)\); denote its self-adjoint
closure by \(G_n\).  The corresponding unitary displacement is
\begin{equation}
W_n=\exp(iG_n/\hbar).
\label{eq:Wn_from_Gn}
\end{equation}
\end{definition}

\begin{proposition}[Arithmetic commutator]
\label{prop:arithmetic_generator_commutator}
For all \(m,n\ge1\), the restrictions of \(G_m\) and \(G_n\) to \(\Sch(\R)\) satisfy
\begin{align}
[G_m,G_n]
&=i\hbar(x_my_n-y_mx_n)I
\label{eq:GmGn_commutator_xy}\\
&=i\hbar\ell\kappa
\bigl(\log m\,\Omega(n)-\Omega(m)\log n\bigr)I.
\label{eq:GmGn_commutator_arithmetic}
\end{align}
\end{proposition}

\begin{proof}
Expand the commutator and use \([Q,P]=i\hbar I\).  Terms involving
\([Q,Q]\) and \([P,P]\) vanish.
\end{proof}

\begin{theorem}[Arithmetic Heisenberg uncertainty law]
\label{thm:arithmetic_uncertainty}
For every normalized \(\psi\in\Sch(\R)\) and all \(m,n\ge1\),
\begin{equation}
\boxed{
\Delta_\psi G_m\,\Delta_\psi G_n
\ge
\frac{\hbar\ell\kappa}{2}
\left|\Omega(m)\log n-\log m\,\Omega(n)\right|.}
\label{eq:arithmetic_uncertainty}
\end{equation}
For distinct primes \(p,q\),
\begin{equation}
\boxed{
\Delta_\psi G_p\,\Delta_\psi G_q
\ge
\frac{\hbar\ell\kappa}{2}
\left|\log\frac{q}{p}\right|.}
\label{eq:prime_pair_uncertainty}
\end{equation}
For two powers of the same prime the lower bound is zero.
\end{theorem}

\begin{proof}
Apply the Robertson inequality to the self-adjoint pair \(G_m,G_n\):
\begin{equation}
\Delta_\psi G_m\,\Delta_\psi G_n
\ge
\frac12
\left|\langle\psi,[G_m,G_n]\psi\rangle\right|.
\label{eq:Robertson_GmGn}
\end{equation}
Substitution of \eqref{eq:GmGn_commutator_arithmetic} gives
\eqref{eq:arithmetic_uncertainty}.  If \(p,q\) are distinct primes then
\(\Omega(p)=\Omega(q)=1\), which gives \eqref{eq:prime_pair_uncertainty}.
If \(m=p^a,n=p^b\), the determinant
\(\Omega(m)\log n-\log m\,\Omega(n)\) vanishes.
\end{proof}

\begin{remark}[The common invariant]
\label{rem:common_symplectic_invariant}
Define the arithmetic symplectic area
\begin{equation}
\mathcal A(m,n)
=
\Omega(m)\log n-\log m\,\Omega(n).
\label{eq:arithmetic_symplectic_area}
\end{equation}
The uncertainty lower bound is proportional to \( |\mathcal A(m,n)| \),
while the twisted product below carries the phase
\(\exp(i\vartheta\mathcal A(m,n)/2)\).  Thus uncertainty and twisted
M\"obius inversion are controlled by the same arithmetic invariant, although
they are distinct consequences of the CCR.
\end{remark}

\subsection{The semiclassical and arithmetic-commutative limits are distinct}
\label{subsec:two_limits}

The parameter controlling the Weyl cocycle is
\begin{equation}
\vartheta(\hbar,\ell,\kappa)
=
\frac{\ell\kappa}{\hbar}.
\label{eq:vartheta_scaling}
\end{equation}
The uncertainty scale and the deformation scale therefore depend on the same
three parameters in different ways.

\begin{proposition}[Opposite fixed-scale limits]
\label{prop:opposite_limits}
Fix \(\ell>0\) and \(\kappa>0\).  For every fixed pair \(m,n\),
\begin{equation}
\frac{\hbar\ell\kappa}{2}|\mathcal A(m,n)|
\longrightarrow0
\qquad(\hbar\to0^+),
\label{eq:uncertainty_classical_limit}
\end{equation}
whereas
\begin{equation}
|\vartheta|
=
\frac{\ell\kappa}{\hbar}
\longrightarrow\infty.
\label{eq:twist_fixed_scale_limit}
\end{equation}
Consequently the ordinary semiclassical limit \(\hbar\to0^+\) with
\(\ell,\kappa\) fixed does not approach the untwisted Dirichlet product.
\end{proposition}

\begin{proof}
Equation~\eqref{eq:uncertainty_classical_limit} follows directly from
Theorem~\ref{thm:arithmetic_uncertainty}, while
\eqref{eq:twist_fixed_scale_limit} follows from
\eqref{eq:vartheta_scaling}.
\end{proof}

\begin{proposition}[Joint semiclassical--commutative scaling]
\label{prop:joint_scaling}
Let \(\ell=\ell(\hbar)>0\) and
\(\kappa=\kappa(\hbar)>0\).  Then
\begin{equation}
\vartheta(\hbar)\longrightarrow0
\qquad(\hbar\to0^+)
\label{eq:theta_to_zero_condition}
\end{equation}
if and only if
\begin{equation}
\ell(\hbar)\kappa(\hbar)=o(\hbar).
\label{eq:joint_scaling_condition}
\end{equation}
Under this scaling the twisted product converges coefficientwise to the
ordinary Dirichlet product, and the arithmetic uncertainty lower bound obeys
\begin{equation}
\frac{\hbar\ell(\hbar)\kappa(\hbar)}{2}
|\mathcal A(m,n)|
=o(\hbar^2)
\label{eq:joint_uncertainty_decay}
\end{equation}
for every fixed \(m,n\).
\end{proposition}

\begin{proof}
The equivalence of \eqref{eq:theta_to_zero_condition} and
\eqref{eq:joint_scaling_condition} is immediate from
\eqref{eq:vartheta_scaling}.  For fixed \(m,n\),
\(\chi_{\vartheta}(m,n)\to1\) as \(\vartheta\to0\), so the
finite divisor sum defining \(f\star_{\vartheta}g\) converges to ordinary
Dirichlet convolution coefficientwise.  Multiplying
\(\ell(\hbar)\kappa(\hbar)=o(\hbar)\) by \(\hbar/2\) gives
\eqref{eq:joint_uncertainty_decay}.
\end{proof}

\begin{remark}[Two independent notions of classicality]
\label{rem:two_classicalities}
The limit \(\hbar\to0^+\) suppresses the Robertson--Heisenberg lower
bound, whereas \(\vartheta\to0\) removes the arithmetic twist.  They are
not the same operation.  With fixed \(\ell,\kappa\) they move in opposite
directions; simultaneous suppression requires the joint scaling
\eqref{eq:joint_scaling_condition}.
\end{remark}

\section{The Heisenberg-twisted Dirichlet algebra}
\label{sec:twisted_dirichlet_algebra}

\subsection{Definition and associativity}
\label{subsec:twisted_convolution_definition}

\begin{definition}[Heisenberg-twisted Dirichlet convolution]
\label{def:twisted_convolution}
For arithmetic functions $f,g:\N\to\C$, define
\begin{equation}
(f\twconv g)(k)
=
\sum_{d\mid k}
f(d)g(k/d)\chi_{\vartheta}(d,k/d).
\label{eq:twisted_convolution}
\end{equation}
The sum is finite for every $k$.
\end{definition}

\begin{lemma}[Cocycle identity]
\label{lem:cocycle_identity}
For all $a,b,c\in\N$,
\begin{equation}
\chi_{\vartheta}(a,b)\chi_{\vartheta}(ab,c)
=
\chi_{\vartheta}(b,c)\chi_{\vartheta}(a,bc).
\label{eq:cocycle_identity}
\end{equation}
\end{lemma}

\begin{proof}
The exponent in \eqref{eq:chi_definition} is the alternating bilinear form
obtained from the additive map $n\mapsto z_n$.  Expanding both sides using
$\log(ab)=\log a+\log b$ and
$\Om(ab)=\Om(a)+\Om(b)$ gives the same exponent.
\end{proof}

\begin{proposition}[Associative noncommutative arithmetic algebra]
\label{prop:twisted_associative}
The product $\twconv$ is associative and has identity $\deltaone$, where
\begin{equation}
\deltaone(1)=1,
\qquad
\deltaone(n)=0\quad(n>1).
\label{eq:delta_one_definition}
\end{equation}
For $\vartheta\neq0$ it is generally noncommutative.
\end{proposition}

\begin{proof}
Associativity follows from Lemma~\ref{lem:cocycle_identity} after grouping
the finite divisor sums by triples $abc=k$.  The identity property follows
from $\chi_{\vartheta}(1,n)=\chi_{\vartheta}(n,1)=1$.  Noncommutativity
follows from Corollary~\ref{cor:prime_commutator}; equivalently,
\begin{equation}
(\delta_p\twconv\delta_q)(pq)
=
\chi_{\vartheta}(p,q)
\label{eq:delta_pq_forward}
\end{equation}
whereas
\begin{equation}
(\delta_q\twconv\delta_p)(pq)
=
\chi_{\vartheta}(q,p)
=
\overline{\chi_{\vartheta}(p,q)}.
\label{eq:delta_pq_reverse}
\end{equation}
\end{proof}

\subsection{Weighted Banach algebra and operator representation}
\label{subsec:weighted_banach_algebra}

\begin{definition}[Weighted arithmetic Banach space]
\label{def:weighted_banach_algebra}
For $\sigma>0$, set
\begin{equation}
\mathcal A_{\sigma}
=
\left\{f:\N\to\C:
\|f\|_{\sigma}:=
\sum_{n\ge1}\frac{|f(n)|}{n^{\sigma}}<\infty
\right\}.
\label{eq:A_sigma_definition}
\end{equation}
\end{definition}

\begin{proposition}[Banach-algebra estimate]
\label{prop:banach_algebra_estimate}
For $f,g\in\mathcal A_{\sigma}$,
\begin{equation}
\|f\twconv g\|_{\sigma}
\le
\|f\|_{\sigma}\|g\|_{\sigma}.
\label{eq:banach_algebra_estimate}
\end{equation}
Hence $(\mathcal A_{\sigma},\twconv)$ is a unital Banach algebra.
\end{proposition}

\begin{proof}
Since $|\chi_{\vartheta}|=1$,
\begin{align}
\|f\twconv g\|_{\sigma}
&\le
\sum_{k\ge1}\frac{1}{k^{\sigma}}
\sum_{d\mid k}|f(d)||g(k/d)|
\label{eq:banach_estimate_step1}\\
&=
\sum_{d,e\ge1}
\frac{|f(d)|}{d^{\sigma}}
\frac{|g(e)|}{e^{\sigma}}
\label{eq:banach_estimate_step2}\\
&=
\|f\|_{\sigma}\|g\|_{\sigma}.
\label{eq:banach_estimate_step3}
\end{align}
Completeness is the completeness of weighted $\ell^1$.
\end{proof}

\begin{definition}[Weyl--Dirichlet representation]
\label{def:Weyl_Dirichlet_representation}
For $f\in\mathcal A_{\sigma}$, define
\begin{equation}
\Pi_{\sigma,\vartheta}(f)
=
\sum_{n\ge1}\frac{f(n)}{n^{\sigma}}W_n,
\label{eq:Pi_definition}
\end{equation}
where the series converges absolutely in operator norm.
\end{definition}

\begin{proposition}[Representation theorem]
\label{prop:representation_theorem}
For $f,g\in\mathcal A_{\sigma}$,
\begin{align}
\|\Pi_{\sigma,\vartheta}(f)\|
&\le\|f\|_{\sigma},
\label{eq:Pi_norm_bound}\\
\Pi_{\sigma,\vartheta}(f)
\Pi_{\sigma,\vartheta}(g)
&=
\Pi_{\sigma,\vartheta}(f\twconv g).
\label{eq:Pi_homomorphism}
\end{align}
\end{proposition}

\begin{proof}
The norm bound follows from $\|W_n\|=1$.  Absolute norm convergence permits
multiplication and regrouping of the double series.  The coefficient of
$W_k$ is exactly \eqref{eq:twisted_convolution} by
Theorem~\ref{thm:integer_weyl_cocycle}.
\end{proof}

\section{M\"obius inversion deformed by the CCR}
\label{sec:mobius_inversion}

\subsection{The M\"obius operator}
\label{subsec:mobius_operator}

\begin{definition}[M\"obius operator]
\label{def:mobius_operator}
For $\sigma>1$, define
\begin{equation}
\mathcal M_{\sigma,\vartheta}
=
\Pi_{\sigma,\vartheta}(\mu)
=
\sum_{n\ge1}
\frac{\mu(n)}{n^{\sigma}}W_n.
\label{eq:mobius_operator}
\end{equation}
The series converges absolutely in operator norm, so
$\mathcal M_{\sigma,\vartheta}\in\Bop(\Hil)$ and
\begin{equation}
\Dom(\mathcal M_{\sigma,\vartheta})=\Hil.
\label{eq:mobius_operator_domain}
\end{equation}
\end{definition}

\begin{remark}[Concrete coordinate action]
\label{rem:mobius_coordinate_action}
For $\psi\in L^2(\R)$,
\begin{equation}
(\mathcal M_{\sigma,\vartheta}\psi)(q)
=
\sum_{n\ge1}
\frac{\mu(n)}{n^{\sigma}}
\exp\!\left[
\frac{i}{\hbar}y_n\left(q-\frac{x_n}{2}\right)
\right]
\psi(q-x_n),
\label{eq:mobius_coordinate_action}
\end{equation}
with convergence in $L^2$.  Thus the operator is a M\"obius-weighted
superposition of simultaneous translations in position and momentum; it is
not a Fourier multiplier when $\kappa\neq0$.
\end{remark}

\subsection{Formal twisted inverse}
\label{subsec:formal_twisted_inverse}

\begin{proposition}[Existence and uniqueness of arithmetic inverses]
\label{prop:formal_inverse}
Let $f:\N\to\C$ satisfy $f(1)\neq0$.  Then there exists a unique arithmetic
function $g$ such that
\begin{equation}
f\twconv g=g\twconv f=\deltaone.
\label{eq:two_sided_inverse}
\end{equation}
Its coefficients are determined recursively from $g(1)=f(1)^{-1}$.
\end{proposition}

\begin{proof}
For $n>1$, the equation $(f\twconv r)(n)=0$ contains the term
$f(1)r(n)$ and otherwise only values $r(d)$ with $d<n$.  It therefore
constructs a unique right inverse $r$ satisfying $f\twconv r=\deltaone$.
Separately, the equations $(\ellft\twconv f)(n)=0$ construct a unique left
inverse $\ellft$ satisfying $\ellft\twconv f=\deltaone$.  Associativity
then gives
\begin{equation}
\ellft
=\ellft\twconv(f\twconv r)
=(\ellft\twconv f)\twconv r
=r,
\label{eq:left_right_inverse_coincide}
\end{equation}
so the inverse is bilateral and unique.
\end{proof}

\begin{definition}[Heisenberg-zeta coefficients]
\label{def:heisenberg_zeta_coefficients}
Let $z_{\vartheta}:\N\to\C$ be the unique two-sided inverse of $\mu$ for
$\twconv$:
\begin{equation}
\mu\twconv z_{\vartheta}
=
z_{\vartheta}\twconv\mu
=
\deltaone.
\label{eq:ztheta_inverse_definition}
\end{equation}
Thus
\begin{equation}
z_{\vartheta}(1)=1,
\label{eq:ztheta_one}
\end{equation}
and for $n>1$,
\begin{equation}
z_{\vartheta}(n)
=
-\sum_{\substack{d\mid n\\ d>1}}
\mu(d)z_{\vartheta}(n/d)
\chi_{\vartheta}(d,n/d).
\label{eq:ztheta_recursion}
\end{equation}
\end{definition}

\begin{proposition}[Commuting boundary]
\label{prop:commuting_boundary}
At $\vartheta=0$,
\begin{equation}
z_0(n)=1
\label{eq:z0_constant_one}
\end{equation}
for every $n\ge1$.
\end{proposition}

\begin{proof}
When $\vartheta=0$, $\chi_0\equiv1$ and $\twconv$ is ordinary Dirichlet
convolution.  Equation~\eqref{eq:ztheta_inverse_definition} becomes
$\mu*z_0=\deltaone$.  Classical M\"obius inversion gives $z_0=\mathbf1$.
See \cite{Apostol1976}.
\end{proof}

\begin{proposition}[Ordered-factorization formula]
\label{prop:ordered_factorization_formula}
For \(n>1\), the twisted inverse has the finite coefficient expansion
\begin{align}
z_\vartheta(n)
=
\sum_{r\ge1}(-1)^r
\sum_{\substack{n_1\cdots n_r=n\\ n_j>1}}
&\left(\prod_{j=1}^r\mu(n_j)\right)
\notag\\
&\times
\exp\left[
\frac{i\vartheta}{2}
\sum_{1\le i<j\le r}
\bigl(
\Omega(n_i)\log n_j-
\log n_i\,\Omega(n_j)
\bigr)
\right].
\label{eq:ordered_factorization_formula}
\end{align}
Only finitely many \(r\) and ordered factorizations contribute for fixed
\(n\).
\end{proposition}

\begin{proof}
Write \(\mu=\deltaone+h\) with \(h(1)=0\).  Coefficientwise the formal
Neumann expansion is
\begin{equation}
z_\vartheta
=
\sum_{r\ge0}(-h)^{\star_\vartheta r}.
\label{eq:formal_neumann_coefficientwise}
\end{equation}
For fixed \(n\), a product of more than \(\log_2 n\) factors exceeding
one cannot equal \(n\), so the coefficient sum is finite.  Iterating the
bicharacter cocycle yields the pairwise phase in
\eqref{eq:ordered_factorization_formula}.
\end{proof}

\subsection{A bounded inverse in a rigorous half-plane}
\label{subsec:bounded_inverse}

\begin{definition}[Uniform Neumann threshold]
\label{def:neumann_threshold}
Let \(\sigma_*\) be the unique real number \(>1\) satisfying
\begin{equation}
\frac{\zeta(\sigma_*)}{\zeta(2\sigma_*)}=2.
\label{eq:neumann_threshold_equation}
\end{equation}
Numerically,
\begin{equation}
\sigma_*=1.578017063034\ldots.
\label{eq:neumann_threshold_numeric}
\end{equation}
\end{definition}

\begin{theorem}[Uniform Neumann region]
\label{thm:neumann_region}
For every real \(\vartheta\) and every \(\sigma>\sigma_*\),
\begin{equation}
z_\vartheta\in\mathcal A_\sigma.
\label{eq:ztheta_in_Asigma}
\end{equation}
Consequently
\begin{equation}
\mathcal Z_{\sigma,\vartheta}
=
\Pi_{\sigma,\vartheta}(z_\vartheta)
=
\sum_{n\ge1}
\frac{z_\vartheta(n)}{n^\sigma}W_n
\label{eq:zeta_operator_deformed}
\end{equation}
converges absolutely in operator norm and satisfies
\begin{equation}
\mathcal M_{\sigma,\vartheta}
\mathcal Z_{\sigma,\vartheta}
=
\mathcal Z_{\sigma,\vartheta}
\mathcal M_{\sigma,\vartheta}
=I.
\label{eq:operator_inverse_identity}
\end{equation}
\end{theorem}

\begin{proof}
Write \(\mu=\deltaone+h\), where \(h(1)=0\).  For \(\sigma>1\),
\begin{align}
\|h\|_\sigma
&=
\sum_{n\ge2}\frac{\mu(n)^2}{n^\sigma}
\label{eq:h_norm_sigma_step1}\\
&=
\frac{\zeta(\sigma)}{\zeta(2\sigma)}-1.
\label{eq:h_norm_sigma_step2}
\end{align}
The right-hand side is continuous and strictly decreasing from \(+\infty\)
to \(0\) as \(\sigma\) runs from \(1^+\) to \(+\infty\), so
\(\sigma_*\) is unique and \(\|h\|_\sigma<1\) exactly when
\(\sigma>\sigma_*\).  The Banach-algebra Neumann series
\begin{equation}
(\deltaone+h)^{-1}
=
\sum_{r=0}^{\infty}(-h)^{\star_\vartheta r}
\label{eq:neumann_inverse}
\end{equation}
then converges in \(\mathcal A_\sigma\).  Uniqueness identifies it with
\(z_\vartheta\), and the representation theorem gives
\eqref{eq:operator_inverse_identity}.
\end{proof}

\begin{warning}[Exact border of the proved inverse]
\label{warn:inverse_border}
The M\"obius operator exists as a bounded operator for every \(\sigma>1\).
The argument above proves a bounded inverse only for
\(\sigma>\sigma_*\).  At \(\sigma=\sigma_*\) the Neumann norm equals
one and this argument is inconclusive.  No claim is made here for
\(1<\sigma\le\sigma_*\).
\end{warning}

\section{The first Heisenberg arithmetic law}
\label{sec:prime_pair_interference}

\begin{theorem}[Prime rays are rigid; distinct primes interfere]
\label{thm:prime_pair_interference}
Let $p,q$ be distinct primes and let $r\ge1$.  Then
\begin{align}
z_{\vartheta}(p^r)
&=1,
\label{eq:ztheta_prime_power}\\
z_{\vartheta}(pq)
&=
2\cos\!\left(\frac{\vartheta}{2}\log\frac{q}{p}\right)-1,
\label{eq:ztheta_semiprime_cos}\\
z_{\vartheta}(pq)-1
&=
-4\sin^2\!\left(\frac{\vartheta}{4}\log\frac{q}{p}\right).
\label{eq:ztheta_semiprime_defect}
\end{align}
\end{theorem}

\begin{proof}
For $p^r$, only the divisors $1$ and $p$ contribute to the recursion because
$\mu(p^j)=0$ for $j\ge2$.  Moreover,
\begin{equation}
\Om(p)\log(p^{r-1})
-
\log p\,\Om(p^{r-1})
=0,
\label{eq:prime_power_symplectic_zero}
\end{equation}
so $\chi_{\vartheta}(p,p^{r-1})=1$.  Equation
\eqref{eq:ztheta_recursion} gives
\begin{equation}
z_{\vartheta}(p^r)=z_{\vartheta}(p^{r-1}),
\label{eq:prime_power_recursion}
\end{equation}
and \eqref{eq:ztheta_prime_power} follows from $z_{\vartheta}(1)=1$.

For $pq$, the four divisors $1,p,q,pq$ give
\begin{equation}
0
=
z_{\vartheta}(pq)
-
\chi_{\vartheta}(p,q)
-
\chi_{\vartheta}(q,p)
+1.
\label{eq:pq_inverse_equation}
\end{equation}
Since
\begin{equation}
\chi_{\vartheta}(p,q)
=
\exp\!\left(\frac{i\vartheta}{2}\log\frac{q}{p}\right),
\label{eq:chi_pq}
\end{equation}
and $\chi_{\vartheta}(q,p)$ is its complex conjugate,
\eqref{eq:ztheta_semiprime_cos} follows.  The identity
\eqref{eq:ztheta_semiprime_defect} is the elementary formula
$2(\cos u-1)=-4\sin^2(u/2)$.
\end{proof}

\begin{corollary}[Loss of ordinary multiplicativity]
\label{cor:loss_of_multiplicativity}
Generically,
\begin{equation}
z_{\vartheta}(pq)
\neq
z_{\vartheta}(p)z_{\vartheta}(q)=1.
\label{eq:loss_of_multiplicativity}
\end{equation}
Thus the Heisenberg deformation preserves every one-prime ray but destroys
ordinary multiplicativity through cross-prime interference.
\end{corollary}

\begin{proof}
The prime-pair formula gives
\(z_\vartheta(pq)=2\cos(\frac\vartheta2\log(q/p))-1\).  This equals one
only on the stated resonance set; away from it the product of the two
one-prime values, which is identically one, differs from the two-prime
coefficient.
\end{proof}

\begin{corollary}[Small-area expansion]
\label{cor:small_area_expansion}
For fixed distinct primes $p,q$,
\begin{equation}
z_{\vartheta}(pq)
=
1-
\frac{\vartheta^2}{4}
\left(\log\frac{q}{p}\right)^2
+O(\vartheta^4)
\qquad(\vartheta\to0).
\label{eq:small_theta_pq}
\end{equation}
There is no linear correction.
\end{corollary}

\begin{proof}
Use \(\cos x=1-x^2/2+O(x^4)\) with
\(x=(\vartheta/2)\log(q/p)\) in the exact prime-pair formula.  Multiplying
out gives \eqref{eq:small_theta_pq}; evenness of the cosine excludes an odd
term.
\end{proof}

\begin{proposition}[Two-prime lattice recurrence]
\label{prop:two_prime_recurrence}
Let \(p\ne q\) be primes, set
\begin{equation}
Z_{a,b}=z_\vartheta(p^aq^b),
\qquad
\lambda=\frac{\vartheta}{2}\log\frac{q}{p}.
\label{eq:Zab_lambda}
\end{equation}
Then for \(a,b\ge1\),
\begin{equation}
Z_{a,b}
=
e^{ib\lambda}Z_{a-1,b}
+e^{-ia\lambda}Z_{a,b-1}
-e^{i(b-a)\lambda}Z_{a-1,b-1},
\label{eq:two_prime_recurrence}
\end{equation}
with boundary values
\begin{equation}
Z_{a,0}=Z_{0,b}=1.
\label{eq:two_prime_boundary}
\end{equation}
In particular,
\begin{equation}
z_\vartheta(p^2q)
=1+2\cos(2\lambda)-2\cos\lambda.
\label{eq:p2q_formula}
\end{equation}
\end{proposition}

\begin{proof}
Apply the right-inverse recursion \eqref{eq:ztheta_recursion} to
\(p^aq^b\).  Since \(\mu(d)\neq0\) only when the exponents of \(p,q\)
in \(d\) are at most one, only \(d=p,q,pq\) contribute.  Evaluation of
the three cocycle factors gives \eqref{eq:two_prime_recurrence}; the boundary
values are Theorem~\ref{thm:prime_pair_interference} on prime rays.
\end{proof}

\begin{proposition}[Evenness of the full inverse]
\label{prop:theta_evenness}
For every $n\ge1$,
\begin{equation}
z_{-\vartheta}(n)=z_{\vartheta}(n).
\label{eq:ztheta_evenness}
\end{equation}
\end{proposition}

\begin{proof}
Changing $\vartheta$ to $-\vartheta$ reverses the cocycle:
\begin{equation}
(f\star_{-\vartheta}g)(n)
=
(g\star_{\vartheta}f)(n).
\label{eq:opposite_algebra_identity}
\end{equation}
Since $z_{\vartheta}$ is a two-sided inverse of $\mu$ for
$\star_{\vartheta}$, it is a two-sided inverse of $\mu$ for the opposite
product $\star_{-\vartheta}$.  Uniqueness of the inverse gives
\eqref{eq:ztheta_evenness}.
\end{proof}

\begin{proposition}[Reality for real deformation]
\label{prop:ztheta_real}
For real \(\vartheta\) and every \(n\ge1\),
\begin{equation}
z_\vartheta(n)\in\mathbb R.
\label{eq:ztheta_real}
\end{equation}
\end{proposition}

\begin{proof}
Complex conjugation changes \(\chi_\vartheta\) to
\(\chi_{-\vartheta}\).  Since \(\mu\) is real, conjugating the inverse
equation shows that \(\overline{z_\vartheta}\) is the inverse for the
\(-\vartheta\) product.  Proposition~\ref{prop:theta_evenness} therefore
gives \(\overline{z_\vartheta}=z_{-\vartheta}=z_\vartheta\).
\end{proof}

\begin{remark}[Why this is the non-tautological step]
\label{rem:non_tautological_step}
At $\vartheta=0$ the inverse coefficients are identically $1$.  At nonzero
$\vartheta$, equation~\eqref{eq:ztheta_semiprime_defect} creates a new
arithmetic coefficient depending on the logarithmic separation of two
primes.  This dependence is forced by the symplectic area in the Weyl
product and disappears if the CCR is removed.  It is therefore the first
result in the notebook where Heisenberg noncommutativity changes the
arithmetic inverse itself rather than merely constraining an unrelated
state space.
\end{remark}

\section{Curvature of the cocyclically deformed M\"obius inverse}
\label{sec:cohomological_curvature}

Set
\begin{equation}
\mathcal A(d,e)=\Omega(d)\log e-\log d\,\Omega(e),
\label{eq:arith_symplectic_area_curv}
\end{equation}
so that
\begin{equation}
\chi_\vartheta(d,e)=\exp\!\left(\frac{i\vartheta}{2}\mathcal A(d,e)\right).
\label{eq:chi_area_curv}
\end{equation}

\begin{proposition}[Coefficientwise entire dependence]
\label{prop:ztheta_entire_coefficientwise}
For each fixed \(n\ge1\), the map \(\vartheta\mapsto z_\vartheta(n)\) is
entire. It is even, and it is real-valued on the real axis.
\end{proposition}

\begin{proof}
Formula~\eqref{eq:ordered_factorization_formula} is a finite sum for fixed
\(n\), each term being an exponential \(e^{ic\vartheta}\). Evenness and
reality follow from Propositions~\ref{prop:theta_evenness} and
\ref{prop:ztheta_real}.
\end{proof}

\begin{proposition}[Second derivative of the inverse]
\label{prop:second_derivative_inverse}
Define
\begin{equation}
g(k)=\frac14\sum_{d\mid k}\mu(d)\mathcal A(d,k/d)^2.
\label{eq:gk_definition}
\end{equation}
Then
\begin{equation}
\boxed{z_0''(n)=\sum_{k\mid n}g(k).}
\label{eq:zsecond_divisor_sum}
\end{equation}
\end{proposition}

\begin{proof}
Differentiate the inverse equation
\begin{equation}
0=\sum_{d\mid n}\mu(d)z_\vartheta(n/d)
\exp\!\left(\frac{i\vartheta}{2}\mathcal A(d,n/d)\right)
\label{eq:inverse_equation_theta}
\end{equation}
twice at \(\vartheta=0\). Since \(z_0\equiv1\) and \(z'_0\equiv0\),
\begin{equation}
(\mu*z_0'')(n)=\frac14\sum_{d\mid n}\mu(d)\mathcal A(d,n/d)^2=g(n).
\label{eq:mu_convolution_zsecond}
\end{equation}
Convolving with \(\mathbf1=\mu^{-1}\) gives
\eqref{eq:zsecond_divisor_sum}.
\end{proof}

\begin{lemma}[Support collapse]
\label{lem:gk_support_collapse}
Let \(k=\prod_{j=1}^{r}p_j^{a_j}\) with distinct primes \(p_j\). Then
\begin{equation}
g(k)=0\qquad\text{for }r=0,\ r=1,\text{ or }r\ge3.
\label{eq:gk_support_zero}
\end{equation}
For two distinct primes,
\begin{equation}
\boxed{g(p^\alpha q^\beta)
=-\frac{\alpha\beta}{2}\log^2\!\frac qp.}
\label{eq:g_two_prime_exact}
\end{equation}
\end{lemma}

\begin{proof}
A contributing divisor is square-free and therefore corresponds to a subset
\(S\subseteq\{1,\ldots,r\}\). Put
\begin{equation}
R=\Omega(k)=\sum_j a_j,\qquad L=\log k,
\qquad X_S=\sum_{j\in S}\log p_j.
\label{eq:R_L_XS}
\end{equation}
Then
\begin{equation}
\mathcal A(d,k/d)=|S|L-RX_S.
\label{eq:A_subset_linear}
\end{equation}
Its square has Boolean degree at most two, so its full alternating difference
vanishes when \(r\ge3\). For \(r=1\) the area is zero. For \(r=2\), direct
evaluation of the four subsets gives
\begin{equation}
\sum_{d\mid p^\alpha q^\beta}\mu(d)
\mathcal A(d,p^\alpha q^\beta/d)^2
=-2\alpha\beta\log^2(q/p),
\label{eq:four_term_two_prime_sum}
\end{equation}
which yields \eqref{eq:g_two_prime_exact}.
\end{proof}

\begin{theorem}[Exact deformation curvature]
\label{thm:exact_curvature_formula}
If
\begin{equation}
n=\prod_{j=1}^{r}p_j^{a_j},\qquad T(a)=\frac{a(a+1)}2,
\label{eq:n_factorization_triangular}
\end{equation}
then
\begin{equation}
\boxed{
z_0''(n)=-\frac12\sum_{1\le i<j\le r}
T(a_i)T(a_j)\log^2\!\frac{p_j}{p_i}.}
\label{eq:zsecond_general_formula}
\end{equation}
Thus \(-z_0''(n)\ge0\), with equality exactly for \(n=1\) or a prime
power.
\end{theorem}

\begin{proof}
Only divisors supported on exactly two primes survive by
Lemma~\ref{lem:gk_support_collapse}. Summing
\eqref{eq:g_two_prime_exact} over
\(1\le\alpha\le a_i\), \(1\le\beta\le a_j\) gives the factor
\(T(a_i)T(a_j)\).
\end{proof}

\begin{corollary}[Square-free prime-log variance]
\label{cor:squarefree_prime_variance}
If \(n=p_1\cdots p_r\) is square-free, then
\begin{align}
-z_0''(n)
&=\frac12\sum_{i<j}(\log p_j-\log p_i)^2
\label{eq:squarefree_pair_variance}\\
&=\frac12\left[r\sum_{p\mid n}(\log p)^2-(\log n)^2\right]
\label{eq:squarefree_variance_identity}\\
&=\frac{r^2}{2}\operatorname{Var}_{p\mid n}(\log p).
\label{eq:squarefree_variance_stat}
\end{align}
\end{corollary}

\begin{proof}
In Theorem~\ref{thm:exact_curvature_formula}, square-freeness gives
\(T(v_p(n))=T(1)=1\) for every \(p\mid n\).  The first equality follows.
Expanding the pairwise squares and counting each \((\log p)^2\) exactly
\(r-1\) times gives
\(\sum_{i<j}(a_j-a_i)^2=r\sum_i a_i^2-(\sum_i a_i)^2\), which yields the
second equality.  Dividing by \(r\) identifies the usual empirical variance
and gives the third.
\end{proof}

\begin{corollary}[Infinitesimal symplectic size]
\label{cor:cohomological_singular_value}
For square-free \(n=p_1\cdots p_r\), let
\begin{equation}
B_{ij}=\log p_j-\log p_i.
\label{eq:B_matrix_logs}
\end{equation}
Then \(B\) has rank at most two. If \(r\ge2\), its two nonzero singular
values are both
\begin{equation}
s(n)=\sqrt{r\sum_{p\mid n}(\log p)^2-(\log n)^2},
\label{eq:singular_value_s}
\end{equation}
and
\begin{equation}
\boxed{-z_0''(n)=\frac12s(n)^2.}
\label{eq:curvature_singular_value}
\end{equation}
\end{corollary}

\begin{proof}
Write \(B=\mathbf1a^T-a\mathbf1^T\), \(a_i=\log p_i\). The two nonzero
singular values of \(uv^T-vu^T\) equal
\begin{equation}
\sqrt{\|u\|^2\|v\|^2-(u\cdot v)^2}.
\label{eq:wedge_singular_value}
\end{equation}
Apply this with \(u=\mathbf1\), \(v=a\), and then use
Corollary~\ref{cor:squarefree_prime_variance}.
\end{proof}

\begin{corollary}[Exact Taylor law for each integer]
\label{cor:coefficient_taylor_exact}
For every fixed \(n\),
\begin{equation}
z_\vartheta(n)=1+\frac{\vartheta^2}{2}z_0''(n)+O_n(\vartheta^4),
\label{eq:ztheta_general_taylor}
\end{equation}
and the Taylor series has infinite radius of convergence. For square-free
\(n\),
\begin{equation}
\boxed{
z_\vartheta(n)=1-\frac{\vartheta^2}{4}
\left[\omega(n)\sum_{p\mid n}(\log p)^2-(\log n)^2\right]
+O_n(\vartheta^4).}
\label{eq:ztheta_squarefree_taylor}
\end{equation}
\end{corollary}

\begin{proof}
Coefficientwise entire dependence was proved in
Proposition~\ref{prop:ztheta_entire_coefficientwise}.  Evenness gives the Taylor
expansion with no odd powers.  The coefficient of \(\vartheta^2\) is
\(z_0''(n)/2\), proving the first formula.  Substitution of the square-free
formula from Corollary~\ref{cor:squarefree_prime_variance} yields the boxed
identity.
\end{proof}

\section{Multiplicative geometry of the curvature}
\label{sec:curvature_multiplicative_geometry}

Define
\begin{equation}
\mathcal C(n):=-z_0''(n)
=\frac12\sum_{p<q}T(v_p(n))T(v_q(n))\log^2\frac qp.
\label{eq:Cn_definition_curvature}
\end{equation}
This section studies \(\mathcal C\) before any summation over \(n\).

\begin{proposition}[Prime-ray rigidity]
\label{prop:curvature_prime_ray_rigidity}
\begin{equation}
\boxed{\mathcal C(n)=0\iff n=1\ \text{or}\ n=p^a.}
\label{eq:curvature_zero_prime_ray}
\end{equation}
\end{proposition}

\begin{proof}
Every summand in \eqref{eq:Cn_definition_curvature} is nonnegative.  It is
strictly positive exactly when two distinct primes occur with positive
valuation.
\end{proof}

\begin{theorem}[Multiplicative superadditivity]
\label{thm:curvature_superadditivity}
For all positive integers \(m,n\),
\begin{equation}
\boxed{\mathcal C(mn)\ge \mathcal C(m)+\mathcal C(n).}
\label{eq:curvature_superadditivity}
\end{equation}
If \(m,n>1\), equality holds exactly when \(m=p^a\) and \(n=p^b\) for the
same prime \(p\).
\end{theorem}

\begin{proof}
Put \(a_p=v_p(m)\), \(b_p=v_p(n)\),
\(A_p=T(a_p)\), \(B_p=T(b_p)\).  Since
\(T(a+b)=T(a)+T(b)+ab\), subtraction of
\(\mathcal C(m)+\mathcal C(n)\) from \(\mathcal C(mn)\) gives
\begin{equation}
\mathcal C(mn)-\mathcal C(m)-\mathcal C(n)
=\frac12\sum_{p<q}\Xi_{p,q}(m,n)\log^2\frac qp,
\label{eq:curvature_superadditivity_expansion}
\end{equation}
where
\begin{align}
\Xi_{p,q}(m,n)={}&A_pB_q+B_pA_q
 +(A_p+B_p)a_qb_q +(A_q+B_q)a_pb_p
 +a_pb_p a_qb_q.
\label{eq:Xi_pq_definition}
\end{align}
All terms are nonnegative.  If \(m,n>1\), all cross terms vanish exactly
when both supports consist of the same single prime.
\end{proof}

\begin{corollary}[Monotonicity on the divisibility poset]
\label{cor:curvature_divisibility_monotonicity}
If \(m\mid n\), then
\begin{equation}
\boxed{\mathcal C(m)\le\mathcal C(n).}
\label{eq:curvature_divisibility_monotonicity}
\end{equation}
\end{corollary}

\begin{proof}
Write \(n=mq\).  Theorem~\ref{thm:curvature_superadditivity} gives
\(\mathcal C(n)\ge\mathcal C(m)+\mathcal C(q)\ge\mathcal C(m)\), because
\(\mathcal C(q)\ge0\).  This is exactly
\eqref{eq:curvature_divisibility_monotonicity}.
\end{proof}

\begin{proposition}[Exact power law]
\label{prop:curvature_power_law}
If \(n=\prod_p p^{a_p}\), then for every integer \(k\ge1\),
\begin{equation}
\boxed{
\mathcal C(n^k)=\frac18\sum_{p<q}
\left[k^4a_p^2a_q^2+k^3(a_p^2a_q+a_pa_q^2)+k^2a_pa_q\right]
\log^2\frac qp.}
\label{eq:curvature_power_exact}
\end{equation}
Hence, if \(n\) contains at least two distinct primes,
\begin{equation}
\boxed{
\lim_{k\to\infty}\frac{\mathcal C(n^k)}{k^4}
=\frac18\sum_{p<q}a_p^2a_q^2\log^2\frac qp>0.}
\label{eq:curvature_power_asymptotic}
\end{equation}
\end{proposition}

\begin{proof}
For every prime exponent \(a_p\),
\begin{equation}
T(ka_p)=\frac{k^2a_p^2+ka_p}{2}.
\label{eq:T_ka_expansion}
\end{equation}
Thus for each pair \(p<q\),
\begin{align}
\frac12T(ka_p)T(ka_q)
&=\frac18
\bigl[k^4a_p^2a_q^2+k^3(a_p^2a_q+a_pa_q^2)+k^2a_pa_q\bigr].
\label{eq:T_pair_power_expansion}
\end{align}
Substitution into \eqref{eq:Cn_definition_curvature} gives
\eqref{eq:curvature_power_exact}.  Division by \(k^4\) and termwise passage
to the limit is legitimate because the sum contains only finitely many prime
pairs.  If the support has at least two primes, at least one coefficient
\(a_p^2a_q^2\log^2(q/p)\) is strictly positive, proving
\eqref{eq:curvature_power_asymptotic}.
\end{proof}

\begin{proposition}[Coprime block decomposition]
\label{prop:curvature_coprime_decomposition}
Assume \((m,n)=1\).  Define
\begin{align}
W(m)&=\sum_{p\mid m}T(v_p(m)),\\
\overline L(m)&=\frac1{W(m)}\sum_{p\mid m}T(v_p(m))\log p,\\
V(m)&=\frac1{W(m)}\sum_{p\mid m}T(v_p(m))
(\log p-\overline L(m))^2,
\end{align}
with the analogous definitions for \(n\).  Then
\begin{equation}
\boxed{
\mathcal C(mn)=\mathcal C(m)+\mathcal C(n)
+\frac12W(m)W(n)\left[V(m)+V(n)+(\overline L(m)-\overline L(n))^2\right].}
\label{eq:curvature_coprime_ANOVA}
\end{equation}
\end{proposition}

\begin{proof}
The cross-prime contribution is
\(\frac12\sum_{p\mid m,q\mid n}A_pB_q(\log q-\log p)^2\).
Expanding around the two weighted means gives
\eqref{eq:curvature_coprime_ANOVA}.  Abstractly this last algebraic identity
is the standard weighted within/between variance decomposition; its use here
is only to expose the multiplicative geometry of \(\mathcal C\).
\end{proof}

\section{Dirichlet summation of the curvature}
\label{sec:curvature_dirichlet_series}

Define
\begin{equation}
c(n)=-z_0''(n)\ge0.
\label{eq:cn_curvature}
\end{equation}

\begin{theorem}[Exact curvature Dirichlet series]
\label{thm:curvature_dirichlet_series}
For \(\Re s>1\),
\begin{equation}
\mathcal K(s)=\sum_{n\ge1}\frac{c(n)}{n^s}
\label{eq:K_dirichlet_definition}
\end{equation}
converges absolutely and
\begin{equation}
\boxed{
\mathcal K(s)=\frac{\zeta(s)}2\left[S_0(s)S_2(s)-S_1(s)^2\right],}
\label{eq:K_variance_determinant}
\end{equation}
where
\begin{equation}
S_j(s)=\sum_p\frac{p^{-s}}{(1-p^{-s})^2}(\log p)^j,
\qquad j=0,1,2.
\label{eq:Sj_prime_sum}
\end{equation}
Equivalently,
\begin{equation}
S_0S_2-S_1^2
=\sum_{p<q}
\frac{p^{-s}q^{-s}}{(1-p^{-s})^2(1-q^{-s})^2}
\log^2\!\frac qp.
\label{eq:weighted_prime_variance}
\end{equation}
\end{theorem}

\begin{proof}
Use Theorem~\ref{thm:exact_curvature_formula} and
\begin{equation}
\sum_{a\ge1}T(a)x^a=\frac{x}{(1-x)^3}.
\label{eq:triangular_generating_function}
\end{equation}
Factoring the unrestricted Euler factors gives \(\zeta(s)\), while a selected
prime contributes
\begin{equation}
\frac{p^{-s}}{(1-p^{-s})^2}.
\label{eq:selected_prime_weight}
\end{equation}
The pair sum is exactly the weighted variance determinant
\(S_0S_2-S_1^2\). Absolute convergence follows for \(\Re s>1\).
\end{proof}

\begin{corollary}[Prime-zeta reduction]
\label{cor:prime_zeta_representation}
Let \(P(s)=\sum_p p^{-s}\) for \(\Re s>1\). Then
\begin{align}
S_0(s)&=\sum_{m\ge1}mP(ms),
\label{eq:S0_prime_zeta}\\
S_1(s)&=-\sum_{m\ge1}mP'(ms),
\label{eq:S1_prime_zeta}\\
S_2(s)&=\sum_{m\ge1}mP''(ms).
\label{eq:S2_prime_zeta}
\end{align}
Using
\begin{equation}
P(s)=\sum_{k\ge1}\frac{\mu(k)}{k}\log\zeta(ks),
\label{eq:prime_zeta_mobius_inversion}
\end{equation}
gives
\begin{align}
S_0(s)&=\sum_{r\ge1}\frac{J_2(r)}{r}\log\zeta(rs),
\label{eq:S0_logzeta}\\
S_1(s)&=-\sum_{r\ge1}\varphi(r)\frac{\zeta'(rs)}{\zeta(rs)},
\label{eq:S1_logderivative}\\
S_2(s)&=\left(\frac{\zeta'}{\zeta}\right)'(s),
\label{eq:S2_logderivative_exact}
\end{align}
where \(J_2(r)=r^2\prod_{p\mid r}(1-p^{-2})\).
\end{corollary}

\begin{proof}
Expand \(x/(1-x)^2=\sum_{m\ge1}mx^m\), substitute
\eqref{eq:prime_zeta_mobius_inversion}, and group by \(r=mk\). The needed
divisor identities are
\begin{align}
r\sum_{k\mid r}\frac{\mu(k)}{k^2}&=\frac{J_2(r)}r,
\label{eq:J2_divisor_identity}\\
r\sum_{k\mid r}\frac{\mu(k)}k&=\varphi(r),
\label{eq:phi_divisor_identity}\\
r\sum_{k\mid r}\mu(k)&=\begin{cases}1,&r=1,\\0,&r>1.\end{cases}
\label{eq:mobius_divisor_delta}
\end{align}
The last identity collapses \(S_2\) to
\eqref{eq:S2_logderivative_exact}.
\end{proof}

\begin{proposition}[Singularity at \(s=1\)]
\label{prop:K_near_one}
As \(s\to1^+\),
\begin{equation}
\boxed{
\mathcal K(s)\sim\frac1{2(s-1)^3}\log\frac1{s-1}.}
\label{eq:K_near_one_asymptotic}
\end{equation}
\end{proposition}

\begin{proof}
From \eqref{eq:prime_zeta_mobius_inversion},
\begin{equation}
P(s)=\log\frac1{s-1}+O(1).
\label{eq:P_near_one}
\end{equation}
Hence
\begin{equation}
S_0(s)=\log\frac1{s-1}+O(1),\quad
S_1(s)=\frac1{s-1}+O(1),\quad
S_2(s)=\frac1{(s-1)^2}+O(1).
\label{eq:S012_near_one}
\end{equation}
Combine with \(\zeta(s)\sim(s-1)^{-1}\).
\end{proof}

\section{Continuation of the curvature and detection of zeta zeros}
\label{sec:curvature_zero_detection}

\begin{theorem}[Local continuation to \(\Re s>1/2\)]
\label{thm:K_local_continuation}
Let \(U\subset\{\Re s>1/2\}\) be simply connected and contain no zero or pole
of \(\zeta(s)\). Choose a branch of \(\log\zeta(s)\) on \(U\). Then
\eqref{eq:S0_logzeta}--\eqref{eq:S2_logderivative_exact}, substituted into
\eqref{eq:K_variance_determinant}, continue \(\mathcal K\) holomorphically to
\(U\).
\end{theorem}

\begin{proof}
For \(r\ge2\), \(\Re(rs)>1\), so the corresponding terms are holomorphic and
converge normally on compact subsets. The \(r=1\) terms are holomorphic by the
hypothesis on \(U\).
\end{proof}

\begin{theorem}[Zeta zeros are logarithmic curvature singularities]
\label{thm:zero_curvature_singularity}
Let \(\rho\) be a zero of multiplicity \(m\ge1\) with
\(1/2<\Re\rho<1\), and write
\begin{equation}
\zeta(s)=c_\rho(s-\rho)^m(1+O(s-\rho)),\qquad c_\rho\ne0.
\label{eq:zeta_local_zero_m}
\end{equation}
Then
\begin{equation}
\boxed{
\mathcal K(s)=-\frac{c_\rho m^2}{2}
(s-\rho)^{m-2}\log(s-\rho)
+O\!\left((s-\rho)^{m-2}\right).}
\label{eq:K_zero_local_form}
\end{equation}
For a simple zero,
\begin{equation}
\boxed{
\mathcal K(s)=-\frac{\zeta'(\rho)}2
\frac{\log(s-\rho)}{s-\rho}
+O\!\left(\frac1{s-\rho}\right).}
\label{eq:K_simple_zero_form}
\end{equation}
\end{theorem}

\begin{proof}
For \(r\ge2\), the terms in
\eqref{eq:S0_logzeta}--\eqref{eq:S1_logderivative} are holomorphic near
\(\rho\). With \(t=s-\rho\), the singular pieces are
\begin{equation}
S_0=m\log t+O(1),\qquad
S_1=-\frac mt+O(1),\qquad
S_2=-\frac m{t^2}+O(1).
\label{eq:S012_near_zero}
\end{equation}
Thus
\begin{equation}
S_0S_2-S_1^2=-\frac{m^2\log t}{t^2}+O(t^{-2}).
\label{eq:det_near_zero}
\end{equation}
Multiply by \(\zeta(s)/2\).
\end{proof}

\begin{corollary}[Curvature criterion equivalent to RH]
\label{cor:curvature_RH_criterion}
RH is equivalent to the absence of logarithmic branch singularities of
\(\mathcal K\) in
\begin{equation}
\boxed{\frac12<\Re s<1.}
\label{eq:RH_curvature_halfstrip}
\end{equation}
\end{corollary}

\begin{proof}
Theorem~\ref{thm:zero_curvature_singularity} detects every zero in the right
half of the critical strip, while Theorem~\ref{thm:K_local_continuation}
gives holomorphy on zero-free neighbourhoods. The functional equation makes
the absence of zeros in this half-strip equivalent to RH.
\end{proof}


\section{Where the logarithmic branching comes from}
\label{sec:prime_zeta_branching}

The logarithmic singularity in Theorem~\ref{thm:zero_curvature_singularity}
does not come from a meromorphic function such as \(\zeta\), \(\zeta'\), or
\(\zeta''\).  It comes from the prime-zeta component of \(S_0\).

\begin{proposition}[Exact source of the logarithm]
\label{prop:source_of_logarithm}
For \(\Re s>1\),
\begin{equation}
S_0(s)=\sum_{m\ge1}mP(ms),
\qquad
P(s)=\sum_{r\ge1}\frac{\mu(r)}{r}\log\zeta(rs).
\label{eq:S0_prime_zeta_source}
\end{equation}
Consequently
\begin{equation}
S_0(s)=\sum_{k\ge1}\frac{J_2(k)}{k}\log\zeta(ks).
\label{eq:S0_branch_expansion}
\end{equation}
If \(\rho\) is a zero of \(\zeta\) of multiplicity \(m\) and
\(1/2<\Re\rho<1\), then in a punctured simply connected neighbourhood of
\(\rho\),
\begin{equation}
S_0(s)=m\log(s-\rho)+O(1).
\label{eq:S0_local_log_source}
\end{equation}
\end{proposition}

\begin{proof}
The first identity follows by expanding
\(p^{-s}(1-p^{-s})^{-2}=\sum_{m\ge1}m p^{-ms}\) and summing over primes.
The second follows from M\"obius inversion for the prime zeta function and
collecting the coefficient of \(\log\zeta(ks)\), exactly as in
Corollary~\ref{cor:prime_zeta_representation}.  If
\(1/2<\Re\rho<1\), every term with \(k\ge2\) is holomorphic near \(\rho\),
because \(\Re(k\rho)>1\).  The \(k=1\) term is
\(\log\zeta(s)=m\log(s-\rho)+O(1)\).
\end{proof}

\begin{remark}[Global branch structure]
\label{rem:global_branch_structure}
For the classical prime-zeta function, only square-free dilation indices occur directly because of the factor \(\mu(r)\).  By contrast, \(S_0\) carries the nonzero weight \(J_2(r)/r\) for every \(r\).  The classical prime-zeta continuation generated by
\begin{equation}
P(s)=\sum_{r\ge1}\frac{\mu(r)}r\log\zeta(rs)
\label{eq:prime_zeta_global_branches}
\end{equation}
has logarithmic branch points at scaled poles and zeros, including points
\(1/r\) and \(\rho/r\), and the line \(\Re s=0\) is a natural boundary for
\(P\); see Fr\"oberg~\cite{Froberg1968}.  This global phenomenon does not
invalidate Corollary~\ref{cor:curvature_RH_criterion}: in the open half-strip
\(1/2<\Re s<1\), scaled points \(\rho/r\) with \(r\ge2\) cannot occur,
because \(\Re(\rho/r)\le1/2\).  Hence the only logarithmic branches there
are the unscaled \(r=1\) branches attached to zeros of \(\zeta\) itself.
\end{remark}

\section{Local monodromy and the Jordan-totient weight}
\label{sec:local_monodromy}

The branch structure is used here only locally.  No globally reordered double
sum over zeta zeros is asserted.

\begin{proposition}[Local monodromy at an isolated scaled zero]
\label{prop:local_monodromy}
Let \(\rho\) be a zero of \(\zeta\) of multiplicity \(m_\rho\), let
\(r\ge1\), and set \(s_0=\rho/r\).  Suppose a sufficiently small punctured
neighbourhood of \(s_0\) contains no other branch point contributing to
\eqref{eq:S0_branch_expansion}.  Analytic continuation once positively around
\(s_0\) changes \(S_0\) by
\begin{equation}
\boxed{
\Delta_{\rho/r} S_0
=2\pi i\,m_\rho\frac{J_2(r)}{r}.}
\label{eq:S0_local_monodromy}
\end{equation}
Since \(S_1\) and \(S_2\) are meromorphic and hence single-valued, the
corresponding monodromy of \(\mathcal K\) is
\begin{equation}
\boxed{
\Delta_{\rho/r}\mathcal K(s)
=\pi i\,m_\rho\frac{J_2(r)}{r}\,
\zeta(s)\left(\frac{\zeta'}{\zeta}\right)'(s).}
\label{eq:K_local_monodromy}
\end{equation}
If several scaled singularities collide at the same point, their local
monodromy contributions are added before applying this formula.
\end{proposition}

\begin{proof}
In \eqref{eq:S0_branch_expansion}, the branch of
\(\log\zeta(rs)\) changes by \(2\pi i m_\rho\) around
\(s_0=\rho/r\).  Multiplication by its coefficient \(J_2(r)/r\) gives
\eqref{eq:S0_local_monodromy}.  The functions \(S_1\) and \(S_2\) are
built from logarithmic derivatives of \(\zeta\), so they are meromorphic and
carry no logarithmic monodromy.  Substitution in
\(\mathcal K=\zeta(S_0S_2-S_1^2)/2\) gives
\eqref{eq:K_local_monodromy}.
\end{proof}

\begin{proposition}[Growth of the local branch weight]
\label{prop:J2_weight_growth}
For every \(r\ge1\),
\begin{equation}
\frac{r}{\zeta(2)}\le \frac{J_2(r)}{r}
=r\prod_{p\mid r}\left(1-\frac1{p^2}\right)\le r.
\label{eq:J2_weight_bounds}
\end{equation}
Consequently the individual local weight is of linear size in \(r\).  In
particular, the unregularized cumulative weight cannot be finite; more
precisely,
\begin{equation}
\sum_{r\le R}\frac{J_2(r)}r
=\frac{R^2}{2\zeta(3)}+O(R\log R).
\label{eq:J2_weight_average}
\end{equation}
\end{proposition}

\begin{proof}
The product in \eqref{eq:J2_weight_bounds} is bounded below by
\(\prod_p(1-p^{-2})=1/\zeta(2)\).  For the average, use
\(J_2(r)/r=r\sum_{d\mid r}\mu(d)d^{-2}\), interchange the finite sums, and
sum \(m\) up to \(R/d\).  The main term is
\(R^2\sum_{d\ge1}\mu(d)d^{-3}/2=R^2/(2\zeta(3))\); elementary estimation of
the remaining terms gives the displayed error.
\end{proof}

\begin{remark}[Why \(S_0\) nevertheless converges away from its branches]
\label{rem:J2_decay_compensation}
For fixed \(s\) with \(\Re s>0\), once \(r\Re s>1\),
\(\log\zeta(rs)=O(2^{-r\Re s})\).  This exponential decay dominates the
linear factor \(J_2(r)/r\).  Thus the growth of the local monodromy weight
is compatible with convergence of \eqref{eq:S0_branch_expansion} away from
its singular points.
\end{remark}

\begin{warning}[No global zero-kernel identity is claimed]
\label{warn:no_global_zero_kernel}
A formal substitution of a Hadamard expansion for \(\log\zeta\) into
\eqref{eq:S0_branch_expansion} produces nested sums over \(r\) and zeta
zeros.  Those sums require prescribed symmetric ordering in the zero
variable and separate convergence control in \(r\).  Version 8 does not
interchange them and does not claim a globally convergent identity of the
form \(\sum_{\rho,\rho'}\mathcal K_s(\rho,\rho')\).  The local monodromy
Proposition~\ref{prop:local_monodromy} does not require such a rearrangement.
\end{warning}


\section{Gauge invariance of the arithmetic cocycle}
\label{sec:gauge_invariance}

The previous sections use the completely additive coordinate
\(A(n)=\Omega(n)\).  It is useful to record an exact invariance that holds for
any completely additive \(A\).

\begin{proposition}[Exact gauge invariance]
\label{prop:gauge_invariance}
Let
\begin{equation}
\mathcal A_A(m,n)=A(m)\log n-\log m\,A(n)
\label{eq:general_additive_area}
\end{equation}
and let \(\chi_{\vartheta,A}=\exp(i\vartheta\mathcal A_A/2)\).
For every \(c\in\mathbb R\),
\begin{equation}
\boxed{
\mathcal A_{A+c\log}(m,n)=\mathcal A_A(m,n),
\qquad
\chi_{\vartheta,A+c\log}=\chi_{\vartheta,A}.}
\label{eq:gauge_invariance}
\end{equation}
Consequently every twisted product and twisted inverse constructed only from
this cocycle depends on \(A\) only through its class modulo
\(\mathbb R\log\).
\end{proposition}

\begin{proof}
The additional contribution is
\begin{equation}
c\log m\log n-c\log m\log n=0.
\label{eq:gauge_added_term_zero}
\end{equation}
\end{proof}

\begin{remark}
\label{rem:gauge_obstruction}
This is an obstruction as well as an invariance: the deformation cannot
recover the component of an additive coordinate parallel to \(\log\) from
the cocycle alone.  The statement concerns the chosen cocycle itself and
should not be confused with invariance of the coefficient curvature under
arbitrary replacement by a cohomologous cocycle.
\end{remark}

\section{Cohomological obstruction to scalar Euler locality}
\label{sec:euler_locality_obstruction}

Write \(e_n=\delta_n\) for the canonical arithmetic basis.  For distinct
primes \(p<q\),
\begin{equation}
e_p\star_\vartheta e_q
=
e^{\,i\vartheta\log(q/p)/2}e_{pq},
\qquad
e_q\star_\vartheta e_p
=
e^{-\,i\vartheta\log(q/p)/2}e_{pq}.
\label{eq:prime_basis_twist}
\end{equation}

\begin{theorem}[No scalar prime-local Euler factorization]
\label{thm:no_scalar_euler_product}
Fix \(\vartheta\in\mathbb R\).  Suppose that there exist scalars
\(a_p(\vartheta)\) and a fixed increasing order of the primes such that the
square-free coefficients of \(\mu\) are produced by an ordered twisted
product
\begin{equation}
\mu=
\mathop{\star_\vartheta}_{p}^{\longrightarrow}
(1-a_p e_p).
\label{eq:putative_scalar_euler_mu}
\end{equation}
Then \(\vartheta=0\).  Hence, for every \(\vartheta\ne0\), the nontrivial
arithmetic cocycle obstructs any scalar Euler factorization into independent
one-prime factors of this form.
\end{theorem}

\begin{proof}
For primes \(p<q\), comparison of the \(pq\)-coefficient gives
\begin{equation}
a_pa_q
e^{\,i\vartheta\log(q/p)/2}=1.
\label{eq:ap_aq_consistency}
\end{equation}
For \(p<q<r\), the three corresponding equations imply
\begin{equation}
a_p^2=e^{-i\vartheta\log(q/p)}.
\label{eq:ap_square_resonance}
\end{equation}
The left side is independent of \(q\).  Taking two primes \(q_1,q_2>p\)
therefore yields
\begin{equation}
e^{-i\vartheta\log(q_1/q_2)}=1.
\label{eq:universal_prime_resonance}
\end{equation}
If \(\vartheta\ne0\), applying this simultaneously to independent prime
ratios forces a nontrivial rational linear relation among logarithms of
distinct primes, contradicting unique factorization.
\end{proof}

\begin{remark}[What is and is not claimed]
\label{rem:euler_obstruction_scope}
Nontrivial 2-cocycles and twisted algebras are classical.  The theorem is a
specific obstruction inside the present arithmetic deformation.  It does not
say that every nontrivial cocycle destroys every conceivable noncommutative
or adelic Euler product; it rules out scalar independent one-prime factors
compatible with the stated twisted product.
\end{remark}

\section{Ordered Euler object and the global Euler-locality defect}
\label{sec:ordered_euler_defect}

For \(\sigma>1\), let \(\mathcal A_\sigma\) be the Banach space of arithmetic
functions with
\begin{equation}
\|f\|_\sigma=\sum_{n\ge1}|f(n)|n^{-\sigma}.
\label{eq:Asigma_norm}
\end{equation}
Since \(|\chi_\vartheta|=1\),
\begin{equation}
\|f\star_\vartheta g\|_\sigma
\le \|f\|_\sigma\|g\|_\sigma.
\label{eq:twisted_Banach_submultiplicative}
\end{equation}

For a prime \(p\), define
\begin{equation}
L_p=\sum_{a\ge0}e_{p^a}.
\label{eq:local_ray_factor}
\end{equation}
The cocycle is trivial on a single prime ray, so
\begin{equation}
\boxed{L_p=(1-e_p)_{\star_\vartheta}^{-1}.}
\label{eq:local_ray_inverse}
\end{equation}
Moreover,
\begin{equation}
\sum_p\|L_p-1\|_\sigma
=
\sum_p\frac{p^{-\sigma}}{1-p^{-\sigma}}<\infty.
\label{eq:ordered_euler_norm_sum}
\end{equation}

\begin{proposition}[Norm-convergent ordered Euler object]
\label{prop:ordered_euler_object}
For \(\sigma>1\), the increasing-prime product
\begin{equation}
\boxed{
E_\vartheta
=
\mathop{\star_\vartheta}_{p}^{\longrightarrow}L_p
}
\label{eq:Etheta_ordered_product}
\end{equation}
converges in \(\mathcal A_\sigma\).  If
\(n=\prod_{j=1}^r p_j^{a_j}\) with \(p_1<\cdots<p_r\), then
\begin{equation}
\boxed{
E_\vartheta(n)
=
\exp\!\left[
\frac{i\vartheta}{2}
\sum_{i<j}a_i a_j\log\frac{p_j}{p_i}
\right].}
\label{eq:Etheta_coefficient}
\end{equation}
\end{proposition}

\begin{proof}
The norm convergence follows from
\eqref{eq:ordered_euler_norm_sum} and the standard convergence criterion for
ordered infinite products in a Banach algebra.  The coefficient formula
follows by repeated use of the cocycle identity; all self-ray phases vanish.
\end{proof}

For \(\sigma>\sigma_*\), the twisted M\"obius inverse \(z_\vartheta\) is
already an element of \(\mathcal A_\sigma\).  Define
\begin{equation}
\boxed{R_\vartheta=z_\vartheta-E_\vartheta.}
\label{eq:Rtheta_def}
\end{equation}

\begin{theorem}[Exact vanishing criterion for the Euler-locality defect]
\label{thm:Rtheta_vanishing}
For real \(\vartheta\),
\begin{equation}
\boxed{R_\vartheta\equiv0\iff\vartheta=0.}
\label{eq:Rtheta_vanishing}
\end{equation}
In fact, products of two primes already detect nonvanishing.
\end{theorem}

\begin{proof}
At \(\vartheta=0\), both coefficient systems equal \(1\).
For \(p<q\), put
\(\lambda=\frac{\vartheta}{2}\log(q/p)\).  Then
\begin{equation}
R_\vartheta(pq)
=
2\cos\lambda-1-e^{i\lambda}
=
e^{-i\lambda}-1.
\label{eq:Rtheta_pq}
\end{equation}
If this vanished for every prime pair, then
\(\vartheta\log(q/p)\in4\pi\mathbb Z\) for every \(p<q\), contradicting the
same universal-resonance argument used in
Theorem~\ref{thm:no_scalar_euler_product} unless \(\vartheta=0\).
\end{proof}


\section{The odd Euler-locality defect: existence and exact initial domain}
\label{sec:odd_euler_defect}

Because \(z_\vartheta=z_{-\vartheta}\), one has \(z'_0(n)=0\).  Define
\begin{equation}
\Phi_{\rm ord}(n)
=
\sum_{p<q}v_p(n)v_q(n)\log\frac qp.
\label{eq:Phi_ord}
\end{equation}
Differentiating the coefficient formula for the ordered Euler object gives
\begin{equation}
R'_0(n)=-\frac{i}{2}\Phi_{\rm ord}(n).
\label{eq:Rprime0}
\end{equation}

\begin{definition}[Odd Euler-locality Dirichlet defect]
\label{def:O_s}
For \(\Re s>1\), set
\begin{equation}
\mathcal O(s)
=
\frac12\sum_{n\ge1}\frac{\Phi_{\rm ord}(n)}{n^s}.
\label{eq:O_dirichlet}
\end{equation}
Whenever the differentiated defect series is used, it is understood through
this absolutely convergent Dirichlet series; no interchange with
\(\partial_\vartheta\) is made outside the domain justified below.
\end{definition}

Put
\begin{equation}
u_p(s)=\frac{p^{-s}}{1-p^{-s}}=\frac1{p^s-1},
\qquad
A(s)=\sum_pu_p(s),
\qquad
B(s,t)=\sum_pu_p(s)p^{it}.
\label{eq:A_B_def}
\end{equation}

\begin{theorem}[Absolute convergence, holomorphy, and differentiation]
\label{thm:O_initial_domain}
The series \eqref{eq:O_dirichlet}, \(A(s)\), and \(B(s,t)\) converge
absolutely and locally uniformly for \(\Re s>1\), uniformly in real \(t\) on
compact sets.  The function \(\mathcal O\) is holomorphic there.  For every
nonnegative integers \(a,b\), the derivatives
\(\partial_s^a\partial_t^bB(s,t)\) are obtained by termwise
differentiation on compact subsets of \(\Re s>1\) and bounded real
\(t\)-intervals.
\end{theorem}

\begin{proof}
Fix \(\varepsilon>0\) and restrict to \(\Re s\ge1+\varepsilon\).  Since
\(|1-p^{-s}|\ge1-p^{-1-\varepsilon}\),
\begin{equation}
|u_p(s)|
\le
\frac{p^{-1-\varepsilon}}{1-p^{-1-\varepsilon}}
\le C_\varepsilon p^{-1-\varepsilon},
\label{eq:up_uniform_bound}
\end{equation}
with a constant independent of \(p\).  Therefore
\(\sum_p|u_p(s)|(\log p)^j\) converges uniformly on the stated compacta for
every fixed \(j\ge0\), by comparison with
\(\sum_{n\ge2}n^{-1-\varepsilon}(\log n)^j\).
This proves normal convergence of \(A\), \(B\), and all the displayed
termwise derivatives.

For \(\mathcal O\), use
\begin{equation}
0\le \Phi_{\rm ord}(n)
\le
\sum_{p<q}v_p(n)v_q(n)(\log p+\log q).
\label{eq:Phi_ord_majorant}
\end{equation}
Summing first over \(n\) with fixed distinct \(p,q\), absolute convergence of
Euler products gives
\begin{equation}
\sum_{n\ge1}\frac{v_p(n)v_q(n)}{n^s}
=
\zeta(s)u_p(s)u_q(s),
\qquad \Re s>1.
\label{eq:vpvq_euler_identity}
\end{equation}
Consequently the absolute majorant for \eqref{eq:O_dirichlet} is bounded by a
constant multiple of
\begin{equation}
\left(\sum_p|u_p(s)|\right)
\left(\sum_q|u_q(s)|\log q\right),
\label{eq:O_absolute_majorant}
\end{equation}
which is locally uniformly finite by \eqref{eq:up_uniform_bound}.  The
Weierstrass theorem then gives holomorphy and termwise differentiation in
\(s\).
\end{proof}

\begin{theorem}[Prime-pair and Fourier representations]
\label{thm:O_prime_pair_fourier}
For \(\Re s>1\),
\begin{align}
\mathcal O(s)
&=
\frac{\zeta(s)}2
\sum_{p<q}u_p(s)u_q(s)\log\frac qp,
\label{eq:O_prime_pair}\\
&=
\frac{\zeta(s)}{2\pi}
\int_0^\infty
\frac{A(s)^2-B(s,t)B(s,-t)}{t^2}\,dt.
\label{eq:O_t_integral}
\end{align}
The integral is absolutely convergent.  Moreover,
\begin{equation}
B(s,t)=\sum_{k\ge1}P(ks-it),
\qquad \Re s>1,
\label{eq:B_prime_zeta}
\end{equation}
with absolute and locally uniform convergence for real \(t\).
\end{theorem}

\begin{proof}
Identity \eqref{eq:vpvq_euler_identity} follows by multiplying the local
geometric series at \(p\) and \(q\):
\begin{equation}
\sum_{a\ge0}a p^{-as}=\frac{p^{-s}}{(1-p^{-s})^2},
\label{eq:vp_local_sum}
\end{equation}
and division by the two ordinary local Euler factors leaves one factor
\(u_p(s)u_q(s)\).  Inserting this identity into the absolutely convergent
series of Theorem~\ref{thm:O_initial_domain} proves
\eqref{eq:O_prime_pair}.

Next, normal convergence permits multiplication of the two prime sums:
\begin{equation}
A(s)^2-B(s,t)B(s,-t)
=
2\sum_{p<q}u_p(s)u_q(s)
\bigl[1-\cos(t\log(q/p))\bigr].
\label{eq:A2_minus_BB_pair_expansion}
\end{equation}
For real \(a\), integration by parts gives
\begin{equation}
\int_0^\infty\frac{1-\cos(at)}{t^2}\,dt
=
\frac{\pi}{2}|a|.
\label{eq:cosine_energy_identity}
\end{equation}
The boundary term vanishes because
\((1-\cos(at))/t=O(t)\) at zero and is \(O(t^{-1})\) at infinity; the
remaining integral is the classical Dirichlet integral.
Furthermore,
\begin{align}
&\sum_{p<q}|u_p(s)u_q(s)|
\int_0^\infty
\frac{|1-\cos(t\log(q/p))|}{t^2}\,dt
\notag\\
&\qquad\le
\frac{\pi}{2}
\sum_{p<q}|u_p(s)u_q(s)|(\log p+\log q)<\infty,
\label{eq:fubini_energy_majorant}
\end{align}
by \eqref{eq:O_absolute_majorant}.  Fubini's theorem is therefore applicable,
and \eqref{eq:O_t_integral} follows.

Finally, \(u_p(s)=\sum_{k\ge1}p^{-ks}\) absolutely for \(\Re s>1\).  Thus
\begin{equation}
\sum_p\sum_{k\ge1}|p^{-ks}|
\le
\sum_p\frac{p^{-\sigma}}{1-p^{-\sigma}}<\infty,
\label{eq:B_double_absolute}
\end{equation}
so Tonelli--Fubini gives \eqref{eq:B_prime_zeta}.  The same majorant is
independent of real \(t\), proving the asserted local uniformity.
\end{proof}

\begin{definition}[Normalized Fourier energy]
\label{def:I_s}
In \(\Re s>1\), define
\begin{equation}
\boxed{
\mathcal I(s)
:=
\int_0^\infty
\frac{A(s)^2-B(s,t)B(s,-t)}{t^2}\,dt
=
\frac{2\pi\mathcal O(s)}{\zeta(s)}.}
\label{eq:I_definition}
\end{equation}
\end{definition}

\subsection{Endpoint regularity and transverse derivatives}
\label{subsec:I_endpoint}

For
\begin{equation}
M_j(s)=\sum_pu_p(s)(\log p)^j,
\label{eq:Mj_u_weights}
\end{equation}
normal convergence from Theorem~\ref{thm:O_initial_domain} yields
\begin{equation}
B(s,t)=A(s)+iM_1(s)t-\frac12M_2(s)t^2+O_K(t^3)
\label{eq:B_taylor_t}
\end{equation}
uniformly for \(s\) in every compact \(K\subset\{\Re s>1\}\).  Hence
\begin{equation}
\lim_{t\to0}
\frac{A(s)^2-B(s,t)B(s,-t)}{t^2}
=
A(s)M_2(s)-M_1(s)^2.
\label{eq:t0_removable}
\end{equation}
At infinity, \(|B(s,t)|\le\sum_p|u_p(s)|\), so the integrand is
\(O_K(t^{-2})\).  These two estimates give a second direct proof of absolute
integrability in \eqref{eq:I_definition}.

\begin{proposition}[Exact first transverse derivative]
\label{prop:B_first_transverse}
For \(\Re s>1\),
\begin{equation}
\boxed{
\partial_tB(s,0)=-i\frac{\zeta'}{\zeta}(s).}
\label{eq:B_t_derivative_logzeta}
\end{equation}
For every integer \(a\ge0\),
\begin{equation}
\boxed{
\partial_s^a\partial_tB(s,0)
=-i\left(\frac{\zeta'}{\zeta}\right)^{(a)}(s).}
\label{eq:B_mixed_derivative_logzeta}
\end{equation}
Both identities extend meromorphically wherever the right-hand side does.
\end{proposition}

\begin{proof}
Termwise differentiation is justified by
Theorem~\ref{thm:O_initial_domain}; hence
\begin{equation}
\partial_tB(s,0)
=i\sum_p\frac{\log p}{p^s-1}.
\label{eq:B_t_prime_sum}
\end{equation}
The absolutely convergent Euler-product logarithmic derivative is
\begin{equation}
-\frac{\zeta'}{\zeta}(s)
=\sum_p\sum_{a\ge1}(\log p)p^{-as}
=\sum_p\frac{\log p}{p^s-1},
\label{eq:logderivative_prime_sum}
\end{equation}
which proves \eqref{eq:B_t_derivative_logzeta}.  Differentiating in \(s\)
termwise on compacta gives \eqref{eq:B_mixed_derivative_logzeta}.  The final
statement follows from uniqueness of meromorphic continuation.
\end{proof}

\begin{proposition}[Higher pure transverse derivatives]
\label{prop:B_higher_transverse}
For every integer \(j\ge1\) and \(\Re s>1\),
\begin{equation}
\partial_t^jB(s,0)
=(-i)^j\sum_{n\ge1}c_j(n)
\left(\frac{\zeta'}{\zeta}\right)^{(j-1)}(ns),
\label{eq:B_higher_t_derivatives}
\end{equation}
where
\begin{equation}
c_j(n)=\sum_{d\mid n}\mu(d)d^{j-1}
=\prod_{p\mid n}(1-p^{j-1}).
\label{eq:cj_definition}
\end{equation}
In particular, \(c_1(n)=\delta_{n1}\), explaining why only the first
transverse derivative is a complete replica filter.
\end{proposition}

\begin{proof}
Insert the prime-zeta M\"obius inversion
\begin{equation}
P(w)=\sum_{m\ge1}\frac{\mu(m)}m\Log\zeta(mw),
\qquad \Re w>1,
\label{eq:P_mobius_inversion_v8}
\end{equation}
into \eqref{eq:B_prime_zeta}.  Absolute convergence in \(\Re s>1\) permits
termwise differentiation, giving
\begin{equation}
\partial_t^jB(s,0)
=(-i)^j\sum_{k,m\ge1}\mu(m)m^{j-1}
L^{(j-1)}(mks),
\label{eq:B_higher_double}
\end{equation}
with \(L=\zeta'/\zeta\).  Grouping by \(n=mk\) is legitimate because the
same differentiated series is absolutely convergent on compacta.  The inner
divisor sum is \(c_j(n)\), proving \eqref{eq:B_higher_t_derivatives}.  Since
\(\mu\) is supported on square-free integers, multiplicativity gives the
Euler factor \(1-p^{j-1}\) in \eqref{eq:cj_definition}.  For \(j=1\), the
standard identity \(\sum_{d\mid n}\mu(d)=\delta_{n1}\) completes the proof.
\end{proof}

\section{Continuation of the Fourier energy and complete replica accounting}
\label{sec:mobius_replica_filter}

The prime-zeta identity \eqref{eq:P_mobius_inversion_v8} is classical
\cite{Froberg1968}.  What is needed here is not merely the location of its
branch points, but a proof that the energy integral can be continued while
retaining control of all sums, derivatives, contour singularities, and
boundary terms.

\subsection{Grouped log-zeta expansion and its normal tail}
\label{subsec:B_grouped_continuation}

For \(\Re s>1\), regrouping the absolutely convergent double series gives
\begin{equation}
\boxed{
B(s,t)=
\sum_{n\ge1}\sum_{d\mid n}\frac{\mu(d)}d
\Log\zeta(ns-idt).}
\label{eq:B_grouped_logzeta}
\end{equation}

\begin{lemma}[Normal convergence of the large-dilation tail]
\label{lem:B_tail_normal}
Fix \(s_0\) with \(\sigma_0:=\Re s_0>0\).  There are a neighbourhood
\(U\) of \(s_0\) and an integer \(N_0\) such that the part of
\eqref{eq:B_grouped_logzeta} with \(n>N_0\) is single-valued and holomorphic
in \(s\in U\), is \(C^\infty\) in real \(t\), and all mixed derivatives of
fixed finite order converge normally on \(U\times\mathbb R\) after
multiplication by any fixed polynomially decaying weight in \(t\).
\end{lemma}

\begin{proof}
Choose \(U\) so that \(\Re s\ge\sigma_0/2\) on \(U\), and choose
\(N_0\) with \((N_0+1)\sigma_0/2>1+\delta\) for some \(\delta>0\).  Then for
\(n>N_0\),
\begin{equation}
\Re(ns-idt)=n\Re s\ge1+\delta.
\label{eq:tail_right_halfplane}
\end{equation}
In that half-plane the Euler product defines the principal holomorphic
logarithm and
\begin{equation}
|\Log\zeta(w)|
\le
\sum_{m\ge2}\frac{C_\delta}{m^{\Re w}}
\le C'_\delta 2^{-\Re w}.
\label{eq:logzeta_exponential_tail}
\end{equation}
Every fixed derivative of \(\Log\zeta(w)\) satisfies the same estimate with
an additional polynomial factor in \(\Re w\), obtained by differentiating
the absolutely convergent von Mangoldt series.  The divisor coefficient
satisfies
\begin{equation}
\sum_{d\mid n}\frac1d\le \tau(n),
\label{eq:divisor_coeff_bound}
\end{equation}
and \(\tau(n)\) has subpolynomial growth.  Hence the resulting tail is
majorized by a convergent series of the form
\(\sum_{n>N_0}n^C\tau(n)2^{-n\sigma_0/2}\), uniformly in real \(t\).
The Weierstrass theorem proves normal convergence and permits all asserted
finite-order differentiations.
\end{proof}

\begin{definition}[Continued branch of \(\mathcal I\)]
\label{def:I_sharp}
A branch \(\mathcal I^\sharp\) near a point \(s_0\) with \(\Re s_0>0\) is
obtained by analytic continuation from \(\Re s>1\) of the finitely many
head terms \(n\le N_0\) in \eqref{eq:B_grouped_logzeta}, while the normal
tail of Lemma~\ref{lem:B_tail_normal} is kept on its single-valued branch.
Whenever a moving logarithmic singularity meets the real \(t\)-axis, the
\(t\)-contour is deformed continuously around it.  Different continuation
paths may change lower-order branch constants; the principal coefficients
computed below are independent of those choices.
\end{definition}

\begin{lemma}[Convergence of the continued contour integral]
\label{lem:I_contour_convergence}
Every local branch from Definition~\ref{def:I_sharp} can be represented by a
locally finite deformed \(t\)-contour on which the continued integrand is
absolutely integrable at infinity and locally integrable at every logarithmic
branch point.  The contour may be chosen so that the contribution of all
indentations above height \(T\) tends to zero as \(T\to\infty\).
\end{lemma}

\begin{proof}
Only the finite head \(n\le N_0\) can contribute zeros or the pole of
\(\zeta\) to the \(t\)-contour.  For each fixed head pair \((n,d)\), those
points are obtained from
\begin{equation}
ns-idt=\rho'
\label{eq:head_contour_zero_equation}
\end{equation}
with \(\rho'\) a zero or pole of \(\zeta\).  The classical zero-counting
estimate
\begin{equation}
N_\zeta(T)=O(T\log T)
\label{eq:zeta_zero_counting_bound}
\end{equation}
is unconditional; see \cite{Titchmarsh1986}.  Hence the number of contour
singularities of height at most \(T\) is \(O(T\log T)\).

Near a zero of multiplicity \(m\), local factorization gives
\begin{equation}
\Log\zeta(\rho'+w)=m\Log w+h(w),
\label{eq:local_logzeta_factorization}
\end{equation}
with \(h\) holomorphic.  Thus each singularity is logarithmic and
\(|\Log|t-t_0||^2\) is locally integrable.  Choose pairwise compatible
indentations with radii no larger than \((1+|t_0|)^{-4}\), shrinking them
further when necessary to avoid overlaps.  On the complementary contour,
standard bounds obtained from the Hadamard product for \(\xi\) imply
\begin{equation}
|\Log\zeta(\sigma+it)|
\le C_K\log^2(2+|t|)
\label{eq:logzeta_contour_growth}
\end{equation}
for the finitely many vertical strips involved, after the chosen disks have
been removed; this standard estimate is sufficient here and follows from the
zero count and local factorization, again as in \cite{Titchmarsh1986}.
Consequently
\begin{equation}
\frac{|A(s)^2-B(s,t)B(s,-t)|}{|t|^2}
\le
C_K\frac{1+\log^4(2+|t|)}{|t|^2}
\label{eq:I_tail_contour_majorant}
\end{equation}
away from the indentation arcs.  The right-hand side is integrable.

On an indentation of radius \(r_0\) around height \(t_0\), the logarithmic
factor is \(O(|\log r_0|+\log(2+|t_0|))\), the arc length is \(O(r_0)\), and
the weight is \(O(t_0^{-2})\) for \(|t_0|\ge1\).  The sum of arc
contributions is therefore dominated by
\begin{equation}
\sum_{t_0}
\frac{r_0\,[1+|\log r_0|+\log^2(2+|t_0|)]}{1+t_0^2},
\label{eq:indentation_sum_majorant}
\end{equation}
which converges because \(r_0\le(1+|t_0|)^{-4}\) and the number of
\(t_0\)'s in a unit interval is \(O(\log(2+|t_0|))\) on average by
\eqref{eq:zeta_zero_counting_bound}.  Local logarithmic integrability handles
the finitely many points with \(|t_0|<1\).  This proves absolute convergence
and the vanishing of the indentation tail.
\end{proof}

\begin{theorem}[Existence of the continued energy germ]
\label{thm:I_continuation_exists}
For every \(s_0\) with \(\Re s_0>0\), every branch obtained from
Definition~\ref{def:I_sharp} defines a sectorial analytic germ
\(\mathcal I^\sharp\) at \(s_0\) away from its local logarithmic cuts.  It
coincides with \(\mathcal I\) in the common domain \(\Re s>1\).  No RH,
zero-simplicity, or nonresonance assumption is used.
\end{theorem}

\begin{proof}
Lemma~\ref{lem:B_tail_normal} reduces the continuation to finitely many
logarithms of meromorphic functions.  On every compact subarc of the deformed
contour not meeting an indentation, these logarithms are holomorphic in
\(s\), and their products are dominated by the integrable majorant
\eqref{eq:I_tail_contour_majorant}.  On each indentation, use the local
factorization \eqref{eq:local_logzeta_factorization}; after translating the
local contour, the integral is a finite linear combination of parameter
integrals of \(\Log(a(s-s_*)+b\tau)\) and products of two such logarithms.
These are holomorphic in a slit sector by differentiation under the finite
local integral.  The normal convergence of the sum of indentation
contributions follows from \eqref{eq:indentation_sum_majorant}.  Therefore the
sum of the compact-contour part, indentation parts, and tail is analytic in
the slit sector.  Continuation from \(\Re s>1\) fixes the branch on overlaps,
so uniqueness of analytic continuation gives agreement with
\eqref{eq:I_definition} in the initial half-plane.
\end{proof}

\subsection{All contributions arriving at a scaled point}
\label{subsec:all_replica_contributions}

Fix \(s_0\) with \(0<\Re s_0<1\), and define
\begin{equation}
\mathcal R(s_0)
=
\{N\in\mathbb N:\zeta(Ns_0)=0\},
\qquad
q_N=\operatorname{ord}_{w=Ns_0}\zeta(w).
\label{eq:Rset_qN}
\end{equation}
The set \(\mathcal R(s_0)\) is finite, because a nontrivial zero has real
part in \((0,1)\), so \(N\Re s_0<1\).
Define
\begin{align}
C(s_0)
&=
\sum_{N\in\mathcal R(s_0)}q_N\frac{\varphi(N)}N,
\label{eq:C_s0_definition}\\
D(s_0)
&=
\sum_{N\in\mathcal R(s_0)}q_N\frac{\varphi(N)}N\log N,
\label{eq:D_s0_definition}
\end{align}
and, for real \(u\),
\begin{equation}
G_{s_0}(u)
=
\sum_{N\in\mathcal R(s_0)}q_N
\sum_{d\mid N}\frac{\mu(d)}d
\Log\left(1-i\frac dN u\right).
\label{eq:G_s0_definition}
\end{equation}

\begin{lemma}[Complete endpoint singular block]
\label{lem:complete_endpoint_block}
With \(z=s-s_0\), the complete contribution of all log-zeta terms singular
at \((s,t)=(s_0,0)\) is
\begin{equation}
S(z,t)
=
\sum_{N\in\mathcal R(s_0)}q_N
\sum_{d\mid N}\frac{\mu(d)}d
\Log(Nz-idt).
\label{eq:S_complete_endpoint}
\end{equation}
After subtracting \(S\), the remainder \(R(z,t)=B(s_0+z,t)-S(z,t)\) is
holomorphic in \(z\) and twice continuously differentiable in real \(t\) in
some product neighbourhood of \((0,0)\).
\end{lemma}

\begin{proof}
In the finite head of \eqref{eq:B_grouped_logzeta}, a term indexed by
\((n,d)\) is singular at \((s_0,0)\) precisely when \(\zeta(ns_0)=0\).
Thus \(n=N\in\mathcal R(s_0)\), while every divisor \(d\mid N\) contributes
the same zero at \(t=0\).  Local factorization
\eqref{eq:local_logzeta_factorization} then yields exactly
\eqref{eq:S_complete_endpoint}.  Since the head is finite and all endpoint
singular factors have been removed, discreteness of the zero set permits a
small product neighbourhood containing no further head singularity.  The
normal tail is smooth by Lemma~\ref{lem:B_tail_normal}.  Hence the remainder
has the asserted regularity.
\end{proof}

\begin{lemma}[Interior contour singularities are lower order]
\label{lem:interior_lower_order}
Let \(t_0>0\) be any logarithmic singularity of a continued head term.  Its
local contribution to \(\mathcal I^\sharp(s_0+z)\), including products of
two logarithmic factors, is \(O(1)+O(z\log^2 z)\) in a slit sector.  In
particular, no interior singularity can contribute either
\(z^{-1}\log z\) or \(z^{-1}\).
\end{lemma}

\begin{proof}
Write \(t=t_0+x\).  The weight is analytic there:
\begin{equation}
\frac1{t^2}=\frac1{t_0^2}+O(x).
\label{eq:t0_weight_expansion}
\end{equation}
A local logarithmic factor has the form
\(\Log(az-ibx)+h(z,x)\), with \(h\) holomorphic.  On a symmetric local
contour of fixed radius, set \(x=zu\) only in the central part
\(|x|\le c|z|\); the Jacobian contributes a factor \(z\), while the weight
has no inverse power of \(z\).  This region is therefore
\(O(z\log^2 z)\).  On the complementary part \(c|z|<|x|<\delta\), expand
\begin{equation}
\Log(az-ibx)
=
\Log(-ibx)+O\!\left(\frac{|z|}{|x|}\right)
\label{eq:interior_log_expansion}
\end{equation}
uniformly on the chosen slit sector.  The error integrates to
\(O(z|\log z|)\), and products of two logarithms to
\(O(z\log^2 z)\).  The zeroth-order terms give a finite constant because
\(\Log|x|\) and \(\Log^2|x|\) are locally integrable.  Therefore no negative
power of \(z\) is produced.
\end{proof}

\begin{theorem}[Exact coefficient of the logarithmic-over-pole term]
\label{thm:principal_replica_filter}
For every continued branch \(\mathcal I^\sharp\) at \(s_0\), the coefficient
of \(\log z/z\), \(z=s-s_0\), is
\begin{equation}
\boxed{
[\log z/z]\,\mathcal I^\sharp(s_0+z)
=-\pi q_1 C(s_0),}
\label{eq:general_log_over_pole_coefficient}
\end{equation}
where \(q_1=0\) when \(1\notin\mathcal R(s_0)\).  Consequently,
\begin{equation}
\boxed{
[\log z/z]\,\mathcal I^\sharp=0
\iff
\zeta(s_0)\ne0.}
\label{eq:true_zero_detector_I}
\end{equation}
No singularity with \(t_0>0\) can change this coefficient.
\end{theorem}

\begin{proof}
From \eqref{eq:S_complete_endpoint}, for \(t=zu\) one has
\begin{equation}
S(z,zu)=C(s_0)\Log z+D(s_0)+G_{s_0}(u),
\qquad
S(z,0)=C(s_0)\Log z+D(s_0).
\label{eq:S_scaled_general}
\end{equation}
The singular--singular numerator equals
\begin{align}
&S(z,0)^2-S(z,zu)S(z,-zu)
\notag\\
&= -\bigl(C\Log z+D\bigr)
\bigl(G(u)+G(-u)\bigr)-G(u)G(-u).
\label{eq:SS_scaled_numerator}
\end{align}
The contour rotation from the ray \(u=t/z\) to the positive real axis is
legitimate in every closed subsector because the branch points
\(u=\pm iN/d\) lie on the imaginary axis; the large circular arc contributes
zero since the integrand is \(O(\log^2|u|/|u|^2)\).  Therefore the
coefficient of \(\Log z/z\) is
\begin{equation}
-C(s_0)
\int_0^\infty\frac{G(u)+G(-u)}{u^2}\,du.
\label{eq:alpha_integral_general}
\end{equation}
For real \(u\),
\begin{equation}
G(u)+G(-u)
=
\sum_{N\in\mathcal R(s_0)}q_N
\sum_{d\mid N}\frac{\mu(d)}d
\log\left(1+\frac{d^2u^2}{N^2}\right).
\label{eq:G_plus_Gminus}
\end{equation}
Integration by parts gives, for \(a>0\),
\begin{equation}
\int_0^\infty\frac{\log(1+a^2u^2)}{u^2}\,du=\pi a.
\label{eq:log_quadratic_integral}
\end{equation}
The boundary term vanishes at zero and infinity.  Thus
\begin{align}
\int_0^\infty\frac{G(u)+G(-u)}{u^2}\,du
&=
\pi\sum_{N\in\mathcal R(s_0)}\frac{q_N}{N}
\sum_{d\mid N}\mu(d)
\notag\\
&=\pi q_1,
\label{eq:G_sum_mobius_collapse}
\end{align}
because \(\sum_{d\mid N}\mu(d)=\delta_{N1}\).  This proves
\eqref{eq:general_log_over_pole_coefficient}.  Since all summands in
\(C(s_0)\) are positive, \(C(s_0)>0\) whenever
\(\mathcal R(s_0)\ne\varnothing\), so the coefficient vanishes exactly when
\(q_1=0\), equivalently \(\zeta(s_0)\ne0\).  Finally,
Lemma~\ref{lem:interior_lower_order} excludes any restoration of the same
order by interior contour singularities.
\end{proof}

\begin{corollary}[Blockwise M\"obius cancellation]
\label{cor:blockwise_mobius_filter}
Let \(\rho\) be a zero of multiplicity \(m_\rho\), let \(r\ge1\), and set
\(s_0=\rho/r\).  The singular--singular block attached specifically to
\((\rho,r)\) contributes
\begin{equation}
\boxed{
-\pi m_\rho^2\frac{\varphi(r)}{r^2}
\left(\sum_{d\mid r}\mu(d)\right)
\frac{\Log(s-s_0)}{s-s_0}}
\label{eq:blockwise_requested_filter_formula}
\end{equation}
to the logarithmic-over-pole scale of \(\mathcal I^\sharp\).  Hence this
block is zero for every \(r>1\), whereas for \(r=1\) it equals
\begin{equation}
-\pi m_\rho^2\frac{\Log(s-\rho)}{s-\rho}.
\label{eq:r1_surviving_block}
\end{equation}
\end{corollary}

\begin{proof}
Restrict \eqref{eq:S_complete_endpoint} to a single dilation index \(r\).
Then \(C=m_\rho\varphi(r)/r\), while the integral
\eqref{eq:G_sum_mobius_collapse} for that single block is
\(\pi m_\rho r^{-1}\sum_{d\mid r}\mu(d)\).  Their negative product gives
\eqref{eq:blockwise_requested_filter_formula}.  The divisor identity
\(\sum_{d\mid r}\mu(d)=0\) for \(r>1\) and equals one for \(r=1\), proving
both conclusions.
\end{proof}

\subsection{Pure replicas: the next coefficient and the complete local germ}
\label{subsec:pure_replica_beta}

Assume now that \(\mathcal R(s_0)\ne\varnothing\) but
\(\zeta(s_0)\ne0\), so \(q_1=0\).  Define
\begin{equation}
\beta(s_0)
:=-\int_0^\infty\frac{|G_{s_0}(u)|^2}{u^2}\,du.
\label{eq:beta_s0_energy}
\end{equation}

\begin{lemma}[Convergence and strict negativity of the collision energy]
\label{lem:beta_s0_negative}
Under \(q_1=0\), the integral \eqref{eq:beta_s0_energy} converges and
\begin{equation}
\boxed{\beta(s_0)<0.}
\label{eq:beta_s0_negative}
\end{equation}
\end{lemma}

\begin{proof}
At zero,
\begin{align}
G'_{s_0}(0)
&=-i\sum_{N\in\mathcal R(s_0)}\frac{q_N}{N}
\sum_{d\mid N}\mu(d)
=0,
\label{eq:Gprime_zero_collision}
\end{align}
because every \(N\in\mathcal R(s_0)\) is greater than one.  Since
\(G(0)=0\), it follows that \(G(u)=O(u^2)\) and the integrand is \(O(u^2)\)
at zero.  At infinity,
\begin{equation}
G(u)=C(s_0)\log u+O(1),
\label{eq:G_infinity_collision}
\end{equation}
so the integrand is \(O(\log^2u/u^2)\).  Thus the integral converges.

It remains to exclude \(G\equiv0\).  Let
\(N_{\max}=\max\mathcal R(s_0)\).  Differentiating \eqref{eq:G_s0_definition}
meromorphically in \(u\), the \(d=1\) term belonging to \(N_{\max}\) has a
simple pole at \(u=-iN_{\max}\).  Every pole from a term with
\(M<N_{\max}\) has modulus \(M/d\le M<N_{\max}\), and terms with
\(N=N_{\max}\), \(d>1\), have modulus \(N_{\max}/d<N_{\max}\).  Hence that
pole cannot cancel.  Therefore \(G\not\equiv0\), its squared modulus is
positive on a set of positive measure, and \eqref{eq:beta_s0_negative}
follows.
\end{proof}

\begin{theorem}[Complete local expansion at a pure replica point]
\label{thm:pure_replica_full_germ}
Suppose \(\mathcal R(s_0)\ne\varnothing\) and \(\zeta(s_0)\ne0\).  Every
continued branch satisfies, in every closed slit subsector,
\begin{equation}
\boxed{
\mathcal I^\sharp(s_0+z)
=
\frac{\beta(s_0)}{z}
+
\pi C(s_0)\frac{\zeta'}{\zeta}(s_0)\Log z
+O(1).}
\label{eq:pure_replica_full_expansion}
\end{equation}
In particular the pole coefficient is nonzero and no term
\(\Log z/z\) is present.
\end{theorem}

\begin{proof}
Decompose \(B=S+R\) as in Lemma~\ref{lem:complete_endpoint_block}.  The
singular--singular part is given by \eqref{eq:SS_scaled_numerator}.  Since
\(q_1=0\), equation \eqref{eq:G_sum_mobius_collapse} makes the entire term
linear in \(C\Log z+D\) vanish.  Therefore its negative-power contribution is
exactly
\begin{equation}
\frac1z\left(-\int_0^\infty\frac{G(u)G(-u)}{u^2}\,du\right)
=
\frac{\beta(s_0)}z,
\label{eq:SS_beta_exact}
\end{equation}
because on the positive real axis \(G(-u)=\overline{G(u)}\).  The contour
rotation used in Theorem~\ref{thm:principal_replica_filter} shows the same
coefficient in every slit subsector.

Near \((0,0)\), write
\begin{equation}
R(z,\pm t)=R_0(z)\pm h_1t+O(t^2+|z|t),
\label{eq:R_local_taylor}
\end{equation}
with \(h_1=\partial_tR(0,0)\).  The constant cross term is proportional to
\begin{equation}
\int_0^\infty\frac{G(u)+G(-u)}{u^2}\,du=0,
\label{eq:constant_cross_zero}
\end{equation}
so it produces neither \(z^{-1}\) nor \(z^{-1}\Log z\).  For the odd linear
term, one obtains
\begin{equation}
h_1\int\frac{S(z,t)-S(z,-t)}t\,dt.
\label{eq:odd_cross_integral}
\end{equation}
As \(u\to+\infty\),
\begin{equation}
G(u)-G(-u)=-i\pi C(s_0)+O(u^{-1}),
\label{eq:G_odd_infinity}
\end{equation}
so integration from \(u=1\) to \(u=\delta/|z|\) gives
\begin{equation}
i\pi C(s_0)h_1\Log z+O(1).
\label{eq:odd_cross_log_term}
\end{equation}
All Taylor terms of order \(t^2\) or higher are integrable at the endpoint
and contribute \(O(1)\); terms involving \(z\) are also \(O(1)\).  By
Lemma~\ref{lem:interior_lower_order}, every interior contour singularity is
bounded up to \(O(z\log^2z)\), and the contour tail is normally convergent by
Lemma~\ref{lem:I_contour_convergence}.  Thus no additional unbounded term of
order \(z^{-1}\), \(z^{-1}\Log z\), or \(\Log z\) is present.

Finally, differentiating \(S\) at \(t=0\) gives
\begin{equation}
\partial_tS(z,0)
=-\frac{i}{z}
\sum_{N\in\mathcal R(s_0)}\frac{q_N}{N}
\sum_{d\mid N}\mu(d)=0,
\label{eq:S_t_derivative_pure_replica}
\end{equation}
so \(h_1=\partial_tB(s_0,0)\).  Proposition~\ref{prop:B_first_transverse}
and meromorphic continuation give
\begin{equation}
h_1=-i\frac{\zeta'}{\zeta}(s_0),
\label{eq:h1_logderivative}
\end{equation}
because \(\zeta(s_0)\ne0\).  Substitution into
\eqref{eq:odd_cross_log_term} proves \eqref{eq:pure_replica_full_expansion}.
\end{proof}

\subsection{Single-replica divisor energy}
\label{subsec:single_replica_energy}

For a single zero \(\rho\) of multiplicity \(m_\rho\) and one dilation
index \(r>1\), define
\begin{equation}
G_r(u)=
\sum_{d\mid r}\frac{\mu(d)}d
\Log\left(1-i\frac d r u\right).
\label{eq:Gr_definition}
\end{equation}

\begin{proposition}[Replica energy, finite formula, and sign]
\label{prop:beta_r}
Define
\begin{equation}
\boxed{
\beta_r=-\int_0^\infty\frac{|G_r(u)|^2}{u^2}\,du.}
\label{eq:beta_r_energy}
\end{equation}
Then \(\beta_r<0\) for every \(r>1\), and
\begin{equation}
\boxed{
\beta_r=
-\frac{\pi}{r}
\sum_{d,e\mid r}\frac{\mu(d)\mu(e)}{de}
\Big[(d+e)\log(d+e)-d\log d-e\log e\Big].}
\label{eq:beta_r_finite}
\end{equation}
For an isolated pure replica point \(s_0=\rho/r\), the pole coefficient in
\eqref{eq:pure_replica_full_expansion} is \(m_\rho^2\beta_r\).
\end{proposition}

\begin{proof}
The convergence and strict sign follow from the proof of
Lemma~\ref{lem:beta_s0_negative}; for a single \(r>1\),
\(G_r(0)=G_r'(0)=0\) and \(G_r\not\equiv0\).

For the finite formula, set
\begin{equation}
F(a,b)=\Re\int_0^\infty
\frac{\Log(1-iau)\Log(1+ibu)}{u^2}\,du,
\qquad a,b>0.
\label{eq:F_ab_definition}
\end{equation}
The integral converges absolutely after taking the real part of the paired
expression.  Differentiate first with respect to \(a\) and \(b\), justified
on compact parameter rectangles by dominated convergence after splitting
\((0,1)\cup(1,\infty)\).  One obtains
\begin{equation}
\frac{\partial^2 F}{\partial a\partial b}
=\Re\int_0^\infty
\frac{du}{(1-iau)(1+ibu)}
=\frac{\pi}{a+b}.
\label{eq:F_cross_derivative}
\end{equation}
The last equality follows from partial fractions and evaluation of the
logarithms at a finite cutoff before taking the limit.  Since
\(F(a,0)=F(0,b)=0\), two integrations give
\begin{equation}
F(a,b)=
\pi\big[(a+b)\log(a+b)-a\log a-b\log b\big].
\label{eq:F_ab_closed}
\end{equation}
Expanding \(|G_r(u)|^2\), using
\(a=d/r\), \(b=e/r\), and the homogeneity of
\eqref{eq:F_ab_closed}, gives \eqref{eq:beta_r_finite}.
\end{proof}

\begin{corollary}[Radical law]
\label{cor:beta_radical_law}
For every \(r>1\),
\begin{equation}
\boxed{
\beta_r=
\frac{\operatorname{rad}(r)}r
\beta_{\operatorname{rad}(r)}.}
\label{eq:beta_radical_law}
\end{equation}
Equivalently, \(r\beta_r\) depends only on the square-free kernel of \(r\).
\end{corollary}

\begin{proof}
Put \(q=\operatorname{rad}(r)\).  In the rescaled form of
\eqref{eq:beta_r_energy}, let \(u=rx\).  Then
\begin{equation}
\beta_r
=-\frac1r\int_0^\infty
\frac{1}{x^2}
\left|\sum_{d\mid r}\frac{\mu(d)}d\Log(1-idx)\right|^2dx.
\label{eq:beta_r_rescaled}
\end{equation}
Only square-free divisors contribute, and those are exactly the divisors of
\(q\).  The integral in \eqref{eq:beta_r_rescaled} is therefore the same as
the corresponding integral for \(q\), while the exterior factors are
\(1/r\) and \(1/q\).  Hence \(r\beta_r=q\beta_q\), which is
\eqref{eq:beta_radical_law}.
\end{proof}

\begin{corollary}[Prime and prime-power replica asymptotics]
\label{cor:beta_prime_asymptotic}
For a prime \(p\),
\begin{equation}
\boxed{
\beta_p=
-\frac{2\pi}{p^2}
\left[p\log p-(p+1)\log(p+1)+(p+1)\log2\right].}
\label{eq:beta_p_formula_exact}
\end{equation}
As \(p\to\infty\),
\begin{equation}
\beta_p
=-\frac{2\pi\log2}{p}
+O\!\left(\frac{\log p}{p^2}\right).
\label{eq:beta_p_asymptotic}
\end{equation}
For every integer \(a\ge1\),
\begin{equation}
\boxed{\beta_{p^a}=p^{1-a}\beta_p.}
\label{eq:beta_prime_power}
\end{equation}
\end{corollary}

\begin{proof}
In \eqref{eq:beta_r_finite}, for \(r=p\) the only divisors are \(1,p\).
Writing out the four terms and collecting logarithms gives
\eqref{eq:beta_p_formula_exact}.  Expanding
\((p+1)\log(p+1)=p\log p+\log p+1+O(p^{-1})\) yields
\eqref{eq:beta_p_asymptotic}; the weaker displayed error is sufficient and
uniform.  Finally, \(\operatorname{rad}(p^a)=p\), so
Corollary~\ref{cor:beta_radical_law} gives \eqref{eq:beta_prime_power}.
\end{proof}

\subsection{Actual zeta zeros, \(\mathcal O\), and monodromy}
\label{subsec:O_vs_I_singularities}

\begin{theorem}[Survival at a genuine zero]
\label{thm:r1_survival}
Let \(\rho\) be a zero of multiplicity \(m_\rho\).  For every continued
branch,
\begin{equation}
\mathcal I^\sharp(\rho+z)
=
-\pi m_\rho C(\rho)\frac{\Log z}{z}
+O\!\left(\frac1z\right),
\label{eq:I_true_zero_general}
\end{equation}
where
\begin{equation}
C(\rho)
=
\sum_{N\in\mathcal R(\rho)}q_N\frac{\varphi(N)}N
\ge m_\rho.
\label{eq:C_rho_true_zero}
\end{equation}
If no scaled zero other than \(N=1\) collides at \(\rho\), then
\begin{equation}
\boxed{
\mathcal I^\sharp(\rho+z)
=-\pi m_\rho^2\frac{\Log z}{z}+O(z^{-1}).}
\label{eq:I_true_zero_noncollision}
\end{equation}
\end{theorem}

\begin{proof}
At \(s_0=\rho\), one has \(q_1=m_\rho>0\).  Theorem
\ref{thm:principal_replica_filter} gives the coefficient
\(-\pi q_1C(\rho)\), proving \eqref{eq:I_true_zero_general}.  All other
terms are at most \(O(z^{-1})\): the singular--singular expansion
\eqref{eq:SS_scaled_numerator} has no stronger scale, cross terms contain at
most one singular logarithm, and interior singularities are lower order by
Lemma~\ref{lem:interior_lower_order}.  In the absence of collisions,
\(C(\rho)=m_\rho\), yielding \eqref{eq:I_true_zero_noncollision}.
\end{proof}

\begin{corollary}[Correct singular orders for \(\mathcal O\)]
\label{cor:O_correct_singular_orders}
Let \(\zeta(\rho+z)=a_mz^m+O(z^{m+1})\), \(a_m\ne0\), with
\(m=m_\rho\).  In the noncolliding case,
\begin{equation}
\boxed{
\mathcal O(\rho+z)
=-\frac{a_m m^2}{2}z^{m-1}\Log z+O(z^{m-1}).}
\label{eq:O_true_zero_order}
\end{equation}
For a simple zero,
\begin{equation}
\boxed{
\mathcal O(\rho+z)
=-\frac{\zeta'(\rho)}2\Log z+O(1).}
\label{eq:O_simple_zero_log}
\end{equation}
At a pure replica \(s_0=\rho/r\), \(r>1\), with no collision with an
actual zero at \(s_0\),
\begin{equation}
\boxed{
\mathcal O(s_0+z)
=
\frac{\zeta(s_0)\,\beta(s_0)}{2\pi z}
+
\frac{C(s_0)\zeta'(s_0)}2\Log z
+O(1).}
\label{eq:O_pure_replica_expansion}
\end{equation}
\end{corollary}

\begin{proof}
Use the identity
\(\mathcal O=\zeta\mathcal I/(2\pi)\), initially valid for \(\Re s>1\)
and then on every continued branch by analytic continuation.  Multiplying
\eqref{eq:I_true_zero_noncollision} by
\(a_mz^m+O(z^{m+1})\) gives \eqref{eq:O_true_zero_order}; the case \(m=1\)
is \eqref{eq:O_simple_zero_log}.  At a pure replica, \(\zeta(s_0)\ne0\), so
multiplication of \eqref{eq:pure_replica_full_expansion} gives
\eqref{eq:O_pure_replica_expansion}.
\end{proof}

\begin{warning}[The requested \(\log z/z\) formula belongs to \(\mathcal I\)]
\label{warn:requested_formula_normalization}
A formula of the form
\begin{equation}
C_{\rho,r}
\left(\sum_{d\mid r}\mu(d)\right)
\frac{\Log(s-\rho/r)}{s-\rho/r}
\label{eq:requested_formula_shape}
\end{equation}
is correct for the normalized energy \(\mathcal I=2\pi\mathcal O/\zeta\),
with
\begin{equation}
C_{\rho,r}
=-\pi m_\rho^2\frac{\varphi(r)}{r^2}
\label{eq:C_rho_r_requested}
\end{equation}
for the individual singular--singular block.  It is not the correct local
order of \(\mathcal O\) itself at \(r=1\), because \(\zeta\) vanishes there
and cancels the pole.  Equations \eqref{eq:O_true_zero_order} and
\eqref{eq:O_pure_replica_expansion} are the correct statements for
\(\mathcal O\).
\end{warning}

\begin{corollary}[Monodromy detector for genuine zeros]
\label{cor:monodromy_zero_detector}
Let \(\Delta\mathcal I^\sharp\) denote one positive turn of the local branch.
Then
\begin{equation}
\operatorname*{Res}_{s=s_0}\Delta\mathcal I^\sharp(s)
=
-2\pi^2 i\,q_1C(s_0).
\label{eq:monodromy_residue_general}
\end{equation}
Consequently
\begin{equation}
\boxed{
\zeta(s_0)=0
\iff
\Delta\mathcal I^\sharp
\text{ has a pole at }s_0.}
\label{eq:monodromy_detector_equivalence}
\end{equation}
\end{corollary}

\begin{proof}
Only the term \(-\pi q_1C(s_0)\Log z/z\) changes by a term of order
\(z^{-1}\) after \(\Log z\mapsto\Log z+2\pi i\).  Its change is
\(-2\pi^2i q_1C(s_0)/z\).  All pole terms without a logarithm are
single-valued, and all lower logarithmic terms have bounded monodromy.  Since
\(C(s_0)>0\), the residue is nonzero exactly when \(q_1>0\), i.e. when
\(\zeta(s_0)=0\).
\end{proof}

\subsection{Analytic checklist for the replica theorem}
\label{subsec:replica_checklist}

\begin{proposition}[Eight-point analytic closure]
\label{prop:eight_point_closure}
The replica analysis above establishes the following items without RH,
zero-simplicity, or a nonresonance assumption:
\begin{enumerate}[label=(\arabic*)]
\item \(\mathcal O(s)\) and \(\mathcal I(s)\) exist absolutely for
\(\Re s>1\) and are holomorphic there.
\item Every local branch required for the replica analysis is obtained by
continuation from that half-plane, with a normally convergent large-dilation
tail.
\item All exchanges among prime sums, dilation sums, the \(t\)-integral, and
the differentiations used in the initial half-plane are justified by the
majorants displayed above.
\item The complete singular expansion through the unbounded terms is
\eqref{eq:pure_replica_full_expansion} at pure replica points and
\eqref{eq:I_true_zero_general} at actual zeros.
\item Every contribution reaching \(s_0\) at the endpoint is included by the
finite set \(\mathcal R(s_0)\); every interior contour contribution is covered
by Lemma~\ref{lem:interior_lower_order}.
\item For each individual index \(r>1\), the coefficient of
\(z^{-1}\Log z\) vanishes by
\(\sum_{d\mid r}\mu(d)=0\).
\item No other endpoint or interior term can restore a
\(z^{-1}\Log z\) contribution when \(\zeta(s_0)\ne0\).
\item The corresponding \(r=1\) term survives with nonzero coefficient
\(-\pi m_\rho^2\) in the noncolliding case, and more generally with
\(-\pi q_1C(s_0)\).
\end{enumerate}
\end{proposition}

\begin{proof}
Items (1)--(3) are Theorems~\ref{thm:O_initial_domain},
\ref{thm:O_prime_pair_fourier}, Lemmas~\ref{lem:B_tail_normal} and
\ref{lem:I_contour_convergence}, and Theorem~\ref{thm:I_continuation_exists}.
Item (4) is Theorems~\ref{thm:pure_replica_full_germ} and
\ref{thm:r1_survival}.  Item (5) is Lemmas~\ref{lem:complete_endpoint_block}
and \ref{lem:interior_lower_order}.  Item (6) is
Corollary~\ref{cor:blockwise_mobius_filter}.  Item (7) follows from
Theorem~\ref{thm:principal_replica_filter} together with
Lemma~\ref{lem:interior_lower_order}.  Item (8) is
Theorem~\ref{thm:r1_survival}.  Thus every item is a consequence of a proved
statement in the preceding subsections.
\end{proof}

\section{A two-variable Euler product and the covariance determinant}
\label{sec:hessian_covariance}

The determinant in \eqref{eq:K_variance_determinant} has an intrinsic
interpretation that does not involve a linearly M\"obius-weighted quantum
state.  For \(\Re s>1\) and \(u\) in a neighbourhood of zero define
\begin{equation}
\mathfrak Z(s,u)=\prod_p\left(1-e^u p^{-s}\right)^{-1},
\qquad
F(s,u)=\log\mathfrak Z(s,u).
\label{eq:two_variable_euler_product}
\end{equation}
Then \(\mathfrak Z(s,0)=\zeta(s)\).

\begin{theorem}[Hessian form of the curvature]
\label{thm:hessian_curvature}
For \(\Re s>1\),
\begin{align}
F_{uu}(s,0)&=S_0(s),
\label{eq:Fuu_S0}\\
-F_{su}(s,0)&=S_1(s),
\label{eq:Fsu_S1}\\
F_{ss}(s,0)&=S_2(s).
\label{eq:Fss_S2}
\end{align}
Consequently
\begin{equation}
\boxed{
\frac{2\mathcal K(s)}{\zeta(s)}
=\det\!\begin{pmatrix}
F_{uu}(s,0)&F_{us}(s,0)\\
F_{su}(s,0)&F_{ss}(s,0)
\end{pmatrix}.}
\label{eq:K_hessian_determinant}
\end{equation}
\end{theorem}

\begin{proof}
Differentiate
\begin{equation}
F(s,u)=-\sum_p\log(1-e^u p^{-s})
\label{eq:F_log_euler}
\end{equation}
termwise in a neighbourhood of \(u=0\), justified by normal convergence for
\(\Re s>1\).  The three derivatives are exactly
\eqref{eq:Sj_prime_sum}.  Equation~\eqref{eq:K_hessian_determinant} is then
Theorem~\ref{thm:curvature_dirichlet_series}.
\end{proof}

\begin{theorem}[Zeta-distribution covariance identity]
\label{thm:zeta_covariance_identity}
Let \(\sigma>1\) be real and let \(N_\sigma\) be the zeta-distributed integer
\begin{equation}
\mathbb P_\sigma(N_\sigma=n)=\frac{n^{-\sigma}}{\zeta(\sigma)}.
\label{eq:zeta_distribution}
\end{equation}
Then
\begin{align}
\operatorname{Var}_\sigma\!\bigl(\Omega(N_\sigma)\bigr)&=S_0(\sigma),
\label{eq:var_Omega_S0}\\
\operatorname{Cov}_\sigma\!\bigl(\Omega(N_\sigma),\log N_\sigma\bigr)&=S_1(\sigma),
\label{eq:cov_Omega_log_S1}\\
\operatorname{Var}_\sigma\!\bigl(\log N_\sigma\bigr)&=S_2(\sigma).
\label{eq:var_log_S2}
\end{align}
Therefore
\begin{equation}
\boxed{
\frac{2\mathcal K(\sigma)}{\zeta(\sigma)}
=\det\operatorname{Cov}_\sigma\!\begin{pmatrix}
\Omega(N_\sigma)\\ \log N_\sigma
\end{pmatrix}>0.}
\label{eq:K_covariance_determinant}
\end{equation}
\end{theorem}

\begin{proof}
Under \eqref{eq:zeta_distribution}, the prime valuations \(v_p(N_\sigma)\)
are independent geometric random variables with
\begin{equation}
\operatorname{Var}(v_p)=\frac{p^{-\sigma}}{(1-p^{-\sigma})^2}.
\label{eq:geometric_prime_variance}
\end{equation}
Since
\begin{equation}
\Omega(N_\sigma)=\sum_p v_p(N_\sigma),
\qquad
\log N_\sigma=\sum_p v_p(N_\sigma)\log p,
\label{eq:Omega_log_prime_valuations}
\end{equation}
independence gives \eqref{eq:var_Omega_S0}--\eqref{eq:var_log_S2}.
Strict positivity follows from
\begin{equation}
S_0S_2-S_1^2
=\sum_{p<q}w_p w_q\log^2(q/p),
\qquad
w_p=\frac{p^{-\sigma}}{(1-p^{-\sigma})^2}>0.
\label{eq:strict_covariance_determinant}
\end{equation}
\end{proof}

\section{The positivity route and its exact obstruction}
\label{sec:positivity_barrier}

The coefficients \(c(n)=-z_0''(n)\) are nonnegative.  It is therefore
natural to ask whether positivity itself forces the absence of the logarithmic
branches in \(1/2<\Re s<1\).  It does not.

\begin{proposition}[Complete monotonicity on the real axis]
\label{prop:K_complete_monotone}
For every real \(\sigma>1\) and every integer \(r\ge0\),
\begin{equation}
\boxed{
(-1)^r\mathcal K^{(r)}(\sigma)
=\sum_{n\ge1}\frac{c(n)(\log n)^r}{n^\sigma}\ge0.}
\label{eq:K_complete_monotonicity}
\end{equation}
In particular \(\mathcal K\) is completely monotone on \((1,\infty)\).
\end{proposition}

\begin{proof}
Differentiate the absolutely convergent Dirichlet series termwise.
\end{proof}

\begin{proposition}[Positivity does not remove complex logarithmic branches]
\label{prop:positivity_not_enough}
Nonnegative Dirichlet coefficients, complete monotonicity on the real axis,
and the strict covariance inequality
\eqref{eq:K_covariance_determinant} do not imply a branch-free continuation
of \(\mathcal K\) to \(1/2<\Re s<1\).
\end{proposition}

\begin{proof}
The source term \(P(s)=\sum_p p^{-s}\) itself has nonnegative Dirichlet
coefficients for \(\Re s>1\), yet its continuation is expressed through
\(\log\zeta(rs)\) and has logarithmic branch points at scaled zeros of
\(\zeta\); see \eqref{eq:prime_zeta_global_branches}.  More directly, the
positive-coefficient function \(S_0\) in
\eqref{eq:S0_branch_expansion} is exactly the factor that contributes
\(m\log(s-\rho)\) at a hypothetical zero \(\rho\) with
\(1/2<\Re\rho<1\).  Thus real-axis positivity and complete monotonicity do not
control the monodromy of the analytic continuation.
\end{proof}

\begin{corollary}[What would still be needed to prove RH]
\label{cor:needed_for_RH}
To deduce Corollary~\ref{cor:curvature_RH_criterion} from the curvature side,
one needs an additional property that controls the complex continuation of
\(\mathcal K\), not merely the signs of its Dirichlet coefficients or the
positive-semidefiniteness of the real covariance matrix in
\eqref{eq:K_covariance_determinant}.
\end{corollary}

\begin{proof}
Proposition~\ref{prop:K_complete_monotone} and
Theorem~\ref{thm:zeta_covariance_identity} control only real-axis signs
and positive semidefiniteness.  Proposition~\ref{prop:positivity_not_enough}
exhibits logarithmic branches compatible with those same properties.
Therefore any proof using this route must add information that controls the
complex continuation of \(\mathcal K\), as asserted.
\end{proof}

\begin{remark}[Status of the attempted step]
\label{rem:status_attempted_step}
The expected coefficient-positivity step has therefore been carried out and
closed: \(c(n)\ge0\) gives complete monotonicity and a positive covariance
Hessian, but those properties do not forbid the logarithmic branches whose
absence is equivalent to RH.  The obstruction is visible inside the same
formula: \(S_0\) has positive coefficients and is simultaneously the carrier
of the prime-zeta logarithmic monodromy.
\end{remark}


\section{A curvature--Weil positive kernel after primitive-replica filtering}
\label{sec:curvature_weil_kernel}

This section records the Version 9 positivity construction.  The underlying
positivity criterion is Weil's classical criterion for RH
\cite{Weil1952,Bombieri2000}.  What is specific to the present notebook is the
way the cocyclic curvature supplies a canonical positive Gram factor and the
way the transverse variable in \(B(s,t)\) removes prime-zeta replicas before
the zero kernel is formed.

\subsection{Polarization of the arithmetic curvature}
\label{subsec:curvature_polarization}

For an integer \(n\ge1\), write
\begin{equation}
w_p(n):=T(v_p(n))=\frac{v_p(n)(v_p(n)+1)}{2}.
\label{eq:curvature_weight_wp}
\end{equation}
For primes \(p<q\), define the real feature
\begin{equation}
\Phi_{pq}(n)
:=\frac{1}{\sqrt2}\sqrt{w_p(n)w_q(n)}\,\log\frac{q}{p}.
\label{eq:curvature_feature}
\end{equation}
Only finitely many of these features are nonzero for each fixed \(n\).

\begin{proposition}[Exact Gram polarization of the curvature]
\label{prop:curvature_gram_polarization}
For every integer \(n\ge1\),
\begin{equation}
\boxed{
\mathcal C(n)=\sum_{p<q}|\Phi_{pq}(n)|^2.}
\label{eq:curvature_feature_norm}
\end{equation}
Consequently
\begin{equation}
K_{\mathcal C}(m,n)
:=\sum_{p<q}\Phi_{pq}(m)\Phi_{pq}(n)
\label{eq:curvature_gram_kernel}
\end{equation}
is a well-defined positive-semidefinite kernel on \(\mathbb N\), and
\begin{equation}
K_{\mathcal C}(n,n)=\mathcal C(n).
\label{eq:curvature_kernel_diagonal}
\end{equation}
\end{proposition}

\begin{proof}
The closed curvature formula proved in
Theorem~\ref{thm:exact_curvature_formula} is
\begin{equation}
\mathcal C(n)
=\frac12\sum_{p<q}
T(v_p(n))T(v_q(n))\log^2\frac{q}{p}.
\label{eq:curvature_closed_recalled_for_gram}
\end{equation}
Substitution of \eqref{eq:curvature_weight_wp} and
\eqref{eq:curvature_feature} into the right-hand side of
\eqref{eq:curvature_feature_norm} gives exactly
\eqref{eq:curvature_closed_recalled_for_gram}.  For any finitely supported
complex family \((a_n)\), every sum below is finite and therefore may be
rearranged without a convergence issue:
\begin{align}
\sum_{m,n}\overline{a_m}a_nK_{\mathcal C}(m,n)
&=\sum_{p<q}\sum_{m,n}
\overline{a_m}a_n\Phi_{pq}(m)\Phi_{pq}(n)
\label{eq:curvature_kernel_quadratic_expand}\\
&=\sum_{p<q}\left|\sum_n a_n\Phi_{pq}(n)\right|^2
\ge0.
\label{eq:curvature_kernel_quadratic_positive}
\end{align}
Thus \(K_{\mathcal C}\) is positive semidefinite.  Setting \(m=n\) in
\eqref{eq:curvature_gram_kernel} and using
\eqref{eq:curvature_feature_norm} proves
\eqref{eq:curvature_kernel_diagonal}.
\end{proof}

\begin{remark}[Why the raw squared cocycle is not the kernel]
\label{rem:raw_cocycle_square_indefinite}
The seemingly direct kernel
\(\mathscr A(m,n)^2\), where
\(\mathscr A(m,n)=\Omega(m)\log n-\log m\,\Omega(n)\), is not positive
semidefinite in general.  If \(\mathscr A(m,n)\ne0\), its restriction to the
two indices \(m,n\) has matrix
\begin{equation}
\begin{pmatrix}
0&\mathscr A(m,n)^2\\
\mathscr A(m,n)^2&0
\end{pmatrix},
\label{eq:raw_cocycle_square_matrix}
\end{equation}
whose eigenvalues are
\(\pm\mathscr A(m,n)^2\).  The polarization in
Proposition~\ref{prop:curvature_gram_polarization} is therefore essential.
\end{remark}

\subsection{The exact primitive-replica filter}
\label{subsec:primitive_replica_filter}

Recall from Theorem~\ref{thm:O_initial_domain} that, for \(\Re s>1\),
\begin{equation}
B(s,t)=\sum_p\frac{p^{it}}{p^s-1}.
\label{eq:B_prime_series_recalled_kernel}
\end{equation}

\begin{proposition}[Exact transverse primitivization]
\label{prop:transverse_primitivization_kernel}
For \(\Re s>1\), differentiation in \(t\) is legitimate term by term and
\begin{equation}
\boxed{
\Pi B(s):=
 i\left.\frac{\partial B}{\partial t}(s,t)\right|_{t=0}
=\frac{\zeta'}{\zeta}(s).}
\label{eq:Pi_B_log_derivative}
\end{equation}
Moreover, if the grouped log-zeta representation is decomposed according to
the dilation index \(N=mk\), then \(\Pi\) annihilates every sector
\(N>1\) and retains the primitive sector \(N=1\).
\end{proposition}

\begin{proof}
For any fixed \(\sigma_0>1\) and \(|t|\le T\),
\begin{equation}
\left|\frac{\partial}{\partial t}
\frac{p^{it}}{p^s-1}\right|
\le \frac{\log p}{p^{\sigma_0}-1}.
\label{eq:B_t_derivative_majorant}
\end{equation}
The series on the right converges because
\begin{equation}
\sum_p\frac{\log p}{p^{\sigma_0}-1}
\le\sum_{p}\sum_{a\ge1}(\log p)p^{-a\sigma_0}
=-\frac{\zeta'}{\zeta}(\sigma_0)<\infty.
\label{eq:B_t_derivative_majorant_sum}
\end{equation}
The Weierstrass theorem therefore permits termwise differentiation of
\eqref{eq:B_prime_series_recalled_kernel}, uniformly on compact subsets of
\(\Re s>1\) and bounded \(t\)-intervals.  At \(t=0\),
\begin{align}
\left.\partial_t B(s,t)\right|_{t=0}
&=i\sum_p\frac{\log p}{p^s-1}
\label{eq:B_t_derivative_prime_sum}\\
&=-i\frac{\zeta'}{\zeta}(s),
\label{eq:B_t_derivative_logzeta_kernel}
\end{align}
which proves \eqref{eq:Pi_B_log_derivative}.

For the sector statement, use the absolutely convergent representation in
\(\Re s>1\)
\begin{equation}
B(s,t)=\sum_{k,m\ge1}\frac{\mu(m)}m
\log\zeta(mks-imt).
\label{eq:B_grouped_for_Pi}
\end{equation}
Differentiating as above and grouping by \(N=mk\) gives
\begin{align}
\left.\partial_tB(s,t)\right|_{t=0}
&=-i\sum_{N\ge1}\frac{\zeta'}{\zeta}(Ns)
\sum_{m\mid N}\mu(m).
\label{eq:Pi_grouped_mobius_sum}
\end{align}
All rearrangements are justified by absolute convergence in \(\Re s>1\).
The elementary identity
\begin{equation}
\sum_{m\mid N}\mu(m)=
\begin{cases}
1,&N=1,\\
0,&N>1
\end{cases}
\label{eq:mobius_primitive_projector}
\end{equation}
proves the asserted annihilation of every replica sector.
\end{proof}

\begin{remark}[Completion]
\label{rem:Pi_completion_xi}
The completed logarithmic derivative is recovered from
\eqref{eq:Pi_B_log_derivative} by the standard identity
\begin{equation}
\frac{\xi'}{\xi}(s)
=\Pi B(s)
+\frac1s+\frac1{s-1}
-\frac12\log\pi
+\frac12\psi\!\left(\frac{s}{2}\right),
\label{eq:xi_log_derivative_from_Pi}
\end{equation}
where \(\psi=\Gamma'/\Gamma\).  Thus the nontrivial poles of the completed
primitive object occur at the genuine zeros \(\rho\), not at the prime-zeta
replicas \(\rho/r\) with \(r>1\).
\end{remark}

\subsection{A zero kernel and Weil positivity}
\label{subsec:zero_kernel_weil}

Let \(g\in C_c^\infty((0,\infty))\), and use the Mellin transform
\begin{equation}
G(s):=\int_0^\infty g(x)x^{s-1}\,dx.
\label{eq:Weil_Mellin_transform}
\end{equation}
Define
\begin{equation}
G^\#(s):=\overline{G(1-\overline{s})},
\qquad
W_g(s):=G(s)G^\#(s).
\label{eq:Weil_involution_weight}
\end{equation}
For every fixed vertical strip, \(G(s)\) decreases faster than every inverse
power of \(|\Im s|\); this follows by repeated integration by parts after the
change of variable \(x=e^u\).  Consequently the zero sums below converge
absolutely, because the classical zero-counting estimate is
\(N(T+1)-N(T)=O(\log(2+T))\); see \cite{Titchmarsh1986}.

\begin{definition}[Primitive-zero kernel]
\label{def:primitive_zero_kernel}
For positive integers \(m,n\), define
\begin{equation}
K_g^Z(m,n)
:=\sum_{\rho}W_g(\rho)
\left(\frac{m}{n}\right)^{\rho-1/2},
\label{eq:primitive_zero_kernel}
\end{equation}
where every nontrivial zero is counted with multiplicity and the positive real
logarithm is used in the power.  The absolute convergence just noted makes
\eqref{eq:primitive_zero_kernel} independent of the order of summation.
\end{definition}

\begin{lemma}[Quadratic form of the zero kernel]
\label{lem:zero_kernel_quadratic_form}
For every finitely supported complex family \((a_n)\), set
\begin{equation}
A_a(s):=\sum_n a_n n^{1/2-s},
\qquad
A_a^\#(s):=\overline{A_a(1-\overline{s})}.
\label{eq:Aa_and_involution}
\end{equation}
Then
\begin{equation}
\boxed{
\sum_{m,n}\overline{a_m}a_nK_g^Z(m,n)
=\sum_\rho W_g(\rho)A_a(\rho)A_a^\#(\rho).}
\label{eq:zero_kernel_quadratic_identity}
\end{equation}
\end{lemma}

\begin{proof}
Because \((a_n)\) has finite support and
\eqref{eq:primitive_zero_kernel} is absolutely convergent, the finite sums in
\(m,n\) may be interchanged with the zero sum.  For a fixed zero \(\rho\),
\begin{align}
\sum_{m,n}\overline{a_m}a_n
\left(\frac{m}{n}\right)^{\rho-1/2}
&=\left(\sum_m\overline{a_m}m^{\rho-1/2}\right)
\left(\sum_n a_n n^{1/2-\rho}\right)
\label{eq:zero_kernel_factorization_1}\\
&=A_a(\rho)A_a^\#(\rho),
\label{eq:zero_kernel_factorization_2}
\end{align}
where the final identity is exactly the definition in
\eqref{eq:Aa_and_involution}.  Substitution proves
\eqref{eq:zero_kernel_quadratic_identity}.
\end{proof}

\begin{theorem}[RH implies positivity of the primitive-zero kernel]
\label{thm:RH_implies_zero_kernel_positive}
If RH holds, then for every \(g\in C_c^\infty((0,\infty))\), the kernel
\(K_g^Z\) is positive semidefinite.
\end{theorem}

\begin{proof}
Under RH, every nontrivial zero has the form
\(\rho=1/2+i\gamma\), whence
\begin{equation}
G^\#(\rho)=\overline{G(\rho)},
\qquad
W_g(\rho)=|G(\rho)|^2.
\label{eq:Weil_weight_under_RH}
\end{equation}
Likewise
\begin{equation}
A_a^\#(\rho)=\overline{A_a(\rho)}.
\label{eq:Aa_under_RH}
\end{equation}
Lemma~\ref{lem:zero_kernel_quadratic_form} therefore gives
\begin{equation}
\sum_{m,n}\overline{a_m}a_nK_g^Z(m,n)
=\sum_\rho |G(\rho)|^2|A_a(\rho)|^2\ge0.
\label{eq:zero_kernel_positive_under_RH}
\end{equation}
The zero sum converges absolutely by the decay stated before
Definition~\ref{def:primitive_zero_kernel}.
\end{proof}

\begin{theorem}[Weil equivalence for the primitive-zero kernel]
\label{thm:zero_kernel_Weil_equivalence}
The following are equivalent:
\begin{enumerate}[label=(\roman*)]
\item RH holds.
\item For every \(g\in C_c^\infty((0,\infty))\), the kernel \(K_g^Z\) is
positive semidefinite on finitely supported sequences.
\end{enumerate}
\end{theorem}

\begin{proof}
The implication (i)\(\Rightarrow\)(ii) is
Theorem~\ref{thm:RH_implies_zero_kernel_positive}.  Conversely, assume (ii).
Choose a single index \(N\) and a sequence supported only at \(N\), with
\(a_N=1\).  Then \eqref{eq:zero_kernel_quadratic_identity} gives
\begin{equation}
K_g^Z(N,N)=\sum_\rho W_g(\rho)\ge0
\label{eq:zero_kernel_diagonal_Weil}
\end{equation}
for every \(g\in C_c^\infty((0,\infty))\).  Bombieri's formulation of Weil's
criterion states that RH is equivalent to
\begin{equation}
\sum_\rho G(\rho)\overline{G(1-\overline\rho)}\ge0
\quad\text{for every }g\in C_c^\infty((0,\infty)),
\label{eq:Weil_criterion_used}
\end{equation}
with zeros counted with multiplicity
\cite{Weil1952,Bombieri2000}.  The summand in
\eqref{eq:Weil_criterion_used} is precisely \(W_g(\rho)\) by
\eqref{eq:Weil_involution_weight}.  Thus
\eqref{eq:zero_kernel_diagonal_Weil} for every \(g\) implies RH.
\end{proof}

\subsection{The curvature--zero Schur kernel}
\label{subsec:curvature_zero_schur}

\begin{definition}[Curvature--zero kernel]
\label{def:curvature_zero_kernel}
For \(g\in C_c^\infty((0,\infty))\), define
\begin{equation}
\mathscr K_g(m,n)
:=K_{\mathcal C}(m,n)K_g^Z(m,n).
\label{eq:curvature_zero_kernel}
\end{equation}
\end{definition}

\begin{theorem}[Exact positivity criterion for the combined kernel]
\label{thm:combined_kernel_RH_equivalence}
The following are equivalent:
\begin{enumerate}[label=(\roman*)]
\item RH holds.
\item For every \(g\in C_c^\infty((0,\infty))\), the kernel
\(\mathscr K_g\) is positive semidefinite on finitely supported sequences.
\end{enumerate}
\end{theorem}

\begin{proof}
Assume RH.  Proposition~\ref{prop:curvature_gram_polarization} and
Theorem~\ref{thm:RH_implies_zero_kernel_positive} show that both factors in
\eqref{eq:curvature_zero_kernel} are positive semidefinite.  A direct proof
of positivity of their entrywise product avoids invoking the Schur product
theorem.  Under RH, insert the Gram expansions
\eqref{eq:curvature_gram_kernel} and
\eqref{eq:zero_kernel_positive_under_RH}.  For every finitely supported
\((a_n)\), absolute convergence of the zero sum and finiteness of the support
permit rearrangement, giving
\begin{align}
\sum_{m,n}\overline{a_m}a_n\mathscr K_g(m,n)
&=\sum_{\rho}|G(\rho)|^2
\sum_{p<q}
\left|\sum_n a_n\Phi_{pq}(n)n^{1/2-\rho}\right|^2
\ge0.
\label{eq:combined_kernel_channel_sum}
\end{align}
Thus (i) implies (ii).

Conversely, assume (ii).  Choose any integer \(N\) with at least two distinct
prime factors; then Theorem~\ref{thm:exact_curvature_formula} gives
\(\mathcal C(N)>0\).  Test the quadratic form on the one-point sequence
\(a_N=1\).  By \eqref{eq:curvature_kernel_diagonal},
\begin{equation}
\mathscr K_g(N,N)
=\mathcal C(N)\sum_\rho W_g(\rho)\ge0.
\label{eq:combined_kernel_diagonal_Weil}
\end{equation}
Division by the positive number \(\mathcal C(N)\) gives Weil positivity for
every \(g\).  Theorem~\ref{thm:zero_kernel_Weil_equivalence} then yields RH.
\end{proof}

\begin{remark}[Logical status: the curvature is still a spectator in the converse]
\label{rem:curvature_spectator}
Theorem~\ref{thm:combined_kernel_RH_equivalence} is exact, but it is not a new
proof of RH.  In the converse direction, a single diagonal entry with
\(\mathcal C(N)>0\) already recovers Weil's classical quadratic functional.
Thus the present Schur product attaches the cocyclic curvature to Weil
positivity without yet making the off-diagonal curvature geometry logically
indispensable.
\end{remark}

\begin{question}[The genuinely new positivity target]
\label{q:unconditional_CCR_gram_realization}
Does there exist, for every admissible \(g\), a Hilbert space
\(\mathcal H_g\) and vectors \(\Psi_n^{(g)}\), constructed directly from the
Weyl operators, the arithmetic cocycle, twisted M"obius inversion, and the
transverse primitive filter \(\Pi\), but not from the zero set of \(\zeta\),
such that
\begin{equation}
\boxed{
\mathscr K_g(m,n)
=\langle\Psi_m^{(g)},\Psi_n^{(g)}\rangle_{\mathcal H_g}?}
\label{eq:unconditional_CCR_gram_target}
\end{equation}
An unconditional construction of \eqref{eq:unconditional_CCR_gram_target}
for the full Weil test class would imply RH by
Theorem~\ref{thm:combined_kernel_RH_equivalence}.  No such construction is
proved here.
\end{question}

\begin{remark}[The exact bridge and the remaining gap]
\label{rem:CCR_curvature_Weil_bridge}
The rigorous chain established in this notebook is
\begin{equation}
\boxed{
\text{CCR/cocycle}
\longrightarrow \mathcal C
\longrightarrow K_{\mathcal C}\succeq0,
\qquad
B\xrightarrow{\,\Pi\,}\frac{\zeta'}{\zeta}
\longrightarrow K_g^Z.}
\label{eq:two_branch_bridge}
\end{equation}
The two branches meet in \(\mathscr K_g=K_{\mathcal C}\circ K_g^Z\).
Its positivity for the full Weil class is equivalent to RH, but its known
proof uses RH on the zero side.  The unresolved step is to obtain the same
positivity from the CCR/cocycle side before zeros are used.
\end{remark}

\section{The Heisenberg-deformed zeta series}
\label{sec:deformed_zeta_series}

\begin{definition}[Heisenberg-deformed zeta series]
\label{def:deformed_zeta_series}
For $\Re s>\sigma_*$, define
\begin{equation}
\zeta_{\mathrm H}(s;\vartheta)
=
\sum_{n\ge1}\frac{z_{\vartheta}(n)}{n^s}.
\label{eq:deformed_zeta_series}
\end{equation}
Theorem~\ref{thm:neumann_region} guarantees absolute convergence for every $\Re s>\sigma_*$.
\end{definition}

\begin{proposition}[Classical boundary]
\label{prop:zeta_classical_boundary}
For $\Re s>1$,
\begin{equation}
\zeta_{\mathrm H}(s;0)=\zeta(s).
\label{eq:zeta_classical_boundary}
\end{equation}
\end{proposition}

\begin{proof}
By Proposition~\ref{prop:commuting_boundary}, $z_0(n)=1$ for all $n$.
Substitution into \eqref{eq:deformed_zeta_series} gives the ordinary
Dirichlet series for $\zeta(s)$.
\end{proof}

\begin{proposition}[Even Heisenberg deformation]
\label{prop:zeta_theta_even}
For $\Re s>\sigma_*$,
\begin{equation}
\zeta_{\mathrm H}(s;-\vartheta)
=
\zeta_{\mathrm H}(s;\vartheta).
\label{eq:zeta_theta_even}
\end{equation}
\end{proposition}

\begin{proof}
This follows termwise from Proposition~\ref{prop:theta_evenness} and absolute
convergence.
\end{proof}

\begin{warning}[No analytic continuation is asserted]
\label{warn:no_continuation}
The series \eqref{eq:deformed_zeta_series} is absolutely convergent and holomorphic on the open half-plane $\Re s>\sigma_*$.  No meromorphic continuation, functional equation, scalar Euler product, or critical-strip theory
for $\zeta_{\mathrm H}(s;\vartheta)$ is asserted here.
\end{warning}

\begin{proposition}[Conditional first-order rigidity of simple zeros]
\label{prop:conditional_zero_rigidity}
Assume that, for some simple zero $\rho$ of the ordinary zeta function,
there exists a function $F(s,\vartheta)$ analytic in a neighbourhood of
$(\rho,0)$ such that
\begin{align}
F(s,0)
&=\zeta(s),
\label{eq:F_classical_boundary}\\
F(s,-\vartheta)
&=F(s,\vartheta),
\label{eq:F_evenness}
\end{align}
and that $F$ agrees with $\zeta_{\mathrm H}$ wherever both are defined.
Then the local zero branch $\rho(\vartheta)$ through $\rho$ satisfies
\begin{equation}
\rho'(0)=0.
\label{eq:zero_first_derivative_zero}
\end{equation}
\end{proposition}

\begin{proof}
The implicit-function theorem applies because $\zeta'(\rho)\neq0$.
Evenness \eqref{eq:F_evenness} makes the unique local zero branch even, so
its first derivative at the origin vanishes.
\end{proof}

\begin{warning}
\label{warn:conditional_zero_result}
Proposition~\ref{prop:conditional_zero_rigidity} is conditional on an analytic
continuation in both variables and does not supply such a continuation.  It
is included because it identifies a precise way in which the exact
$\vartheta\mapsto-\vartheta$ symmetry would constrain the motion of classical
Riemann zeros under the Heisenberg deformation if that continuation exists.
\end{warning}

\section{A genuine M\"obius Hamiltonian}
\label{sec:dynamics}

\subsection{Definition and domain}
\label{subsec:hamiltonian_definition}

\begin{definition}[Cosine and sine Weyl quadratures]
\label{def:Weyl_quadratures}
For $n\ge1$, set
\begin{align}
C_n
&=\frac{W_n+W_n^*}{2},
\label{eq:Cn_definition}\\
J_n
&=\frac{W_n-W_n^*}{2i}.
\label{eq:Jn_definition}
\end{align}
Both are bounded self-adjoint operators with norm at most $1$.
\end{definition}

\begin{definition}[M\"obius Hamiltonian]
\label{def:mobius_hamiltonian}
Fix an energy scale $E_0>0$ and $\sigma>1$.  Define
\begin{equation}
H_{\mu,\sigma,\vartheta}
=
E_0\left(\Reop\mathcal M_{\sigma,\vartheta}-I\right)
=
E_0\sum_{n\ge2}
\frac{\mu(n)}{n^{\sigma}}C_n.
\label{eq:mobius_hamiltonian}
\end{equation}
The series converges absolutely in operator norm.
\end{definition}

\begin{proposition}[Self-adjoint bounded dynamics]
\label{prop:hamiltonian_selfadjoint}
For $\sigma>1$, $H_{\mu,\sigma,\vartheta}$ is bounded and self-adjoint on
all of $\Hil$:
\begin{equation}
\Dom(H_{\mu,\sigma,\vartheta})=\Hil.
\label{eq:hamiltonian_domain}
\end{equation}
It therefore generates the unitary group
\begin{equation}
U(t)=\exp\!\left(-\frac{i}{\hbar}tH_{\mu,\sigma,\vartheta}\right).
\label{eq:time_evolution_unitary}
\end{equation}
\end{proposition}

\begin{proof}
The norm series in \eqref{eq:mobius_hamiltonian} is dominated by
$E_0\sum_{n\ge2}|\mu(n)|n^{-\sigma}$, which converges for $\sigma>1$.
Every $C_n$ is self-adjoint, hence so is the norm limit.  The spectral theorem
then defines \eqref{eq:time_evolution_unitary}.
\end{proof}

\begin{proposition}[Explicit Schr\"odinger equation]
\label{prop:schrodinger_equation}
For $\psi_0\in\Hil$, the evolution $\psi(t)=U(t)\psi_0$ satisfies
\begin{equation}
i\hbar\,\partial_t\psi(t,q)
=
(H_{\mu,\sigma,\vartheta}\psi(t))(q)
\label{eq:schrodinger_equation}
\end{equation}
in the strong $L^2$ sense, where
\begin{align}
(H_{\mu,\sigma,\vartheta}\psi)(q)
=
\frac{E_0}{2}
\sum_{n\ge2}
\frac{\mu(n)}{n^{\sigma}}
\Bigg[
& e^{\frac{i}{\hbar}y_n(q-x_n/2)}\psi(q-x_n)
\notag\\
+{}&e^{-\frac{i}{\hbar}y_n(q+x_n/2)}\psi(q+x_n)
\Bigg].
\label{eq:hamiltonian_coordinate_action}
\end{align}
\end{proposition}

\begin{proof}
Equation~\eqref{eq:hamiltonian_coordinate_action} follows from
Proposition~\ref{prop:weyl_coordinate_action} applied to $W_n$ and $W_n^*$.
The series converges in operator norm, so the bounded Hamiltonian gives a
strong Schr\"odinger evolution for every initial vector.
\end{proof}

\subsection{Exact equations of motion for position and momentum}
\label{subsec:equations_of_motion}

\begin{lemma}[Bounded first commutators]
\label{lem:bounded_first_commutators}
For $\sigma>1$, the series
\begin{align}
\sum_{n\ge2}
\frac{|\mu(n)|x_n}{n^{\sigma}},
\label{eq:x_commutator_sum}\\
\sum_{n\ge2}
\frac{|\mu(n)|y_n}{n^{\sigma}}
\label{eq:y_commutator_sum}
\end{align}
converge.  Hence $[H_{\mu,\sigma,\vartheta},Q]$ and
$[H_{\mu,\sigma,\vartheta},P]$, initially defined on $\Sch(\R)$, extend to
bounded operators.
\end{lemma}

\begin{proof}
The first series is bounded by a constant multiple of
$\sum n^{-\sigma}\log n$.  Since
\begin{equation}
\Om(n)\le\frac{\log n}{\log2},
\label{eq:Omega_log_bound}
\end{equation}
the second is bounded by another constant multiple of the same convergent
series.
\end{proof}

\begin{theorem}[Exact Heisenberg equations of motion]
\label{thm:equations_of_motion}
Let
\begin{align}
Q(t)
&=U(t)^*QU(t),
\label{eq:Qt_definition}\\
P(t)
&=U(t)^*PU(t).
\label{eq:Pt_definition}
\end{align}
The bounded commutators in Lemma~\ref{lem:bounded_first_commutators} imply invariance of $\Dom(Q)$ and $\Dom(P)$ under $U(t)$.  Indeed, for $A=Q$ or $P$ and $\psi\in\Dom(A)$, Duhamel's formula first on a core and then by graph-norm closure gives
\begin{equation}
A U(t)\psi=U(t)A\psi+\frac{i}{\hbar}\int_0^t U(t-\tau)[H,A]U(\tau)\psi\,d\tau.
\label{eq:duhamel_domain_invariance}
\end{equation}
whose right-hand side lies in $L^2$ because $[H,A]$ is bounded.  On these domains, and equivalently as quadratic-form identities, the Heisenberg equations are
\begin{align}
\dot Q(t)
&=
\frac{E_0}{\hbar}
U(t)^*
\left[
\sum_{n\ge2}
\frac{\mu(n)x_n}{n^{\sigma}}J_n
\right]
U(t),
\label{eq:Q_equation_motion}\\
\dot P(t)
&=
\frac{E_0}{\hbar}
U(t)^*
\left[
\sum_{n\ge2}
\frac{\mu(n)y_n}{n^{\sigma}}J_n
\right]
U(t).
\label{eq:P_equation_motion}
\end{align}
At $t=0$ this becomes
\begin{align}
\dot Q(0)
&=
\frac{E_0\ell}{\hbar}
\sum_{n\ge2}
\frac{\mu(n)\log n}{n^{\sigma}}J_n,
\label{eq:Qdot_zero}\\
\dot P(0)
&=
\frac{E_0\kappa}{\hbar}
\sum_{n\ge2}
\frac{\mu(n)\Om(n)}{n^{\sigma}}J_n.
\label{eq:Pdot_zero}
\end{align}
\end{theorem}

\begin{proof}
For $A=Q$ or $P$, the bounded-commutator theorem for unitary dynamics gives
\begin{equation}
\dot A(t)
=
\frac{i}{\hbar}U(t)^*[H_{\mu,\sigma,\vartheta},A]U(t).
\label{eq:general_heisenberg_equation}
\end{equation}
From Proposition~\ref{prop:weyl_QP_commutators},
\begin{align}
[Q,C_n]
&=ix_nJ_n,
\label{eq:Q_Cn_commutator}\\
[P,C_n]
&=iy_nJ_n.
\label{eq:P_Cn_commutator}
\end{align}
Therefore
\begin{align}
[H_{\mu,\sigma,\vartheta},Q]
&=
-iE_0\sum_{n\ge2}
\frac{\mu(n)x_n}{n^{\sigma}}J_n,
\label{eq:H_Q_commutator}\\
[H_{\mu,\sigma,\vartheta},P]
&=
-iE_0\sum_{n\ge2}
\frac{\mu(n)y_n}{n^{\sigma}}J_n.
\label{eq:H_P_commutator}
\end{align}
Lemma~\ref{lem:bounded_first_commutators} justifies the sums as bounded
operators.  Substitution into \eqref{eq:general_heisenberg_equation} gives
\eqref{eq:Q_equation_motion} and \eqref{eq:P_equation_motion}.
\end{proof}

\begin{remark}[Physical and arithmetic meaning of the two equations]
\label{rem:meaning_equations_motion}
The position variable is $Q$, the logarithmic-scale coordinate.  The
momentum variable is $P=-i\hbar\partial_q$, the generator conjugate to that
coordinate.  Multiplication by an integer $n$ is represented by a phase-space
displacement of size $\ell\log n$ in $Q$ and $\kappa\Om(n)$ in $P$.
Equations~\eqref{eq:Qdot_zero} and \eqref{eq:Pdot_zero} therefore say that
the M\"obius Hamiltonian drives position by logarithmic size and momentum by
prime-factor count, with the same quantum current operators $J_n$.
\end{remark}

\subsection{The CCR is visible in the dynamics}
\label{subsec:CCR_in_dynamics}

\begin{proposition}[Pairwise dynamical noncommutativity]
\label{prop:CJ_commutator}
Extend $\log$ and $\Om$ additively to positive rationals and define $W_r$ by
\eqref{eq:Wn_definition} for $r\in\Qp$.  For $m,n\in\N$, put
\begin{equation}
\Phi_{m,n}
=
\vartheta\bigl(\Om(m)\log n-\log m\,\Om(n)\bigr).
\label{eq:Phi_mn}
\end{equation}
Then
\begin{align}
[C_m,J_n]
=
\frac12\sin\!\left(\frac{\Phi_{m,n}}{2}\right)
\Bigl(
W_{mn}+W_{m/n}-W_{n/m}-W_{1/(mn)}
\Bigr).
\label{eq:Cm_Jn_commutator}
\end{align}
\end{proposition}

\begin{proof}
Expand $C_m$ and $J_n$ using
\eqref{eq:Cn_definition}--\eqref{eq:Jn_definition}.  Apply
\eqref{eq:weyl_commutator} to the four pairs
$(z_m,z_n)$, $(z_m,-z_n)$, $(-z_m,z_n)$, and $(-z_m,-z_n)$.
\end{proof}

\begin{corollary}[Second-order arithmetic commutator]
\label{cor:second_order_force}
At $t=0$,
\begin{align}
\ddot Q(0)
=
\frac{iE_0^2}{2\hbar^2}
\sum_{m,n\ge2}
&\frac{\mu(m)\mu(n)x_n}{(mn)^{\sigma}}
\sin\!\left(\frac{\Phi_{m,n}}{2}\right)
\notag\\
&\times
\Bigl(
W_{mn}+W_{m/n}-W_{n/m}-W_{1/(mn)}
\Bigr),
\label{eq:Q_second_derivative}
\end{align}
and the analogous momentum equation is
\begin{align}
\ddot P(0)
=
\frac{iE_0^2}{2\hbar^2}
\sum_{m,n\ge2}
&\frac{\mu(m)\mu(n)y_n}{(mn)^{\sigma}}
\sin\!\left(\frac{\Phi_{m,n}}{2}\right)
\notag\\
&\times
\Bigl(
W_{mn}+W_{m/n}-W_{n/m}-W_{1/(mn)}
\Bigr).
\label{eq:P_second_derivative}
\end{align}
The double sums converge in operator norm for $\sigma>1$.
\end{corollary}

\begin{proof}
Differentiate \eqref{eq:Q_equation_motion} and
\eqref{eq:P_equation_motion} once more and use
Proposition~\ref{prop:CJ_commutator}.  Polynomial logarithmic weights are
summable against $(mn)^{-\sigma}$ for $\sigma>1$.
\end{proof}

\begin{corollary}[The interaction disappears with the Weyl cocycle]
\label{cor:interaction_disappears}
If $\vartheta=0$, then
\begin{equation}
\Phi_{m,n}=0
\label{eq:Phi_zero}
\end{equation}
for all $m,n$, and all pairwise terms in
\eqref{eq:Q_second_derivative}--\eqref{eq:P_second_derivative} vanish.
Thus the second-order cross-arithmetic second-order term displayed above is generated
entirely by the exact Heisenberg cocycle.
\end{corollary}

\begin{proof}
If \(\vartheta=0\), the phase \(\Phi_{m,n}\) in
Proposition~\ref{prop:CJ_commutator} is zero for every pair, so every summand
in the displayed second-order cross terms contains a factor
\(\sin(\Phi_{m,n}/2)=0\) or the equivalent commutator factor.  The pairwise
contributions therefore vanish identically.
\end{proof}

\section{Exact non-localization inside the arithmetic dynamics}
\label{sec:nonlocalization}

\begin{proposition}[Schwartz action and distributional transpose]
\label{prop:schwartz_distribution_extension}
For \(\sigma>1\), the Hamiltonian
\(H_{\mu,\sigma,\vartheta}\) maps \(\Sch(\R)\) continuously into
itself.  Its continuous transpose
\(H^{\times}:\Sch'(\R)\to\Sch'(\R)\) is defined by
\begin{equation}
\langle H^{\times}T,\phi\rangle
=
\langle T,H_{\mu,\sigma,\vartheta}\phi\rangle,
\qquad
T\in\Sch'(\R),\quad\phi\in\Sch(\R).
\label{eq:distributional_transpose_definition}
\end{equation}
Then \(H^{\times}\) is continuous on \(\Sch'(\R)\).
\end{proposition}

\begin{proof}
For every Schwartz seminorm
\(p_{a,b}(\phi)=\sup_q|q^a\phi^{(b)}(q)|\), the coordinate action of a
Weyl operator gives
\begin{equation}
p_{a,b}(W_n\phi)
\le
C_{a,b}(1+|x_n|+|y_n|)^{N_{a,b}}
\sum_{j\le b}p_{a+j,b-j}(\phi)
\label{eq:weyl_schwartz_bound}
\end{equation}
for suitable constants.  Since \(x_n,y_n=O(\log n)\) and
\begin{equation}
\sum_{n\ge2}n^{-\sigma}(1+\log n)^N<\infty
\qquad(\sigma>1)
\label{eq:log_weight_summability}
\end{equation}
for every fixed \(N\), the Hamiltonian series converges in every Schwartz
seminorm and defines a continuous map \(\Sch\to\Sch\).  Therefore
\eqref{eq:distributional_transpose_definition} defines the standard
continuous dual action on \(\Sch'\).  No termwise extension to all tempered
distributions is assumed without this transpose construction.
\end{proof}

\begin{proposition}[No exact position point is a stationary generalized mode]
\label{prop:no_delta_eigenstate}
Let \(\sigma>1\), \(\ell>0\), and \(E_0>0\).  For no
\(q_0\in\R\) and no scalar \(E\in\C\) does
\(\delta_{q_0}\in\Sch'(\R)\) satisfy
\begin{equation}
H^{\times}_{\mu,\sigma,\vartheta}\delta_{q_0}
=E\delta_{q_0}.
\label{eq:delta_eigenstate_impossible}
\end{equation}
\end{proposition}

\begin{proof}
The transpose action of a Weyl displacement on a point mass is
\begin{align}
W(x,y)^{\times}\delta_{q_0}
&=
\exp\left(\frac{i}{\hbar}y(q_0-x/2)\right)
\delta_{q_0-x},
\label{eq:W_delta_action}\\
(W(x,y)^*)^{\times}\delta_{q_0}
&=
\exp\left(-\frac{i}{\hbar}y(q_0+x/2)\right)
\delta_{q_0+x}.
\label{eq:Wstar_delta_action}
\end{align}
The \(n=2\) term therefore contributes nonzero masses at
\(q_0\pm\ell\log2\).  No other term produces the same support point as the unique \(n=2\) contribution at either \(q_0-\ell\log2\) or \(q_0+\ell\log2\).
Since \(\mu(2)=-1\), the corresponding coefficient cannot cancel.  Hence
the support of the left side of \eqref{eq:delta_eigenstate_impossible}
contains a point distinct from \(q_0\), while the right side is supported
only at \(q_0\).
\end{proof}

\begin{remark}[Relation to exact Heisenberg impossibility]
\label{rem:nonlocalization_relation}
Theorem~\ref{thm:no_joint_eigenstate} is a universal consequence of the CCR.
Proposition~\ref{prop:no_delta_eigenstate} is more specific: the M\"obius
Hamiltonian itself does not admit an exactly position-localized stationary
mode.  The reason is not a measurement uncertainty but the exact translation
content of the arithmetic Weyl operators.
\end{remark}


\section{Mertens record laws and Heisenberg jets}
\label{sec:mertens_records_heisenberg}

This section develops a Mertens-side consequence of classical oscillation
theory and then couples it exactly to an auxiliary Heisenberg lift.  The
analytic inputs are classical and are identified explicitly.  The deductions
from those inputs and the Weyl-area identities are proved below.  No claim is
made that the record law itself is a new theorem about the zero set.

\subsection{Historical maxima and the zero supremum}
\label{subsec:record_maximum_theta}

\begin{definition}[Historical maximum and record times]
\label{def:Astar_records}
For \(x\ge1\), define
\begin{equation}
A^*(x):=\max_{1\le n\le x}|M(n)|,
\qquad
M(x):=\sum_{n\le x}\mu(n).
\label{eq:Astar_definition}
\end{equation}
Let
\begin{equation}
\Theta:=\sup\{\Re\rho:\zeta(\rho)=0,\ 0<\Re\rho<1\}.
\label{eq:Theta_definition}
\end{equation}
Since the oscillation theorem below implies that \(A^*(x)\to\infty\), and since \(A^*\) is integer-valued and can increase by at most one at each integer because \(|M(n)-M(n-1)|\le1\), it cannot skip an integer height.  Thus for
each integer \(k\ge1\) we may define
\begin{equation}
r_k:=\min\{n\ge1:A^*(n)=k\}.
\label{eq:record_time_definition}
\end{equation}
\end{definition}

\begin{theorem}[Classical oscillation input]
\label{thm:KP_input}
For every \(\delta>0\) and \(\varepsilon>0\), all sufficiently large \(Y\)
have the following property: the interval
\([Y^{1-\delta},Y]\) contains a sign change of the Mertens function with
combined oscillation of size \(x^{\Theta-\varepsilon}\).  In particular,
there is a point \(y\in[Y^{1-\delta},Y]\) for which
\begin{equation}
|M(y)|\gg_{\delta,\varepsilon}y^{\Theta-\varepsilon}.
\label{eq:KP_large_value}
\end{equation}
\end{theorem}

\begin{proof}
This is the specialization \(a=0\) of Corollary~4 of
Kaczorowski--Pintz~\cite{KaczorowskiPintz1986}, combined with the
``combined oscillation'' clause in their Theorem~1.  Their interval is of the
form \([Y^{1-\delta},Y]\), and their singularity exponent for the Mertens
Dirichlet integral is precisely the supremum \(\Theta\) of the real parts of
nontrivial zeros of \(\zeta\).  The definition of combined oscillation gives
values on the two sides of the sign change whose absolute values have the
stated order; retaining either one gives \eqref{eq:KP_large_value}.
\end{proof}

\begin{lemma}[Standard upper bound at the zero supremum]
\label{lem:M_upper_theta}
For every \(\eta>0\),
\begin{equation}
M(x)\ll_\eta x^{\Theta+\eta},
\qquad
A^*(x)\ll_\eta x^{\Theta+\eta}.
\label{eq:M_upper_theta}
\end{equation}
\end{lemma}

\begin{proof}
The identity
\(\sum_{n\ge1}\mu(n)n^{-s}=1/\zeta(s)\) initially holds for
\(\Re s>1\).  The standard Littlewood--Perron argument continues the
corresponding summatory estimate throughout every fixed half-plane
\(\Re s>\Theta+\eta\), because \(1/\zeta(s)\) is holomorphic there and has
the standard vertical growth available away from the zero set.  This gives
the first estimate in \eqref{eq:M_upper_theta}; see
Titchmarsh~\cite{Titchmarsh1986}.  Taking the maximum over \(n\le x\) gives
the second estimate after enlarging the implied constant.
\end{proof}

\begin{theorem}[Uniform growth exponent of the historical maximum]
\label{thm:Astar_theta}
One has
\begin{equation}
\boxed{A^*(x)=x^{\Theta+o(1)}.}
\label{eq:Astar_theta}
\end{equation}
Equivalently,
\begin{equation}
\frac{\log A^*(x)}{\log x}\longrightarrow\Theta.
\label{eq:log_Astar_theta}
\end{equation}
\end{theorem}

\begin{proof}
Fix \(\eta>0\).  Choose \(\delta,\varepsilon>0\) so small that
\begin{equation}
(1-\delta)(\Theta-\varepsilon)>\Theta-\eta.
\label{eq:delta_epsilon_choice}
\end{equation}
By Theorem~\ref{thm:KP_input}, for every sufficiently large \(Y\) there is
\(y\in[Y^{1-\delta},Y]\) satisfying
\(|M(y)|\gg y^{\Theta-\varepsilon}\).  Hence
\begin{equation}
A^*(Y)
\ge |M(y)|
\gg Y^{(1-\delta)(\Theta-\varepsilon)}
\gg Y^{\Theta-\eta}.
\label{eq:Astar_lower_eta}
\end{equation}
Lemma~\ref{lem:M_upper_theta} gives
\begin{equation}
A^*(Y)\ll_\eta Y^{\Theta+\eta}.
\label{eq:Astar_upper_eta}
\end{equation}
Since \(\eta>0\) is arbitrary, \eqref{eq:Astar_theta} follows.
\end{proof}

\subsection{Uniform separation of logarithmic scales}
\label{subsec:scale_separation}

\begin{corollary}[Scale separation]
\label{cor:scale_separation}
For every fixed \(0<\alpha<1\),
\begin{equation}
\boxed{
\frac{A^*(x^\alpha)}{A^*(x)}
=x^{-(1-\alpha)\Theta+o(1)}\longrightarrow0.}
\label{eq:scale_separation}
\end{equation}
\end{corollary}

\begin{proof}
Apply Theorem~\ref{thm:Astar_theta} at \(x\) and \(x^\alpha\):
\begin{equation}
A^*(x^\alpha)=x^{\alpha\Theta+o(1)},
\qquad
A^*(x)=x^{\Theta+o(1)}.
\label{eq:Astar_two_scales}
\end{equation}
Their quotient is the expression in \eqref{eq:scale_separation}.  Since
\(\Theta\ge1/2\) and \(\alpha<1\), its main exponent is strictly negative.
\end{proof}

\begin{remark}[Relation to the earlier record criterion]
\label{rem:record_criterion_closed}
The sufficient condition isolated earlier in the notebook was only the
record-subsequence statement
\(A^*(r_k^\alpha)=o(A^*(r_k))\).  Corollary
\ref{cor:scale_separation} proves the stronger uniform statement for all
\(x\to\infty\).  Thus no separate argument excluding long historical
plateaus is needed.
\end{remark}

\subsection{Logarithmic historical collapse}
\label{subsec:log_history_collapse}

\begin{definition}[Logarithmic history and energy]
\label{def:I_E}
For real \(x\ge1\), with \(M(t)=M(\lfloor t\rfloor)\), put
\begin{align}
I(x)&:=\int_1^x\frac{M(t)}{t}\,dt,
\label{eq:mertens_I_definition}\\
E(x)&:=\int_1^x\frac{M(t)^2}{t}\,dt.
\label{eq:E_definition}
\end{align}
\end{definition}

\begin{theorem}[Uniform logarithmic \(L^p\)-collapse]
\label{thm:uniform_Lp_collapse}
For every fixed \(p>0\),
\begin{equation}
\boxed{
\frac1{A^*(x)^p\log x}
\int_1^x|M(t)|^p\frac{dt}{t}
\longrightarrow0.}
\label{eq:uniform_Lp_collapse}
\end{equation}
Consequently,
\begin{equation}
\boxed{I(x)=o(A^*(x)\log x).}
\label{eq:I_uniform_Astar}
\end{equation}
\end{theorem}

\begin{proof}
Fix \(0<\alpha<1\).  Split the integral at \(x^\alpha\).  By the definition
of \(A^*\),
\begin{align}
\int_1^x|M(t)|^p\frac{dt}{t}
&\le
A^*(x^\alpha)^p\int_1^{x^\alpha}\frac{dt}{t}
+A^*(x)^p\int_{x^\alpha}^{x}\frac{dt}{t}
\notag\\
&=
\alpha A^*(x^\alpha)^p\log x
+(1-\alpha)A^*(x)^p\log x.
\label{eq:Lp_cutoff_bound}
\end{align}
After division by \(A^*(x)^p\log x\), Corollary
\ref{cor:scale_separation} gives
\begin{equation}
\limsup_{x\to\infty}
\frac1{A^*(x)^p\log x}
\int_1^x|M(t)|^p\frac{dt}{t}
\le1-\alpha.
\label{eq:Lp_limsup_alpha}
\end{equation}
Letting \(\alpha\uparrow1\) proves \eqref{eq:uniform_Lp_collapse}.  The
case \(p=1\), together with
\(|I(x)|\le\int_1^x|M(t)|dt/t\), gives \eqref{eq:I_uniform_Astar}.
\end{proof}

\begin{corollary}[Record logarithmic law]
\label{cor:record_log_law}
At the record times \(r_k\),
\begin{equation}
\boxed{
\frac{I(r_k)}{M(r_k)\log r_k}\longrightarrow0.}
\label{eq:record_I_law}
\end{equation}
Moreover, for every \(p>0\),
\begin{equation}
\boxed{
\frac1{\log r_k}\int_1^{r_k}
\left|\frac{M(t)}{M(r_k)}\right|^p\frac{dt}{t}
\longrightarrow0.}
\label{eq:record_Lp_law}
\end{equation}
In particular,
\begin{equation}
\boxed{E(r_k)=o(k^2\log r_k).}
\label{eq:record_energy_small}
\end{equation}
\end{corollary}

\begin{proof}
By definition, \(A^*(r_k)=|M(r_k)|=k\).  Substitute \(x=r_k\) in
Theorem~\ref{thm:uniform_Lp_collapse}.  Equation
\eqref{eq:record_I_law} follows from \eqref{eq:I_uniform_Astar}, and
\eqref{eq:record_energy_small} is the case \(p=2\).
\end{proof}

\begin{corollary}[Logarithmic rarity of values comparable with the record]
\label{cor:record_occupancy}
For every fixed \(c>0\),
\begin{equation}
\boxed{
\frac1{\log r_k}
\int_1^{r_k}
\mathbf 1_{\{|M(t)|\ge ck\}}\frac{dt}{t}
\longrightarrow0.}
\label{eq:record_occupancy}
\end{equation}
\end{corollary}

\begin{proof}
For any fixed \(p>0\), Markov's inequality gives
\begin{equation}
\mathbf 1_{\{|M(t)|\ge ck\}}
\le c^{-p}\left|\frac{M(t)}{k}\right|^p.
\label{eq:occupancy_markov}
\end{equation}
Integrate with respect to \(dt/t\) and apply
\eqref{eq:record_Lp_law}.
\end{proof}

\subsection{The law of the record times}
\label{subsec:record_time_law}

\begin{theorem}[Record-time exponent]
\label{thm:record_time_exponent}
The first record times satisfy
\begin{equation}
\boxed{r_k=k^{1/\Theta+o(1)}.}
\label{eq:record_time_exponent}
\end{equation}
Consequently, for every fixed \(0<\varepsilon<1\),
\begin{equation}
\boxed{
\frac{\log r_{\lceil\varepsilon k\rceil}}
{\log r_k}\longrightarrow1.}
\label{eq:record_log_regular_variation}
\end{equation}
\end{theorem}

\begin{proof}
At \(x=r_k\), Theorem~\ref{thm:Astar_theta} gives
\begin{equation}
k=A^*(r_k)=r_k^{\Theta+o(1)}.
\label{eq:k_rk_theta}
\end{equation}
Taking logarithms yields
\(\log k=(\Theta+o(1))\log r_k\), hence
\eqref{eq:record_time_exponent}.  Applying the same formula to
\(\lceil\varepsilon k\rceil\) gives
\begin{equation}
\frac{\log r_{\lceil\varepsilon k\rceil}}{\log r_k}
=
\frac{(1/\Theta+o(1))\log(\lceil\varepsilon k\rceil)}
{(1/\Theta+o(1))\log k}
\longrightarrow1.
\label{eq:record_time_ratio_proof}
\end{equation}
\end{proof}

\subsection{An auxiliary Heisenberg lift of the Mertens path}
\label{subsec:aux_heisenberg_mertens}

The record theorem above is a consequence of classical analytic number theory.
We now connect it to the arithmetic cocycle by an exact auxiliary Weyl lift.
This lift is not substituted for the original embedding
\(n\mapsto(\ell\log n,\kappa\Omega(n))\); rather, its first transverse jet
recovers that original symplectic determinant.

\begin{definition}[Transverse Mertens deformation]
\label{def:Ft_At}
For real \(t\), define
\begin{align}
a_t(n)&:=\mu(n)e^{t\Omega(n)},
\label{eq:at_definition}\\
F_t(N)&:=\sum_{n\le N}a_t(n),
\label{eq:Ft_definition}\\
L(N)&:=\sum_{n\le N}\mu(n)\log n.
\label{eq:L_definition}
\end{align}
Let
\begin{equation}
v_n(t):=\bigl(a_t(n),\mu(n)\log n\bigr)\in\mathbb R^2
\label{eq:vn_auxiliary}
\end{equation}
and define the ordered area
\begin{equation}
\mathcal A_t(N)
:=
\frac12\sum_{m<n\le N}
\det(v_m(t),v_n(t)).
\label{eq:At_definition}
\end{equation}
Equivalently,
\begin{equation}
\mathcal A_t(N)
=
\frac12\sum_{m<n\le N}\mu(m)\mu(n)
\left(e^{t\Omega(m)}\log n-\log m\,e^{t\Omega(n)}\right).
\label{eq:At_double_sum}
\end{equation}
\end{definition}

\begin{remark}[Weyl interpretation]
\label{rem:At_weyl_interpretation}
For any fixed nonzero scale constant \(c\), the ordered product of Weyl
displacements with increments proportional to \(v_n(t)\) has central phase
\(\exp(ic\mathcal A_t(N))\).  Thus \(\mathcal A_t\) is literally the
central coordinate of the step-two Heisenberg lift of the auxiliary path.
\end{remark}

\begin{proposition}[Exact incremental equation]
\label{prop:At_increment}
For every integer \(N\ge2\),
\begin{equation}
\boxed{
\Delta\mathcal A_t(N)
=
\frac{\mu(N)}2
\left[
F_t(N-1)\log N-e^{t\Omega(N)}L(N-1)
\right],}
\label{eq:At_increment}
\end{equation}
where \(\Delta\mathcal A_t(N)=\mathcal A_t(N)-\mathcal A_t(N-1)\).
\end{proposition}

\begin{proof}
Only pairs whose second entry is \(N\) are added when the cutoff increases
from \(N-1\) to \(N\).  Hence
\begin{align}
2\Delta\mathcal A_t(N)
&=
\sum_{m<N}
\left[a_t(m)\mu(N)\log N-\mu(m)\log m\,a_t(N)\right]
\notag\\
&=
\mu(N)\log N\,F_t(N-1)
-\mu(N)e^{t\Omega(N)}L(N-1),
\label{eq:At_increment_proof}
\end{align}
which is \eqref{eq:At_increment}.
\end{proof}

\begin{proposition}[Order-zero area--energy identity]
\label{prop:A0_energy_identity}
For every integer \(N\ge1\),
\begin{equation}
\boxed{
2\mathcal A_0(N)=M(N)I(N)-E(N).}
\label{eq:A0_energy_identity}
\end{equation}
\end{proposition}

\begin{proof}
At \(t=0\), Proposition~\ref{prop:At_increment} gives
\begin{equation}
\Delta\mathcal A_0(N)
=
\frac{\mu(N)}2
\left[M(N-1)\log N-L(N-1)\right].
\label{eq:A0_increment_first}
\end{equation}
Abel summation gives the exact identity
\begin{equation}
L(N-1)=M(N-1)\log(N-1)-I(N-1),
\label{eq:L_Abel_Nminus1}
\end{equation}
and the step-function convention in Definition~\ref{def:I_E} gives
\begin{equation}
I(N)-I(N-1)=M(N-1)\log\frac{N}{N-1}.
\label{eq:I_increment}
\end{equation}
Therefore
\begin{equation}
\Delta\mathcal A_0(N)=\frac12\mu(N)I(N).
\label{eq:A0_increment_I}
\end{equation}
Summing over \(N\) yields
\begin{equation}
2\mathcal A_0(N)=\sum_{n\le N}\mu(n)I(n).
\label{eq:A0_sum_muI}
\end{equation}
Discrete summation by parts now gives
\begin{align}
\sum_{n\le N}\mu(n)I(n)
&=M(N)I(N)
-\sum_{n<N}M(n)\bigl(I(n+1)-I(n)\bigr)
\notag\\
&=M(N)I(N)
-\sum_{n<N}M(n)^2\log\frac{n+1}{n}
\notag\\
&=M(N)I(N)-E(N),
\label{eq:muI_summation_by_parts}
\end{align}
which proves \eqref{eq:A0_energy_identity}.
\end{proof}

\begin{corollary}[Central area at record times]
\label{cor:A0_record_small}
At the record times,
\begin{equation}
\boxed{\mathcal A_0(r_k)=o(k^2\log r_k).}
\label{eq:A0_record_small}
\end{equation}
\end{corollary}

\begin{proof}
By Corollary~\ref{cor:record_log_law},
\(M(r_k)I(r_k)=o(k^2\log r_k)\) and
\(E(r_k)=o(k^2\log r_k)\).  Apply
Proposition~\ref{prop:A0_energy_identity}.
\end{proof}

\subsection{The first transverse jet is the original arithmetic cocycle}
\label{subsec:first_jet_cocycle}

\begin{proposition}[Cocycle as first transverse jet]
\label{prop:first_jet_cocycle}
One has
\begin{equation}
\boxed{
\mathcal A'_0(N)
=
\frac12\sum_{m<n\le N}\mu(m)\mu(n)\mathscr A(m,n),}
\label{eq:first_jet_cocycle}
\end{equation}
where
\begin{equation}
\mathscr A(m,n)=\Omega(m)\log n-\log m\,\Omega(n).
\label{eq:first_jet_A_definition}
\end{equation}
Moreover,
\begin{equation}
\boxed{
\Delta\mathcal A'_0(N)
=
\frac{\mu(N)}2
\left[
M_\Omega(N-1)\log N-\Omega(N)L(N-1)
\right],}
\label{eq:first_jet_increment}
\end{equation}
where
\begin{equation}
M_\Omega(x):=\sum_{n\le x}\mu(n)\Omega(n).
\label{eq:MOmega_definition}
\end{equation}
\end{proposition}

\begin{proof}
Differentiate the finite sum \eqref{eq:At_double_sum} at \(t=0\).  This gives
\eqref{eq:first_jet_cocycle} immediately.  Differentiating the finite
increment identity \eqref{eq:At_increment} at \(t=0\) gives
\eqref{eq:first_jet_increment}, because
\(F'_0(N-1)=M_\Omega(N-1)\).
\end{proof}

\begin{lemma}[Selberg--Delange input for the weighted Mertens sum]
\label{lem:MOmega_selberg_delange}
As \(x\to\infty\),
\begin{equation}
\boxed{M_\Omega(x)\sim\frac{x}{(\log x)^2}.}
\label{eq:MOmega_asymptotic}
\end{equation}
\end{lemma}

\begin{proof}
For complex \(z\) in a fixed neighborhood of \(1\), the Dirichlet series is
\begin{equation}
D(s,z)
:=\sum_{n\ge1}\frac{\mu(n)z^{\Omega(n)}}{n^s}
=\prod_p(1-zp^{-s})
=\zeta(s)^{-z}G(s,z),
\label{eq:Dsz_selberg_delange}
\end{equation}
where
\begin{equation}
G(s,z):=
\prod_p(1-zp^{-s})(1-p^{-s})^{-z}
\label{eq:Gsz_definition}
\end{equation}
is holomorphic in the standard Selberg--Delange region near \((s,z)=(1,1)\),
and \(G(1,1)=1\).  The uniform Selberg--Delange theorem for compact
\(z\)-sets gives the corresponding asymptotic family; see
Tenenbaum~\cite{Tenenbaum2015}.  Its leading coefficient is
\begin{equation}
\frac{G(1,z)}{\Gamma(-z)}
\frac{x}{(\log x)^{z+1}}.
\label{eq:SD_leading_family}
\end{equation}
Since
\begin{equation}
\frac1{\Gamma(-z)}=(z-1)+O((z-1)^2)
\qquad(z\to1),
\label{eq:gamma_zero_at_one}
\end{equation}
and the Selberg--Delange expansion is locally uniform and holomorphic in
\(z\), differentiating at \(z=1\) yields
\begin{equation}
\frac{\partial}{\partial z}
\sum_{n\le x}\mu(n)z^{\Omega(n)}\Big|_{z=1}
=\frac{x}{(\log x)^2}(1+o(1)).
\label{eq:SD_derivative_at_one}
\end{equation}
The derivative on the left is precisely \(M_\Omega(x)\).
\end{proof}

\begin{theorem}[Local first-jet reconstruction of M\"obius]
\label{thm:first_jet_reconstruction}
As \(N\to\infty\),
\begin{equation}
\boxed{
\Delta\mathcal A'_0(N)
=
\frac{\mu(N)N}{2\log N}
+o\!\left(\frac{N}{\log N}\right).}
\label{eq:first_jet_reconstruction}
\end{equation}
If \(\mu(N)\ne0\), then equivalently
\begin{equation}
\Delta\mathcal A'_0(N)
=
\frac{\mu(N)}2\frac{N}{\log N}(1+o(1)).
\label{eq:first_jet_reconstruction_squarefree}
\end{equation}
If \(\mu(N)=0\), the left side of
\eqref{eq:first_jet_reconstruction} is exactly zero.
\end{theorem}

\begin{proof}
By Lemma~\ref{lem:MOmega_selberg_delange},
\begin{equation}
M_\Omega(N-1)\log N
=\frac{N}{\log N}(1+o(1)).
\label{eq:first_jet_main_term}
\end{equation}
The classical zero-free-region estimate for Mertens sums gives, with some
\(c>0\),
\begin{equation}
M(x)\ll
x\exp\left\{-c(\log x)^{3/5}(\log\log x)^{-1/5}\right\}.
\label{eq:classical_M_zero_free_bound}
\end{equation}
By Abel summation,
\begin{equation}
L(x)=M(x)\log x-I(x)
\ll
x\log x\exp\left\{-c'(\log x)^{3/5}(\log\log x)^{-1/5}\right\}
\label{eq:L_zero_free_bound}
\end{equation}
for a possibly smaller \(c'>0\).  Since
\(\Omega(N)\le\log N/\log2\),
\begin{equation}
\Omega(N)L(N-1)=o\!\left(\frac{N}{\log N}\right).
\label{eq:OmegaL_negligible}
\end{equation}
Insert \eqref{eq:first_jet_main_term} and
\eqref{eq:OmegaL_negligible} into
\eqref{eq:first_jet_increment}.  This proves
\eqref{eq:first_jet_reconstruction}.  If \(\mu(N)=0\), the exact identity
\eqref{eq:first_jet_increment} shows that the increment vanishes exactly.
\end{proof}

\begin{corollary}[First jet at record times]
\label{cor:first_jet_records}
Let
\(\sigma_k:=\operatorname{sgn}M(r_k)\).  Then
\begin{equation}
\mu(r_k)=\sigma_k
\label{eq:mu_at_record}
\end{equation}
and
\begin{equation}
\boxed{
\Delta\mathcal A'_0(r_k)
\sim
\frac{\sigma_k}{2}\frac{r_k}{\log r_k}.}
\label{eq:first_jet_record_asymptotic}
\end{equation}
Consequently,
\begin{equation}
\boxed{
\frac{\log|\Delta\mathcal A'_0(r_k)|}{\log k}
\longrightarrow\frac1\Theta.}
\label{eq:first_jet_record_theta}
\end{equation}
\end{corollary}

\begin{proof}
By minimality of \(r_k\), one has \(|M(r_k-1)|\le k-1\), whereas
\(|M(r_k)|=k\).  Since
\(M(r_k)-M(r_k-1)=\mu(r_k)\in\{-1,0,1\}\), the final increment must have the
same sign as \(M(r_k)\), proving \eqref{eq:mu_at_record}.  Apply
Theorem~\ref{thm:first_jet_reconstruction} at \(N=r_k\) to obtain
\eqref{eq:first_jet_record_asymptotic}.  Finally use
Theorem~\ref{thm:record_time_exponent}; the factors \(2\) and
\(\log r_k\) contribute only \(o(\log k)\) to the logarithm.
\end{proof}

\subsection{The centered jet and the short-interval comparison}
\label{subsec:centered_jet_short_interval}

\begin{definition}[Centered first-jet remainder]
\label{def:centered_jet}
Define
\begin{equation}
\mathcal R_{\mathrm H}(N)
:=
\Delta\mathcal A'_0(N)-\frac{\mu(N)N}{2\log N}.
\label{eq:centered_jet_definition}
\end{equation}
\end{definition}

\begin{corollary}
\label{cor:centered_jet_little_o}
One has
\begin{equation}
\boxed{\mathcal R_{\mathrm H}(N)=o(N/\log N).}
\label{eq:centered_jet_little_o}
\end{equation}
\end{corollary}

\begin{proof}
This is exactly Theorem~\ref{thm:first_jet_reconstruction} with its leading
term subtracted.
\end{proof}

\begin{remark}[The \(0.55\) short-interval theorem]
\label{rem:short_interval_055}
Matom\"aki--Ter\"av\"ainen prove the independent uniform estimate
\begin{equation}
\sum_{x<n\le x+H}\mu(n)=o(H)
\qquad
(H\ge x^\theta,\ \theta>0.55),
\label{eq:MT_055}
\end{equation}
with a quantitative logarithmic saving in their proof
\cite{MatomakiTeravainen2023}.  This theorem is not needed for
Theorems~\ref{thm:Astar_theta}--\ref{thm:first_jet_reconstruction}.  Its role
here is diagnostic: it constrains the ordinary-time width over which a fixed
fraction of a record height can be accumulated, but it does not replace the
Kaczorowski--Pintz oscillation input or determine the centered remainder
\(\mathcal R_{\mathrm H}\).
\end{remark}

\begin{remark}[Version 11 supersedes the centered-jet question]
\label{rem:centered_jet_superseded}
The question posed in Version 10 after \eqref{eq:centered_jet_little_o} is no
longer the correct frontier.  Uniform Selberg--Delange expansion yields a full
asymptotic series for \(M_\Omega\), and hence
\(\mathcal R_{\mathrm H}(N)=O(N/\log^2N)\), with arbitrarily many further
inverse-logarithmic terms available.  Thus the first centered remainder is
still perturbative Selberg--Delange data.  The genuinely nonlocal question is
whether the accumulated area or its Dirichlet transform admits spectral
control not already equivalent to the one-variable Mertens problem; see
Section~\ref{sec:v11_squarefree_critical_global}.
\end{remark}

\begin{remark}[Known-versus-derived boundary for Version 10]
\label{rem:v10_known_derived_boundary}
The Kaczorowski--Pintz oscillation theorem, the standard upper estimate
\(M(x)\ll_\eta x^{\Theta+\eta}\), the Vinogradov--Korobov-type zero-free
region estimate, the Selberg--Delange theorem, and the Matom\"aki--Ter\"av\"ainen
short-interval theorem are literature inputs.  Equations
\eqref{eq:Astar_theta}--\eqref{eq:record_occupancy} are elementary deductions
from those inputs.  The auxiliary Weyl path, the exact area identity
\eqref{eq:A0_energy_identity}, the first-jet identification
\eqref{eq:first_jet_cocycle}, and their coupling to record times are the
notebook-specific structural layer.  No bibliographic-priority claim is made
for any consequence unless explicitly stated otherwise.
\end{remark}



\section{Squarefree critical deformation, full jet control, and a global quadratic transform}
\label{sec:v11_squarefree_critical_global}

This section imports only the useful arithmetic content of the earlier
squarefree-polynomial experiment and removes its matrix terminology.  The
finite polynomial, its roots, and its logarithmic derivatives are retained;
the previous numerical Taylor truncations are replaced by compact-uniform
Selberg--Delange estimates and Rouch\'e arguments.  The section then returns
to the Heisenberg deformation and derives a global Perron transform for the
quadratic first jet.

\subsection{The squarefree polynomial and its signed section}
\label{subsec:v11_squarefree_polynomial}

\begin{definition}[Squarefree generating polynomial]
\label{def:v11_Px}
For \(x\ge2\), set
\begin{equation}
P_x(z):=\sum_{n\le x}\mu(n)^2z^{\Omega(n)}.
\label{eq:v11_Px}
\end{equation}
Because \(\mu(n)^2\) restricts to squarefree integers, \(\Omega(n)=\omega(n)\)
on the support of the sum.  Put
\begin{equation}
Q_{\rm sf}(x):=P_x(1)=\sum_{n\le x}\mu(n)^2.
\label{eq:v11_Qsf}
\end{equation}
\end{definition}

\begin{proposition}[Exact signed section and deformation]
\label{prop:v11_signed_section}
For every real or complex \(t\),
\begin{align}
P_x(-1)&=M(x),
\label{eq:v11_Pminus1}\\
P_x(-e^t)&=F_t(x)
:=\sum_{n\le x}\mu(n)e^{t\Omega(n)}.
\label{eq:v11_Ft_Px}
\end{align}
Moreover, for \(\Re s>1\),
\begin{equation}
\sum_{n\ge1}\frac{\mu(n)^2z^{\Omega(n)}}{n^s}
=\prod_p(1+zp^{-s}).
\label{eq:v11_Px_Euler_surface}
\end{equation}
\end{proposition}

\begin{proof}
If \(\mu(n)^2=1\), then \(n\) is squarefree and
\(\mu(n)=(-1)^{\Omega(n)}\).  This gives
\eqref{eq:v11_Pminus1} and \eqref{eq:v11_Ft_Px} term by term.  Absolute
convergence for \(\Re s>1\) gives the Euler product
\eqref{eq:v11_Px_Euler_surface} because each prime may occur at most once.
\end{proof}

\begin{proposition}[Exact factorial moments at the signed point]
\label{prop:v11_factorial_moments}
For every integer \(j\ge1\),
\begin{equation}
P_x^{(j)}(-1)
=(-1)^j\sum_{n\le x}\mu(n)(\Omega(n))_j,
\label{eq:v11_factorial_moments}
\end{equation}
where \((y)_j=y(y-1)\cdots(y-j+1)\).
\end{proposition}

\begin{proof}
Differentiate the finite polynomial \eqref{eq:v11_Px}.  On squarefree
integers,
\((-1)^{\Omega(n)-j}=(-1)^j\mu(n)\), which gives
\eqref{eq:v11_factorial_moments}.
\end{proof}

\subsection{Uniform Selberg--Delange and the unique critical root}
\label{subsec:v11_uniform_SD_root}

For \(z\) in a fixed compact neighborhood of \(-1\), write
\begin{equation}
\prod_p(1+zp^{-s})
=\zeta(s)^z H(s,z),
\qquad
H(s,z):=\prod_p(1+zp^{-s})(1-p^{-s})^z.
\label{eq:v11_SD_factorization}
\end{equation}
The linear terms in \(p^{-s}\) cancel in \(\log H\), so \(H\) is holomorphic
near \((1,-1)\), with \(H(1,-1)=1\).  Define
\begin{equation}
A(z):=\frac{H(1,z)}{\Gamma(z)}.
\label{eq:v11_Az}
\end{equation}
Then
\begin{equation}
A(-1)=0,
\qquad
A'(-1)=-1.
\label{eq:v11_A_simple_zero}
\end{equation}
The uniform Selberg--Delange theorem gives, on a fixed closed disk
\(D=\{|z+1|\le\delta_0\}\),
\begin{equation}
P_x(z)
=x(\log x)^{z-1}\bigl(A(z)+E_x(z)\bigr),
\qquad
\sup_{z\in D}|E_x(z)|\ll\frac1{\log x}.
\label{eq:v11_uniform_SD}
\end{equation}
We use this standard input in the form proved for compact complex-parameter
sets in Selberg--Delange theory
\cite{Tenenbaum2015,deLaBretecheTenenbaum2021}.

\begin{lemma}[Cauchy control of the Selberg--Delange error]
\label{lem:v11_Ex_derivative}
If \(D'\Subset D\), then
\begin{equation}
\sup_{z\in D'}|E_x'(z)|\ll_{D',D}\frac1{\log x}.
\label{eq:v11_Ex_derivative}
\end{equation}
\end{lemma}

\begin{proof}
The error in \eqref{eq:v11_uniform_SD} is holomorphic in \(z\) on the fixed
parameter disk.  Cauchy's integral formula on a circle contained in \(D\)
and surrounding each point of \(D'\) gives
\eqref{eq:v11_Ex_derivative} from the uniform bound in
\eqref{eq:v11_uniform_SD}.
\end{proof}

\begin{theorem}[Unique critical root near \(-1\)]
\label{thm:v11_unique_critical_root}
There exists \(\delta_1>0\) such that for all sufficiently large \(x\),
\(P_x\) has exactly one zero, counted with multiplicity, in
\(|z+1|<\delta_1\).  Denote it by \(\xi_*(x)\).  Then
\begin{equation}
\xi_*(x)\in\mathbb R,
\qquad
\xi_*(x)\to-1,
\label{eq:v11_xistar_real}
\end{equation}
and the zero is simple.
\end{theorem}

\begin{proof}
Choose \(\delta_1\) so that \(A\) has no zero in
\(0<|z+1|\le\delta_1\).  On \(|z+1|=\delta\le\delta_1\),
\(|A(z)|\ge c_\delta>0\).  For sufficiently large \(x\),
\(|E_x(z)|<|A(z)|\) on that circle.  After division by the nonvanishing
factor \(x(\log x)^{z-1}\), Rouch\'e's theorem shows that \(A+E_x\), and
hence \(P_x\), has exactly one zero in the disk.  Taking arbitrarily small
fixed \(\delta\) yields convergence to \(-1\).  The coefficients of \(P_x\)
are real, so uniqueness forces the zero to be real; counting multiplicity
forces simplicity.
\end{proof}

\begin{theorem}[Sharp critical-root localization without a Taylor-tail assumption]
\label{thm:v11_root_M_localization}
Assume \(M(x)\ne0\) and put
\begin{equation}
h_x:=1+\xi_*(x).
\label{eq:v11_hx}
\end{equation}
Then
\begin{equation}
\boxed{
 h_x
 =\frac{M(x)(\log x)^2}{x}\,(1+o(1)).}
\label{eq:v11_hx_M}
\end{equation}
Equivalently,
\begin{equation}
\boxed{
 M(x)
 =\frac{x}{(\log x)^2}h_x\,(1+o(1)).}
\label{eq:v11_M_from_hx}
\end{equation}
\end{theorem}

\begin{proof}
Set
\begin{equation}
e_x:=E_x(-1).
\label{eq:v11_ex}
\end{equation}
Evaluating \eqref{eq:v11_uniform_SD} at \(z=-1\) and using
\(A(-1)=0\) gives the exact identity
\begin{equation}
e_x=\frac{M(x)(\log x)^2}{x}.
\label{eq:v11_ex_M_exact}
\end{equation}
The classical zero-free-region estimate for \(M\) implies \(e_x\to0\).
For \(w\) in a fixed small disk, \eqref{eq:v11_A_simple_zero} and
Lemma~\ref{lem:v11_Ex_derivative} give
\begin{equation}
A(-1+w)+E_x(-1+w)
=-w+e_x+R_x(w),
\label{eq:v11_local_root_equation}
\end{equation}
with
\begin{equation}
|R_x(w)|
\ll |w|^2+\frac{|w|}{\log x}.
\label{eq:v11_local_root_remainder}
\end{equation}
Fix \(0<\varepsilon<1\) and consider the circle
\(|w-e_x|=\varepsilon|e_x|\).  On that circle,
\(|w|\le(1+\varepsilon)|e_x|\), so
\begin{equation}
\sup |R_x(w)|
=o(|e_x|),
\label{eq:v11_Rouche_small_remainder}
\end{equation}
whereas \(|-w+e_x|=\varepsilon|e_x|\).  Rouch\'e's theorem therefore gives
exactly one zero inside the circle.  By
Theorem~\ref{thm:v11_unique_critical_root} it must be \(h_x=1+\xi_*(x)\).
Since \(\varepsilon\) is arbitrary,
\(h_x=e_x(1+o(1))\).  Use \eqref{eq:v11_ex_M_exact} to obtain
\eqref{eq:v11_hx_M}; inversion gives \eqref{eq:v11_M_from_hx}.
\end{proof}

\begin{corollary}[Record exponent in the critical root]
\label{cor:v11_root_record_theta}
At the first absolute record time \(r_k\),
\begin{equation}
\boxed{
|1+\xi_*(r_k)|=r_k^{\Theta-1+o(1)}.}
\label{eq:v11_root_record_scale}
\end{equation}
Hence
\begin{equation}
\boxed{
\lim_{k\to\infty}
\frac{\log|1+\xi_*(r_k)|}{\log r_k}
=\Theta-1.}
\label{eq:v11_root_record_theta}
\end{equation}
In particular,
\begin{equation}
\mathrm{RH}
\quad\Longleftrightarrow\quad
\lim_{k\to\infty}
\frac{\log|1+\xi_*(r_k)|}{\log r_k}
=-\frac12.
\label{eq:v11_root_RH_equivalence}
\end{equation}
\end{corollary}

\begin{proof}
At \(r_k\), \(|M(r_k)|=k=r_k^{\Theta+o(1)}\) by
Theorem~\ref{thm:record_time_exponent}.  Insert this into
\eqref{eq:v11_hx_M}; the logarithmic factor contributes only
\(o(\log r_k)\).  The equivalence follows because RH is equivalent to
\(\Theta=1/2\).
\end{proof}

\begin{warning}[Status of the root criterion]
\label{warn:v11_root_not_RH_progress}
Equation~\eqref{eq:v11_root_RH_equivalence} is an exact arithmetic
characterization, not an independent route to RH.  Its exponent is obtained
through the already established record law
\(r_k=k^{1/\Theta+o(1)}\).  The criterion becomes a forward route only if a
bound on the critical-root displacement can be proved without importing that
record law or equivalent zero information.
\end{warning}

\subsection{What the remaining roots contribute}
\label{subsec:v11_noncritical_roots}

Write the roots of \(P_x\), counted with multiplicity, as
\(\xi_{x,1},\ldots,\xi_{x,d_x}\), with
\(\xi_{x,1}=\xi_*(x)\).  Since
\begin{equation}
Q_{\rm sf}(x)
=\frac6{\pi^2}x+O(\sqrt x),
\label{eq:v11_Qsf_asymptotic}
\end{equation}
we can separate the Mertens-dependent critical root from the rest.

\begin{proposition}[Exact cancellation product]
\label{prop:v11_root_product_exact}
If \(M(x)\ne0\), then
\begin{equation}
\frac{|M(x)|}{Q_{\rm sf}(x)}
=
\prod_{j=1}^{d_x}
\frac{|1+\xi_{x,j}|}{|1-\xi_{x,j}|}.
\label{eq:v11_root_product_exact}
\end{equation}
\end{proposition}

\begin{proof}
Factor
\(P_x(z)=c_x\prod_j(z-\xi_{x,j})\).  Evaluate at \(z=-1\) and \(z=1\),
take absolute values, and divide.
\end{proof}

\begin{theorem}[Deterministic residual product of the noncritical roots]
\label{thm:v11_noncritical_product}
As \(x\to\infty\) through values with \(M(x)\ne0\),
\begin{equation}
\boxed{
\prod_{j=2}^{d_x}
\frac{|1+\xi_{x,j}|}{|1-\xi_{x,j}|}
\sim
\frac{\pi^2}{3(\log x)^2}.}
\label{eq:v11_noncritical_product}
\end{equation}
Thus the dependence on the size of \(M(x)\) is asymptotically concentrated
in the single critical factor.
\end{theorem}

\begin{proof}
From Theorem~\ref{thm:v11_root_M_localization},
\begin{equation}
\frac{|1+\xi_*(x)|}{|1-\xi_*(x)|}
\sim
\frac{|M(x)|(\log x)^2}{2x},
\label{eq:v11_critical_product_factor}
\end{equation}
because \(\xi_*(x)\to-1\).  Divide
\eqref{eq:v11_root_product_exact} by
\eqref{eq:v11_critical_product_factor} and use
\eqref{eq:v11_Qsf_asymptotic}.
\end{proof}

\begin{corollary}[Additive cancellation coordinate of the noncritical roots]
\label{cor:v11_noncritical_logsum}
Define
\begin{equation}
c_{x,j}:=-\log\frac{|1+\xi_{x,j}|}{|1-\xi_{x,j}|}.
\label{eq:v11_cxj}
\end{equation}
Then
\begin{equation}
\boxed{
\sum_{j=2}^{d_x}c_{x,j}
=2\log\log x+\log\frac{3}{\pi^2}+o(1).}
\label{eq:v11_noncritical_logsum}
\end{equation}
\end{corollary}

\begin{proof}
Take minus logarithms in
\eqref{eq:v11_noncritical_product}.
\end{proof}

\subsection{Logarithmic-derivative traces: the previous trace obstruction collapses}
\label{subsec:v11_trace_collapse}

\begin{definition}[Root traces]
\label{def:v11_root_traces}
For \(M(x)\ne0\) and integer \(m\ge1\), define
\begin{equation}
T_m(x):=\sum_{j=1}^{d_x}(1+\xi_{x,j})^{-m}.
\label{eq:v11_Tm}
\end{equation}
\end{definition}

\begin{proposition}[Exact logarithmic-derivative identity]
\label{prop:v11_trace_logderivative}
For every \(m\ge1\),
\begin{equation}
\boxed{
T_m(x)
=-\frac1{(m-1)!}
\left.\frac{d^m}{dz^m}\log P_x(z)\right|_{z=-1}.}
\label{eq:v11_trace_logderivative}
\end{equation}
\end{proposition}

\begin{proof}
From the root factorization,
\(P_x'(z)/P_x(z)=\sum_j(z-\xi_{x,j})^{-1}\).  Differentiate
\(m-1\) times and evaluate at \(z=-1\).  The two sign factors combine to
give exactly \eqref{eq:v11_trace_logderivative}.
\end{proof}

\begin{theorem}[Uniform critical-root dominance of all traces]
\label{thm:v11_trace_dominance}
Let \(h_x=1+\xi_*(x)\).  Then
\begin{equation}
\boxed{
\sup_{m\ge1}|T_m(x)h_x^m-1|\longrightarrow0}
\label{eq:v11_trace_uniform_dominance}
\end{equation}
as \(x\to\infty\) through \(M(x)\ne0\).  Consequently, uniformly for all
integers \(m\ge1\),
\begin{equation}
T_m(x)
=
\left(\frac{x}{M(x)(\log x)^2}\right)^m(1+o(1)).
\label{eq:v11_trace_M_asymptotic}
\end{equation}
\end{theorem}

\begin{proof}
By Theorem~\ref{thm:v11_unique_critical_root}, choose a fixed \(\delta>0\)
such that for all sufficiently large \(x\), every noncritical root satisfies
\(|1+\xi_{x,j}|\ge\delta\).  Hence
\begin{equation}
T_m(x)
=h_x^{-m}+R_m(x),
\qquad
|R_m(x)|\le d_x\delta^{-m}.
\label{eq:v11_trace_split}
\end{equation}
Therefore, once \(|h_x|<\delta\),
\begin{equation}
|R_m(x)h_x^m|
\le d_x\left(\frac{|h_x|}{\delta}\right)^m
\le \frac{d_x|h_x|}{\delta}
\qquad(m\ge1).
\label{eq:v11_trace_relative_bound}
\end{equation}
The degree \(d_x\) is the largest number of distinct prime factors possible
below \(x\), so \(d_x=O(\log x/\log\log x)\).  The classical
Vinogradov--Korobov-type bound for \(M\), together with
\eqref{eq:v11_hx_M}, gives \(d_x|h_x|\to0\).  This proves
\eqref{eq:v11_trace_uniform_dominance}.  Equation
\eqref{eq:v11_trace_M_asymptotic} follows from
\eqref{eq:v11_hx_M}.
\end{proof}

\begin{corollary}[Consecutive traces reconstruct the critical displacement]
\label{cor:v11_trace_ratio}
Uniformly for integers \(m\ge1\),
\begin{align}
\frac{T_m(x)}{T_{m+1}(x)}
&=h_x(1+o(1)),
\label{eq:v11_trace_ratio_h}\\
M(x)
&=\frac{x}{(\log x)^2}
\frac{T_m(x)}{T_{m+1}(x)}(1+o(1)).
\label{eq:v11_trace_ratio_M}
\end{align}
\end{corollary}

\begin{proof}
Apply Theorem~\ref{thm:v11_trace_dominance} to \(m\) and \(m+1\) and divide.
Then use \eqref{eq:v11_M_from_hx}.
\end{proof}

\begin{remark}[What is closed from the earlier trace programme]
\label{rem:v11_trace_obstruction_closed}
The exact traces do not suffer asymptotically from persistent cancellation
among the noncritical roots: the unique root near \(-1\) dominates every
power simultaneously.  This closes the earlier structural ``trace
obstruction'' at the level of reconstructing \(M\).  It does not by itself
improve any bound on \(M\), because the dominant root displacement is already
asymptotically equivalent to \(M(x)(\log x)^2/x\).
\end{remark}

\subsection{Critical-scale limit and uniform control of the full jet tower}
\label{subsec:v11_critical_scale}

Put
\begin{equation}
L_x:=\log\log x,
\qquad
\Phi_x(u):=
\frac{(\log x)^2L_x}{x}
F_{u/L_x}(x).
\label{eq:v11_Phi_definition}
\end{equation}

\begin{theorem}[Compact-uniform critical-scale limit]
\label{thm:v11_critical_scale_limit}
For every compact \(K\subset\mathbb C\),
\begin{equation}
\boxed{
\Phi_x(u)
=ue^{-u}
+O_K\!\left(
\frac1{L_x}+\frac{L_x}{\log x}
\right)}
\label{eq:v11_critical_scale_error}
\end{equation}
uniformly for \(u\in K\).  In particular,
\begin{equation}
\boxed{
\Phi_x(u)\longrightarrow ue^{-u}}
\label{eq:v11_critical_scale_limit}
\end{equation}
uniformly on compact subsets of \(\mathbb C\).
\end{theorem}

\begin{proof}
By \eqref{eq:v11_Ft_Px},
\(F_{u/L_x}(x)=P_x(z_x(u))\) with
\begin{equation}
z_x(u):=-e^{u/L_x}.
\label{eq:v11_zxu}
\end{equation}
Uniformly for \(u\) in a fixed compact set,
\begin{align}
z_x(u)+1
&=-\frac{u}{L_x}+O_K(L_x^{-2}),
\label{eq:v11_zxu_expansion}\\
A(z_x(u))
&=\frac{u}{L_x}+O_K(L_x^{-2}),
\label{eq:v11_A_zxu}\\
(\log x)^{z_x(u)-1}
&=\frac{e^{-u}}{(\log x)^2}
\left(1+O_K(L_x^{-1})\right).
\label{eq:v11_log_power_scale}
\end{align}
Insert these formulas into \eqref{eq:v11_uniform_SD}.  The main term gives
\(ue^{-u}\) after multiplication by \((\log x)^2L_x/x\), while the
Selberg--Delange error contributes \(O_K(L_x/\log x)\).  This proves
\eqref{eq:v11_critical_scale_error}.
\end{proof}

\begin{corollary}[Uniform Cauchy control of growing jets]
\label{cor:v11_cauchy_jets}
Define
\begin{equation}
M_j(x):=\sum_{n\le x}\mu(n)\Omega(n)^j
=\left.\partial_t^jF_t(x)\right|_{t=0}.
\label{eq:v11_Mj}
\end{equation}
Fix \(R>0\).  There is \(C_R>0\) such that for every \(j\ge1\),
\begin{equation}
\boxed{
\left|
\frac{(\log x)^2M_j(x)}{xL_x^{j-1}}
-j(-1)^{j-1}
\right|
\le
\frac{j!}{(2R)^j}C_R
\left(\frac1{L_x}+\frac{L_x}{\log x}\right).}
\label{eq:v11_growing_jet_bound}
\end{equation}
In particular, for every fixed \(\varepsilon>0\), the asymptotic
\begin{equation}
M_j(x)
=j(-1)^{j-1}\frac{xL_x^{j-1}}{(\log x)^2}(1+o(1))
\label{eq:v11_growing_jet_asymptotic}
\end{equation}
holds uniformly throughout any range of \(j=j(x)\) for which the right side
of \eqref{eq:v11_growing_jet_bound} is \(o(j)\); for example the safe range
\begin{equation}
j\le(1-\varepsilon)
\frac{\log L_x}{\log\log L_x}
\label{eq:v11_safe_j_range}
\end{equation}
is admissible for all sufficiently large \(x\).
\end{corollary}

\begin{proof}
The function \(\Phi_x(u)-ue^{-u}\) is holomorphic.  Apply Cauchy's derivative
estimate on \(|u|\le R\) using the compact-uniform error bound on
\(|u|\le2R\).  Since
\begin{equation}
\Phi_x^{(j)}(0)
=\frac{(\log x)^2}{xL_x^{j-1}}M_j(x)
\label{eq:v11_Phi_derivative_Mj}
\end{equation}
and
\((ue^{-u})^{(j)}|_{u=0}=j(-1)^{j-1}\), we obtain
\eqref{eq:v11_growing_jet_bound}.  Stirling's formula gives the displayed
safe range.
\end{proof}

\begin{corollary}[Uniform control of the whole Taylor tail]
\label{cor:v11_full_jet_tail}
For every \(R>0\) there is \(C_R'>0\) such that for \(|u|\le R\),
\begin{equation}
\boxed{
\left|
\sum_{j>J}\frac{M_j(x)}{j!}
\left(\frac{u}{L_x}\right)^j
\right|
\le
C_R'\frac{x}{(\log x)^2L_x}\,2^{-J}}
\label{eq:v11_full_jet_tail}
\end{equation}
for all sufficiently large \(x\), uniformly in \(J\ge0\).
\end{corollary}

\begin{proof}
By Theorem~\ref{thm:v11_critical_scale_limit}, \(\Phi_x\) is uniformly
bounded on \(|u|\le2R\).  Cauchy's coefficient estimate therefore bounds
the coefficient of \(u^j\) in \(\Phi_x\) by \(C_R(2R)^{-j}\).  Summing the
resulting geometric tail for \(|u|\le R\) and undoing the normalization in
\eqref{eq:v11_Phi_definition} gives \eqref{eq:v11_full_jet_tail}.
\end{proof}

\begin{corollary}[Effective limiting differential equation]
\label{cor:v11_effective_ODE}
On every compact set \(K\subset\mathbb C\),
\begin{equation}
(\partial_u+1)^2\Phi_x(u)\longrightarrow0
\label{eq:v11_effective_ODE}
\end{equation}
uniformly on \(K\).
\end{corollary}

\begin{proof}
Compact-uniform convergence of holomorphic functions implies
compact-uniform convergence of all derivatives on smaller compact sets.
Apply this to \eqref{eq:v11_critical_scale_limit} and use
\((\partial_u+1)^2(ue^{-u})=0\).
\end{proof}

\begin{remark}[Gamma-shape coincidence]
\label{rem:v11_gamma_caution}
For positive real \(u\), the function \(ue^{-u}\) is the density of a
\(\Gamma(2,1)\) law.  Here, however, \(u\) is a complex deformation
parameter.  No probability measure is inferred from this formal coincidence.
\end{remark}


\subsection{Perturbative spectral blindness of every fixed transverse jet}
\label{subsec:v12_fixed_jet_blindness}

The preceding compact-uniform theorem settles a question left open by the
first-jet calculation.  Passing to the second transverse derivative does not
leave the Selberg--Delange regime; neither does any fixed derivative order.
The statement below is deliberately about the \emph{leading term}.  It does
not assert that lower-order remainders contain no information about zeta
zeros.

\begin{theorem}[Universal leading term for every fixed Heisenberg jet]
\label{thm:v12_all_fixed_jets}
For each fixed integer \(j\ge1\), define
\begin{equation}
\mathcal J_j(N):=
\left.\partial_t^j\mathcal A_t(N)\right|_{t=0}.
\label{eq:v12_Jj_definition}
\end{equation}
Then
\begin{equation}
\boxed{
\Delta\mathcal J_j(N)
=
\frac{j(-1)^{j-1}}{2}\,\mu(N)
\frac{N(\log\log N)^{j-1}}{\log N}
+o_j\!\left(
\frac{N(\log\log N)^{j-1}}{\log N}
\right).}
\label{eq:v12_all_fixed_jets}
\end{equation}
If \(\mu(N)=0\), the left side vanishes exactly.
\end{theorem}

\begin{proof}
Differentiate the finite exact identity \eqref{eq:At_increment} exactly
\(j\) times at \(t=0\).  Since
\(\partial_t^jF_t(N-1)|_{t=0}=M_j(N-1)\), one obtains
\begin{equation}
\Delta\mathcal J_j(N)
=
\frac{\mu(N)}2
\left[
M_j(N-1)\log N-\Omega(N)^jL(N-1)
\right].
\label{eq:v12_exact_j_increment}
\end{equation}
For fixed \(j\), Corollary~\ref{cor:v11_cauchy_jets} gives
\begin{equation}
M_j(x)
=
 j(-1)^{j-1}
 \frac{x(\log\log x)^{j-1}}{(\log x)^2}
 (1+o_j(1)).
\label{eq:v12_Mj_fixed}
\end{equation}
The Vinogradov--Korobov-type estimate already used for \(L(x)\) implies,
for every fixed \(B>0\),
\begin{equation}
L(x)=O_B\!\left(\frac{x}{(\log x)^B}\right).
\label{eq:v12_L_superlog}
\end{equation}
Since \(\Omega(N)\le \log N/\log2\), choose \(B>j+2\).  Then
\begin{equation}
\Omega(N)^jL(N-1)
=o_j\!\left(
\frac{N(\log\log N)^{j-1}}{\log N}
\right).
\label{eq:v12_Omega_j_L_negligible}
\end{equation}
Insert \eqref{eq:v12_Mj_fixed} and
\eqref{eq:v12_Omega_j_L_negligible} into
\eqref{eq:v12_exact_j_increment}.  The exact factor \(\mu(N)\) also proves
the final assertion.
\end{proof}

\begin{corollary}[The second jet is still Selberg--Delange]
\label{cor:v12_second_jet}
One has
\begin{equation}
\boxed{
\Delta\mathcal A''_0(N)
=-\mu(N)\frac{N\log\log N}{\log N}
+o\!\left(\frac{N\log\log N}{\log N}\right).}
\label{eq:v12_second_jet}
\end{equation}
Thus the second transverse jet is not the first place where a new exponent
or a new zero-sensitive leading term appears.
\end{corollary}

\begin{proof}
Take \(j=2\) in Theorem~\ref{thm:v12_all_fixed_jets}.
\end{proof}

\begin{corollary}[Common record exponent for the entire fixed tower]
\label{cor:v12_record_all_jets}
At the first absolute record times \(r_k\), for every fixed \(j\ge1\),
\begin{equation}
\boxed{
\frac{\log|\Delta\mathcal J_j(r_k)|}{\log k}
\longrightarrow\frac1\Theta.}
\label{eq:v12_record_all_jets}
\end{equation}
\end{corollary}

\begin{proof}
At a first absolute record,
\(\mu(r_k)=\operatorname{sgn}M(r_k)\ne0\).  Theorem
\ref{thm:v12_all_fixed_jets} gives
\begin{equation}
|\Delta\mathcal J_j(r_k)|
=
 r_k^{1+o(1)},
\label{eq:v12_Jj_record_rscale}
\end{equation}
because fixed powers of \(\log r_k\) and \(\log\log r_k\) contribute
\(r_k^{o(1)}\).  The record law
\(r_k=k^{1/\Theta+o(1)}\) completes the proof.
\end{proof}

\begin{remark}[What ``spectrally blind'' means here]
\label{rem:v12_blindness_meaning}
The leading asymptotics in
\eqref{eq:v12_all_fixed_jets}, including the entire compact critical-scale
profile \(ue^{-u}\), are obtained from the local Selberg--Delange structure at
\(s=1\).  Their coefficients do not involve \(\Theta\) or individual
nontrivial zeros.  Consequently, evaluating those universal asymptotics at
\(N=r_k\) can reflect \(\Theta\) only through the already zero-sensitive
location of the record time.  This is the precise, limited sense in which the
fixed perturbative tower is called spectrally blind.  No assertion is made
about beyond-all-orders remainders or nonperturbative values of \(t\).
\end{remark}

\begin{warning}[Closed route: increasing the fixed jet order]
\label{warn:v12_higher_jet_route_closed}
Moving from \(\mathcal A'_0\) to \(\mathcal A''_0\), or to any other fixed
\(\partial_t^j\mathcal A_t|_{t=0}\), does not by itself create a new route
toward \(\Theta\): every such leading term is already forced by
Theorem~\ref{thm:v12_all_fixed_jets}.  The remaining candidates are genuinely
nonlocal or nonperturbative: the accumulated quantities
\(\mathcal J_j(N)\), the complete \(\mathcal A_t(N)\) outside the
\(t=O(1/\log\log N)\) window, the remainder after subtraction of the complete
local Selberg--Delange expansion, and the global Perron--Mellin transform of
Section~\ref{subsec:v11_global_Perron}.
\end{warning}

\subsection{The inner Mertens scale and a second complete tail estimate}
\label{subsec:v11_inner_scale}

Define, whenever \(M(x)\ne0\),
\begin{equation}
a_x:=\frac{M(x)}{M_\Omega(x)}.
\label{eq:v11_ax}
\end{equation}
Since \(M_\Omega(x)\sim x/(\log x)^2\) and the classical zero-free-region
bound gives \(M(x)=o_A(x/(\log x)^A)\) for every fixed \(A>0\), one has
\(a_x\to0\).

\begin{theorem}[Inner-scale linear limit]
\label{thm:v11_inner_scale}
For every compact \(K\subset\mathbb C\),
\begin{equation}
\boxed{
\frac{F_{v a_x}(x)}{M(x)}
=1+v+o_K(1)}
\label{eq:v11_inner_scale_limit}
\end{equation}
uniformly for \(v\in K\).
\end{theorem}

\begin{proof}
The exact exponential remainder gives
\begin{equation}
F_t(x)
=M(x)+tM_\Omega(x)+R_x(t),
\label{eq:v11_Ft_linear_expansion}
\end{equation}
where
\begin{align}
|R_x(t)|
&\le
\frac{|t|^2}{2}
\sum_{n\le x}\mu(n)^2\Omega(n)^2e^{|t|\Omega(n)}
\label{eq:v11_Rx_exp_bound1}\\
&\ll |t|^2x(\log x)^2
\exp\left(\frac{|t|\log x}{\log2}\right).
\label{eq:v11_Rx_exp_bound2}
\end{align}
For \(t=va_x\), uniformly on compact \(v\)-sets,
\begin{equation}
|t|\log x
\ll_K\frac{|M(x)|(\log x)^3}{x}\to0.
\label{eq:v11_tlog_small}
\end{equation}
Hence
\begin{equation}
\frac{|R_x(va_x)|}{|M(x)|}
\ll_K
\frac{|M(x)|(\log x)^6}{x}\to0.
\label{eq:v11_Rx_over_M}
\end{equation}
Finally \(a_xM_\Omega(x)=M(x)\), giving
\eqref{eq:v11_inner_scale_limit}.
\end{proof}

\begin{corollary}[Logarithmic coordinate of the critical root]
\label{cor:v11_tc}
Let
\begin{equation}
t_c(x):=\log(-\xi_*(x)),
\label{eq:v11_tc_def}
\end{equation}
which is real for all sufficiently large \(x\).  Then
\begin{equation}
\boxed{
 t_c(x)=-a_x(1+o(1))
 =-\frac{M(x)(\log x)^2}{x}(1+o(1)).}
\label{eq:v11_tc_asymptotic}
\end{equation}
\end{corollary}

\begin{proof}
By definition \(F_{t_c(x)}(x)=P_x(\xi_*(x))=0\).  Theorem
\ref{thm:v11_inner_scale} has its unique limiting zero at \(v=-1\); Rouch\'e
on a fixed small circle around \(-1\) gives
\(t_c/a_x\to-1\).  The second formula uses
\(M_\Omega(x)\sim x/(\log x)^2\).
\end{proof}

\subsection{The complete critical-scale law for the Heisenberg increment}
\label{subsec:v11_Heisenberg_scale_limit}

\begin{theorem}[Heisenberg critical-scale profile]
\label{thm:v11_Heisenberg_scale_profile}
Let \(K\) be a compact subset of \(\{u\in\mathbb C:\Re u\le0\}\).  Then
uniformly for \(u\in K\),
\begin{equation}
\boxed{
\Delta\mathcal A_{u/L_N}(N)
=\frac{\mu(N)N}{2\log N\,L_N}
ue^{-u}
+o_K\!\left(\frac{N}{\log N\,L_N}\right),}
\label{eq:v11_Heisenberg_scale_profile}
\end{equation}
where \(L_N=\log\log N\).
\end{theorem}

\begin{proof}
Use the exact increment identity \eqref{eq:At_increment}.  By
Theorem~\ref{thm:v11_critical_scale_limit}, replacing \(N\) by \(N-1\)
changes the renormalized first term by \(o_K(1)\), and
\begin{equation}
F_{u/L_N}(N-1)\log N
=\frac{N}{\log N\,L_N}ue^{-u}
+o_K\!\left(\frac{N}{\log N\,L_N}\right).
\label{eq:v11_first_area_term_scale}
\end{equation}
For \(\Re u\le0\),
\(|e^{(u/L_N)\Omega(N)}|\le1\).  The Vinogradov--Korobov-type estimate used
in the proof of Theorem~\ref{thm:first_jet_reconstruction} gives
\begin{equation}
L(N-1)
=o\!\left(\frac{N}{\log N\,L_N}\right).
\label{eq:v11_L_negligible_scale}
\end{equation}
Insert these estimates into \eqref{eq:At_increment}.
\end{proof}

\begin{corollary}[Critical-scale profile at record times]
\label{cor:v11_Heisenberg_record_profile}
If \(\sigma_k=\operatorname{sgn}M(r_k)\), then for compact
\(K\subset\{\Re u\le0\}\),
\begin{equation}
\boxed{
\frac{2\log r_k\,\log\log r_k}{\sigma_k r_k}
\Delta\mathcal A_{u/\log\log r_k}(r_k)
\longrightarrow ue^{-u}}
\label{eq:v11_Heisenberg_record_profile}
\end{equation}
uniformly for \(u\in K\).
\end{corollary}

\begin{proof}
Use \(\mu(r_k)=\sigma_k\) from Corollary~\ref{cor:first_jet_records} in
Theorem~\ref{thm:v11_Heisenberg_scale_profile}.
\end{proof}

\begin{theorem}[Heisenberg increment at the critical root deformation]
\label{thm:v11_critical_root_area_increment}
At \(t=t_c(N)\), one has the exact identity
\begin{equation}
\boxed{
\Delta\mathcal A_{t_c(N)}(N)
=-\frac{e^{t_c(N)\Omega(N)}}2
\left(\mu(N)^2\log N+\mu(N)L(N-1)\right).}
\label{eq:v11_critical_root_area_exact}
\end{equation}
At the first absolute record times,
\begin{equation}
\boxed{
\Delta\mathcal A_{t_c(r_k)}(r_k)
\sim-\frac12 k\log r_k.}
\label{eq:v11_critical_root_area_record}
\end{equation}
\end{theorem}

\begin{proof}
The equation \(F_{t_c(N)}(N)=0\) gives
\begin{equation}
F_{t_c(N)}(N-1)
=-\mu(N)e^{t_c(N)\Omega(N)}.
\label{eq:v11_Ftc_previous}
\end{equation}
Insert this into \eqref{eq:At_increment} to obtain
\eqref{eq:v11_critical_root_area_exact}.  At \(N=r_k\),
\(\mu(r_k)=\sigma_k\) and
\(M(r_k-1)=\sigma_k(k-1)\).  The global historical-average estimate
\eqref{eq:I_uniform_Astar} gives
\begin{equation}
L(r_k-1)=\sigma_k(k-1)\log r_k+o(k\log r_k).
\label{eq:v11_L_record_previous}
\end{equation}
Also \eqref{eq:v11_tc_asymptotic} and the classical zero-free-region bound
imply \(t_c(r_k)\Omega(r_k)\to0\).  Therefore the exponential factor tends
to one and the bracket in \eqref{eq:v11_critical_root_area_exact} is
\(k\log r_k(1+o(1))\).
\end{proof}

\subsection{The record-jump inversion test and its limitation}
\label{subsec:v11_inversion_test}

\begin{corollary}[A pure-arithmetic RH-equivalent record-jump exponent]
\label{cor:v11_jump_RH_equivalent}
The following are equivalent:
\begin{equation}
\mathrm{RH},
\label{eq:v11_RH_item_a}
\end{equation}
and
\begin{equation}
\boxed{
|\Delta\mathcal A'_0(r_k)|=k^{2+o(1)}.}
\label{eq:v11_jump_k2}
\end{equation}
\end{corollary}

\begin{proof}
Corollary~\ref{cor:first_jet_records} gives
\begin{equation}
\frac{\log|\Delta\mathcal A'_0(r_k)|}{\log k}\to\frac1\Theta.
\label{eq:v11_jump_theta_again}
\end{equation}
Thus \eqref{eq:v11_jump_k2} is equivalent to \(1/\Theta=2\), hence to
\(\Theta=1/2\), which is RH.
\end{proof}

\begin{warning}[Why the local jump does not yet invert the logic]
\label{warn:v11_jump_shadow}
The equivalence above is not a new proof mechanism.  The asymptotic
\(\Delta\mathcal A'_0(r_k)\sim\sigma_k r_k/(2\log r_k)\) shows that the
\(\Theta\)-dependence of the local jump enters through the record location
\(r_k\).  A genuinely new inversion must therefore use a nonlocal quantity
such as the accumulated area \(\mathcal A'_0(N)\), the full deformation
\(\mathcal A_t\), or a transform whose singularities can be analyzed without
first inserting the record law.
\end{warning}

\subsection{A global Perron--Mellin transform for the quadratic first jet}
\label{subsec:v11_global_Perron}

Define
\begin{equation}
a(n):=\mu(n)\Omega(n),
\qquad
b(n):=\mu(n)\log n.
\label{eq:v11_ab_sequences}
\end{equation}
Then the first-jet increment can be written
\begin{equation}
2\Delta\mathcal A'_0(n)
=\sum_{m<n}\bigl(a(m)b(n)-b(m)a(n)\bigr).
\label{eq:v11_firstjet_ordered_area}
\end{equation}
For \(\Re s>1\), define
\begin{align}
U(s)
&:=\sum_{n\ge1}\frac{a(n)}{n^s}
=-\frac1{\zeta(s)}\sum_p\frac1{p^s-1},
\label{eq:v11_Us}\\
V(s)
&:=\sum_{n\ge1}\frac{b(n)}{n^s}
=\frac{\zeta'(s)}{\zeta(s)^2}.
\label{eq:v11_Vs}
\end{align}


\begin{lemma}[Local audit at \(s=1\)]
\label{lem:v12_UV_at_one}
As \(s\to1\) inside a fixed sector avoiding the negative real direction from
\(1\),
\begin{align}
\sum_p\frac1{p^s-1}
&=\log\frac1{s-1}+O(1),
\label{eq:v12_P1_near_one}\\
U(s)
&=-(s-1)\log\frac1{s-1}+O(s-1),
\label{eq:v12_U_near_one}\\
V(s)
&=-1+O(s-1).
\label{eq:v12_V_near_one}
\end{align}
\end{lemma}

\begin{proof}
The difference
\begin{equation}
\sum_p\left(\frac1{p^s-1}-p^{-s}\right)
\label{eq:v12_P1_difference}
\end{equation}
is normally convergent in a neighborhood of \(s=1\), because each summand is
\(O(p^{-2\Re s})\).  The standard prime-zeta expansion gives
\(\sum_pp^{-s}=\log(1/(s-1))+O(1)\) in the stated sector, proving
\eqref{eq:v12_P1_near_one}.  The Laurent expansion
\(\zeta(s)=(s-1)^{-1}+O(1)\) gives
\(1/\zeta(s)=(s-1)+O((s-1)^2)\), hence
\eqref{eq:v12_U_near_one}.  Finally
\(V(s)=-(1/\zeta(s))'=-1+O(s-1)\), which proves
\eqref{eq:v12_V_near_one}.
\end{proof}

\begin{remark}[Why the local jet and the global transform are different]
\label{rem:v12_local_vs_global}
Lemma~\ref{lem:v12_UV_at_one} makes the dichotomy explicit.  The local
weighted sum \(M_\Omega\) and every fixed transverse jet inherit their
leading asymptotics from the logarithmic singular structure at \(s=1\).
The Perron integrand below, however, contains both \(w\) and \(s-w\) and can
in principle develop singularities when nontrivial singularities of the two
factors interact.  Determining whether those candidate spectral terms survive
the antisymmetrization is therefore a genuinely different problem from
computing another fixed jet.
\end{remark}

\begin{proposition}[Perron representation of the global first-jet increments]
\label{prop:v11_global_Perron}
Let \(s\in\mathbb C\) with \(\sigma:=\Re s>2\), and choose
\(c\) with \(1<c<\sigma-1\).  Then the absolutely convergent Dirichlet
series
\begin{equation}
\mathscr D_{\mathcal A}(s)
:=\sum_{n\ge1}\frac{2\Delta\mathcal A'_0(n)}{n^s}
\label{eq:v11_DA_definition}
\end{equation}
admits the Perron-limit representation
\begin{equation}
\boxed{
\mathscr D_{\mathcal A}(s)
=\lim_{T\to\infty}\frac1{2\pi i}
\int_{c-iT}^{c+iT}
\frac{U(w)V(s-w)-V(w)U(s-w)}{w}\,dw.}
\label{eq:v11_global_Perron}
\end{equation}
The limit is understood in the usual Perron sense; no contour displacement is
asserted in this proposition.
\end{proposition}

\begin{proof}
For \(m\ne n\) and \(c>0\), Perron's kernel gives
\begin{equation}
\mathbf 1_{m<n}
=\lim_{T\to\infty}\frac1{2\pi i}
\int_{c-iT}^{c+iT}\left(\frac{n}{m}\right)^w\frac{dw}{w}.
\label{eq:v11_Perron_indicator}
\end{equation}
For fixed \(T\), the double sums obtained after inserting
\eqref{eq:v11_Perron_indicator} converge absolutely because
\(c>1\) and \(\sigma-c>1\).  Their products are exactly
\(U(w)V(s-w)\) and \(V(w)U(s-w)\).  The coefficient series in
\eqref{eq:v11_DA_definition} is absolutely convergent for \(\sigma>2\), since
\(|2\Delta\mathcal A'_0(n)|\ll n(\log n)^2\).  Passing through finite
partial sums and then taking the standard Perron limit gives
\eqref{eq:v11_global_Perron}.
\end{proof}

\begin{remark}[Explicit zeta content of the integrand]
\label{rem:v11_global_Perron_zeta_content}
If
\begin{equation}
\mathcal P_1(s):=\sum_p\frac1{p^s-1},
\label{eq:v11_P1}
\end{equation}
then
\begin{equation}
U(w)V(s-w)-V(w)U(s-w)
=
\frac{
\mathcal P_1(s-w)\zeta'(w)\zeta(s-w)
-\mathcal P_1(w)\zeta'(s-w)\zeta(w)
}{\zeta(w)^2\zeta(s-w)^2}.
\label{eq:v11_integrand_zeta_form}
\end{equation}
Thus the accumulated quadratic geometry is connected directly to zeta and
prime-zeta singularities before any record subsequence is selected.  This is
qualitatively different from the local record-jump identity.
\end{remark}

\begin{question}[Global spectral inversion problem]
\label{q:v11_global_spectral_inversion}
Can the contour in \eqref{eq:v11_global_Perron} be displaced far enough to
classify the surviving singularities in the \(s\)-plane and obtain a sharp
abscissa or oscillation law for \(\mathscr D_{\mathcal A}\) or
\(\mathcal A'_0(N)\)?  In particular, does the antisymmetrization in
\eqref{eq:v11_global_Perron} cancel the candidate contributions associated
with the rightmost zeta singularities, or does a nonzero quadratic spectral
term survive?  No answer is claimed here.
\end{question}

\begin{remark}[Version 12 boundary: structural results versus new bounds]
\label{rem:v11_boundary}
The uniform critical root, the compact-uniform profile \(ue^{-u}\), the full
jet-tail estimate, the noncritical-root product law, uniform trace dominance,
and the fixed-jet asymptotics of Theorem~\ref{thm:v12_all_fixed_jets} are
rigorous consequences of standard Selberg--Delange and classical Mertens
bounds combined with the notebook's exact deformations.  They provide new
structural coordinates for \(M\), but no stronger pointwise bound for
\(M(x)\) or new restriction on \(\Theta\).  In particular, the second jet is
now explicitly removed as a prospective escape from the \(s=1\) mechanism.
The Perron transform \eqref{eq:v11_global_Perron} is the first object in this
layer whose proposed use is genuinely nonlocal: any forward implication for
zeta zeros would have to come from an independent analysis of that transform,
a nonperturbative deformation, or a beyond-all-orders remainder rather than
from re-inserting the already known Mertens record exponent.
\end{remark}


\section{Literature boundary and novelty status through Version 12}
\label{sec:status_classification}

The literature boundary is stated component by component rather than by
assigning novelty to the entire package.

\subsection{Known framework}
\label{subsec:known_framework}

The use of translations and functional calculus to form zeta operators has a
direct literature precedent: Herichi--Lapidus construct spectral operators of
the form \(\zeta(\partial)\), together with truncated shifts and operator
criteria related to RH \cite{HerichiLapidus2014}.  Version 8 therefore does
not claim novelty for translation representations or M\"obius inversion at
the untwisted operator level.

The following are classical or direct specializations of classical theory:
\begin{enumerate}[leftmargin=*]
\item the Schr\"odinger pair, CCR, Weyl operators, and the Robertson
inequality for arbitrary observables \cite{Robertson1929,Folland1989};
\item projective representations and multipliers of abelian groups, including
classification through alternating commutator bicharacters
\cite{Kleppner1965,BaggettKleppner1973};
\item twisted group and semigroup algebras
\cite{PackerRaeburn1992,RigalZadunaisky2015};
\item classical M\"obius inversion, Jordan totients, and additive functions
such as \(\Omega\) \cite{Apostol1976};
\item the prime-zeta identity
\(P(s)=\sum_{r\ge1}\mu(r)\log\zeta(rs)/r\), its scaled logarithmic
singularities, and its natural boundary at \(\Re s=0\)
\cite{Froberg1968};
\item the weighted variance decomposition used in
\eqref{eq:curvature_coprime_ANOVA}.
\end{enumerate}
In particular, Theorem~\ref{thm:arithmetic_uncertainty} is not a new
uncertainty principle: it is Robertson's inequality applied to the specific
quadratures \(G_m,G_n\).

\subsection{Notebook-specific formulas not located in the targeted search}
\label{subsec:notebook_specific_formulas}
The targeted literature search carried out for Version 8 did not locate the
following exact formulas as previously published results:
\begin{enumerate}[leftmargin=*]
\item the arithmetic embedding
\(n\mapsto(\ell\log n,\kappa\Omega(n))\) combined with twisted M\"obius
inversion;
\item the resulting coefficients \(z_\vartheta(n)\), including the
prime-pair and two-prime recurrences;
\item the exact curvature formula
\eqref{eq:zsecond_general_formula} and its prime-log variance form;
\item the curvature Dirichlet determinant
\(\mathcal K=\zeta(S_0S_2-S_1^2)/2\);
\item the local monodromy formula
\eqref{eq:K_local_monodromy} with weight \(m_\rho J_2(r)/r\);
\item the multiplicative superadditivity, prime-ray rigidity, and power laws
\eqref{eq:curvature_superadditivity}--\eqref{eq:curvature_power_asymptotic}.
\end{enumerate}
This is a priority statement only in the limited sense ``not located in the
search performed''; it is not a proof of bibliographic novelty.

\subsection{What is specifically new mathematics inside this manuscript}
\label{subsec:new_math_v8}
Independently of priority, the deductions that are new relative to the
standard ingredients listed above are the exact consequences proved here for
the chosen arithmetic cocycle: the twisted inverse, its coefficientwise
curvature, the closed prime-pair geometry, the determinant Dirichlet series,
its local branch weight, and the multiplicative laws of
\(\mathcal C(n)\).  By contrast, the ambient Weyl/cohomology machinery and
the abstract variance identities are not new.

\section{Corrections and extensions carried into Version 8}
\label{sec:v2_audit_corrections}

\begin{remark}
\label{rem:v2_audit_corrections}
Version 8 incorporates the following mathematical corrections and
clarifications: the cocycle is proved nontrivial rather than merely described
as noncommutative; no faithfulness of \(\Pi_{\sigma,\vartheta}\) is
claimed; left and right inverses are constructed separately before being
identified; the uniform Neumann half-plane is sharpened to
\(\sigma>\sigma_*\); the ordered-factorization formula and a two-prime
recurrence are included; real-valuedness of \(z_\vartheta(n)\) for real
\(\vartheta\) is proved; the deformed Dirichlet series is asserted
holomorphic only on its proved open half-plane; the unitary evolution is not
claimed to preserve Schwartz space; and the action on tempered distributions
is defined by continuous transposition; the signs and supports of the two point-mass transpose formulas are corrected explicitly; and no global rearrangement of conditionally convergent sums over zeta zeros is asserted.
\end{remark}

\section{Open problems}
\label{sec:open_problems}

\begin{question}[Analytic continuation]
\label{q:analytic_continuation}
Does $\zeta_{\mathrm H}(s;\vartheta)$ admit a natural meromorphic continuation
from $\Re s\ge2$ to a larger domain for fixed nonzero $\vartheta$?
\end{question}

\begin{question}[Functional symmetry]
\label{q:functional_symmetry}
Is there a completion of $\zeta_{\mathrm H}(s;\vartheta)$ with a symmetry
that reduces at $\vartheta=0$ to the classical reflection $s\mapsto1-s$?
A nontrivial answer would be required before using the deformation to say
anything substantive about the critical line.
\end{question}

\begin{question}[Spectral zeros]
\label{q:spectral_zeros}
Can the zero set of a continued $\zeta_{\mathrm H}$ be realized as a spectral
or resonance set of an operator built intrinsically from
$H_{\mu,\sigma,\vartheta}$, rather than as zeros inserted into a multiplier?
\end{question}

\begin{question}[Classical-zero rigidity]
\label{q:classical_zero_rigidity}
If an analytic two-variable continuation through $\vartheta=0$ exists, does
the exact evenness in $\vartheta$ impose any second-order constraint on the
motion of the classical Riemann zeros beyond
Proposition~\ref{prop:conditional_zero_rigidity}?
\end{question}

\begin{question}[Choice of the second additive coordinate]
\label{q:additive_coordinate}
Which completely additive arithmetic functions $A(n)$, replacing $\Om(n)$
in $y_n=\kappa A(n)$, produce inequivalent deformations?  The choice
$A=\Om$ is distinguished here by its direct relation to the sign of $\mu$ on
square-free integers, but a classification is open.
\end{question}


\begin{question}[Higher deformation curvatures]
\label{q:higher_curvatures}
Determine closed formulas for
\begin{equation}
\frac{d^{2k}}{d\vartheta^{2k}}z_\vartheta(n)\Big|_{\vartheta=0},
\qquad k\ge2,
\label{eq:higher_curvature_derivatives}
\end{equation}
and for their Dirichlet generating series. The fourth-order case is the first
test of whether the collapse to finitely many one-prime moments persists or
irreducible higher prime geometry appears.
\end{question}

\section{Notebook history}
\label{sec:notebook_history}

\noindent\textbf{Version 12.}
Closed the proposed higher-fixed-jet route.  By differentiating the exact
Heisenberg increment at arbitrary fixed order and using the compact-uniform
critical-scale theorem, proved the universal asymptotic
\[
\Delta\partial_t^j\mathcal A_t(N)|_{t=0}
=\frac{j(-1)^{j-1}}2\mu(N)
\frac{N(\log\log N)^{j-1}}{\log N}
+o_j\!\left(\frac{N(\log\log N)^{j-1}}{\log N}\right).
\]
In particular the second jet remains purely Selberg--Delange at leading
order.  This shows that every fixed jet and the whole
\(t=O(1/\log\log N)\) critical window are generated by the local
\(s=1\) mechanism; record evaluations inherit \(\Theta\) only through the
record locations.  Added a local audit of the Dirichlet factors \(U\) and
\(V\) at \(s=1\), separating this perturbative mechanism from the genuinely
nonlocal Perron--Mellin transform.  The forward programme is therefore narrowed
to accumulated Heisenberg areas, nonperturbative deformation, beyond-all-orders
remainders, and contour analysis of \eqref{eq:v11_global_Perron}.

\noindent\textbf{Version 11.}
Added the squarefree critical-deformation layer.  Recast the useful part of
the earlier squarefree-polynomial work without Leslie-matrix terminology,
proved a unique critical root near \(-1\), and replaced the former fixed-jet
Taylor argument by a Rouch\'e proof of
\(1+\xi_*(x)=M(x)(\log x)^2/x\,(1+o(1))\).  Proved the deterministic product
law for all noncritical roots and uniform dominance of the critical root in
all logarithmic-derivative traces.  Established the compact-uniform
critical-scale limit
\((\log x)^2\log\log x\,F_{u/\log\log x}(x)/x\to ue^{-u}\), Cauchy bounds for
growing jets, and a geometric bound for the full Taylor tail.  Added the
inner scale \(t\asymp M/M_\Omega\), the critical-root Heisenberg increment,
and the half-plane critical-scale limit for \(\Delta\mathcal A_t\).  Finally,
derived the global Perron--Mellin representation
\eqref{eq:v11_global_Perron}; its contour deformation and surviving spectral
singularities are left as the principal new open route.

\noindent\textbf{Version 10.}
Added the Mertens-record and Heisenberg-jet layer.  Using the classical
Kaczorowski--Pintz oscillation theorem and the standard upper estimate at the
zero supremum, proved \(A^*(x)=x^{\Theta+o(1)}\), uniform scale separation,
logarithmic historical \(L^p\)-collapse, the record law
\(r_k=k^{1/\Theta+o(1)}\), and the record logarithmic-average theorem.
Introduced the auxiliary Weyl path \(v_n(t)\), proved the exact central identity
\(2\mathcal A_0=MI-E\), identified the first transverse jet with the original
arithmetic cocycle, and combined Selberg--Delange with the incremental area
identity to obtain
\(\Delta\mathcal A'_0(N)=\mu(N)N/(2\log N)+o(N/\log N)\).
Separated these classical-input consequences from the remaining open problem
of additional cancellation in the centered jet.


\noindent\textbf{Version 9.}
Added the curvature--Weil positivity route: an exact Gram polarization of the
cocyclic curvature, the transverse primitive-replica filter
\(\Pi=i\partial_t|_{0}\), a primitive-zero kernel for the full Weil test
class, and the combined curvature--zero Schur kernel.  Proved that positivity
of the combined kernel for every admissible test function is equivalent to RH,
while recording that the converse presently collapses to Weil positivity on a
single diagonal entry.  The remaining genuinely new target is an
unconditional Gram realization of the combined kernel from the Weyl/cocycle
side alone.


\noindent\textbf{Version 8.}
Added the full analytic audit of \(\mathcal O\) and \(\mathcal I\): absolute
convergence, differentiability, sum--integral interchange, grouped
prime-zeta continuation, all endpoint collisions, interior contour
singularities, exact blockwise M\"obius cancellation, strict replica energy,
radical law, correct distinction between \(\mathcal I\) and \(\mathcal O\),
and the monodromy detector.  Added an abstract and a symbol index, and audited
the bibliography and theorem proofs.


\begin{remark}
\label{rem:history_v1}
Version 1 used M\"obius-weighted logarithmic translations, Fourier
multipliers, an eta multiplier in the critical strip, generalized plane-wave
null modes, and a standard CCR non-localization statement.  Its main defect
was that the arithmetic multiplier and the Heisenberg pair were juxtaposed
rather than coupled.
\end{remark}

\begin{remark}
\label{rem:history_v2}
Version 2 removes the critical-strip multiplier as the central mechanism and
replaces ordinary Dirichlet convolution by the Weyl-cocycle product
\eqref{eq:twisted_convolution}.  Heisenberg noncommutativity now changes the
M\"obius inverse coefficients themselves and enters the exact equations of
motion of the M\"obius Hamiltonian.
\end{remark}




\begin{remark}[Version 7]
\label{rem:history_v7}
Version 7 corrects the preceding audit issues and adds the gauge-invariance
obstruction, the scalar Euler-locality obstruction, the norm-convergent
ordered Euler object, its exact defect \(R_\vartheta\), the odd Dirichlet
defect \(\mathcal O(s)\), and the M\"obius filtering of the leading
logarithmic-over-pole term at rational prime-zeta replicas.
\end{remark}

\section{Literature status of the Version 4/5 formulas}
\label{sec:literature_status_v5}

The general ingredients are established mathematics.  Twisted products by
2-cocycles and projective representations belong to the standard theory of
twisted group algebras; see, for example, Packer--Raeburn~\cite{PackerRaeburn1992}.
The prime-zeta identity
\begin{equation}
P(s)=\sum_{r\ge1}\frac{\mu(r)}r\log\zeta(rs)
\label{eq:literature_prime_zeta_identity}
\end{equation}
and the associated logarithmic branch structure and natural boundary are
classical; see Fr\"oberg~\cite{Froberg1968}.

The specific chain developed here---the Heisenberg arithmetic cocycle based on
\((\log n,\Omega(n))\), its twisted M\"obius inverse \(z_\vartheta(n)\), the
closed coefficient curvature
\eqref{eq:zsecond_general_formula}, and its Dirichlet summation
\eqref{eq:K_variance_determinant}---is not identified here with a prior theorem.
The classical prime-zeta formulas explain the continuation of the resulting
\(S_j\), but they do not by themselves produce the curvature formula or the
Hessian identity \eqref{eq:K_hessian_determinant}.

\begin{remark}
\label{rem:history_v5}
Version 8 replaces the overstrong linear-state no-go claim by the exact sesquilinear statement and adds the exact prime-zeta branching source, Hessian/covariance representation, and positivity obstruction and moves the analysis to
the cocycle itself. Differentiating the twisted M\"obius inverse at
\(\vartheta=0\) produces the prime-factor curvature
\eqref{eq:zsecond_general_formula}. Summing it yields
\eqref{eq:K_variance_determinant}, whose local continuation detects every
zeta zero in \(\Re s>1/2\) through \eqref{eq:K_zero_local_form}.
\end{remark}

\begin{thebibliography}{99}

\bibitem{Robertson1929}
H.~P.~Robertson,
\emph{The uncertainty principle},
Physical Review \textbf{34} (1929), 163--164.
doi:10.1103/PhysRev.34.163.

\bibitem{Kleppner1965}
A.~Kleppner,
\emph{Multipliers on Abelian groups},
Mathematische Annalen \textbf{158} (1965), 11--34.

\bibitem{BaggettKleppner1973}
L.~Baggett and A.~Kleppner,
\emph{Multiplier representations of Abelian groups},
Journal of Functional Analysis \textbf{14} (1973), 299--324.
doi:10.1016/0022-1236(73)90075-X.

\bibitem{RigalZadunaisky2015}
L.~Rigal and P.~Zadunaisky,
\emph{Twisted semigroup algebras},
Algebras and Representation Theory \textbf{18} (2015), no.~5, 1155--1186.
doi:10.1007/s10468-015-9525-z; arXiv:1406.2985.

\bibitem{Heisenberg1927}
W.~Heisenberg,
\emph{\"Uber den anschaulichen Inhalt der quantentheoretischen Kinematik und
Mechanik},
Zeitschrift f\"ur Physik \textbf{43} (1927), 172--198.
doi:10.1007/BF01397280.

\bibitem{Folland1989}
G.~B.~Folland,
\emph{Harmonic Analysis in Phase Space},
Annals of Mathematics Studies 122,
Princeton University Press, Princeton, 1989.

\bibitem{ReedSimon1980}
M.~Reed and B.~Simon,
\emph{Methods of Modern Mathematical Physics, Vol.~I: Functional Analysis},
revised and enlarged ed., Academic Press, 1980.


\bibitem{Froberg1968}
C.-E.~Fr\"oberg,
\emph{On the prime zeta function},
BIT Numerical Mathematics \textbf{8} (1968), 187--202.
doi:10.1007/BF01933420.

\bibitem{PackerRaeburn1992}
J.~A.~Packer and I.~Raeburn,
\emph{On the structure of twisted group $C^*$-algebras},
Transactions of the American Mathematical Society \textbf{334} (1992), 685--718.
doi:10.2307/2154478.

\bibitem{Apostol1976}
T.~M.~Apostol,
\emph{Introduction to Analytic Number Theory},
Springer, New York, 1976.
doi:10.1007/978-1-4757-5579-4.

\bibitem{SzekelyRizzo2013}
G.~J.~Sz\'ekely and M.~L.~Rizzo,
\emph{Energy statistics: A class of statistics based on distances},
Journal of Statistical Planning and Inference \textbf{143} (2013), no.~8,
1249--1272.
doi:10.1016/j.jspi.2013.03.018.


\bibitem{HerichiLapidus2014}
H.~Herichi and M.~L.~Lapidus,
\emph{Truncated infinitesimal shifts, spectral operators and quantized
universality of the Riemann zeta function},
Annales de la Facult\'e des Sciences de Toulouse, S\'erie 6, \textbf{23}
(2014), no.~3, 621--664.
doi:10.5802/afst.1419.

\bibitem{KaczorowskiPintz1986}
J.~Kaczorowski and J.~Pintz,
\emph{Oscillatory properties of arithmetical functions. I},
Acta Mathematica Hungarica \textbf{48} (1986), no.~1--2, 173--185.
doi:10.1007/BF01949062.

\bibitem{deLaBretecheTenenbaum2021}
R.~de la Bret\`eche and G.~Tenenbaum,
\emph{Remarks on the Selberg--Delange method},
Acta Arithmetica \textbf{200} (2021), 349--369.
doi:10.4064/aa201024-26-5.

\bibitem{Tenenbaum2015}
G.~Tenenbaum,
\emph{Introduction to Analytic and Probabilistic Number Theory},
3rd ed., Graduate Studies in Mathematics, vol.~163,
American Mathematical Society, Providence, RI, 2015.
doi:10.1090/gsm/163.

\bibitem{MatomakiTeravainen2023}
K.~Matom\"aki and J.~Ter\"av\"ainen,
\emph{On the M\"obius function in all short intervals},
Journal of the European Mathematical Society \textbf{25} (2023), no.~4,
1207--1225.
doi:10.4171/JEMS/1205; arXiv:1911.09076.

\bibitem{Titchmarsh1986}
E.~C.~Titchmarsh,
\emph{The Theory of the Riemann Zeta-Function},
2nd ed., revised by D.~R.~Heath-Brown,
Oxford University Press, Oxford, 1986.


\bibitem{Weil1952}
A.~Weil,
\emph{Sur les ``formules explicites'' de la th\'eorie des nombres premiers},
Comm. S\'em. Math. Univ. Lund [Medd. Lunds Univ. Mat. Sem.],
Tome Suppl\'ementaire (1952), 252--265.

\bibitem{Bombieri2000}
E.~Bombieri,
\emph{Remarks on Weil's quadratic functional in the theory of prime numbers, I},
Atti Accad. Naz. Lincei Cl. Sci. Fis. Mat. Natur. Rend. Lincei Mat. Appl.
(9) \textbf{11} (2000), no.~3, 183--233.
MR1841692; Zbl 1008.11034.

\end{thebibliography}

\end{document}
